%% ============================================================================
%%  CAPITOLO 1 — I numeri e le funzioni reali
%%  Fonte principale per argomenti e dimostrazioni: Marcellini–Sbordone, cap. 1.
%%  Il capitolo è suddiviso in file di paragrafo nella cartella cap01/.
%% ============================================================================
\aperturacapitolo{I numeri e le funzioni reali}{cap:numeri}

\epigraph{Dio ha creato i numeri interi; tutto il resto è opera dell'uomo.}{Leopold Kronecker}

\lettrine[lines=3, findent=2pt, nindent=0pt]{I}{l metodo} comunemente usato in
Matematica consiste nel precisare senza ambiguità i presupposti, da non cambiare durante
l'elaborazione della teoria, e nel dedurre da essi, in modo logico e coerente, il
maggior numero possibile di informazioni. I presupposti sono le \emph{regole del gioco}:
avrebbero potuto essere diverse, ma una volta iniziato il gioco non vengono più cambiate.
In Matematica tali presupposti si chiamano \emph{postulati} o \emph{assiomi}; da essi,
mediante dimostrazioni, si deducono i risultati, o \emph{teoremi}.

In questo capitolo fissiamo le regole del gioco dell'Analisi Matematica: gli assiomi dei
numeri reali. Vedremo che da una manciata di proprietà --- undici in tutto, e il lettore
le conosce quasi tutte fin dalle scuole elementari --- discende l'intero edificio della
teoria. Tra esse ce n'è una, l'\emph{assioma di completezza}, che il lettore
probabilmente non ha mai incontrato in forma esplicita, e che è il vero protagonista del
corso: è ciò che distingue i numeri reali dai numeri razionali, ed è ciò che rende
possibile parlare di limiti, di derivate e di integrali. Nella seconda parte del
capitolo richiameremo il linguaggio degli insiemi e delle funzioni, e passeremo in
rassegna le funzioni elementari (lineari, potenze, esponenziali, logaritmi,
trigonometriche) che useremo in tutto il libro.

\medskip
\noindent\textbf{In questo capitolo.}
\begin{itemize}
  \item Gli assiomi dei numeri reali e le loro prime conseguenze (\S\ref{sec:assiomi},
  \S\ref{sec:conseguenze}).
  \item Il linguaggio degli insiemi e delle relazioni (\S\ref{sec:insiemi}).
  \item I numeri naturali, interi, razionali; perché $\sqrt2$ non è razionale e perché
  $\Q$ non è completo (\S\ref{sec:nzq}).
  \item Approfondimento: l'assioma di Dedekind (\S\ref{subsec:assioma-dedekind}) e un
  modello dei numeri reali costruito a partire da $\Q$ (\S\ref{subsec:modello-reali}).
  \item Funzioni, grafici, funzioni iniettive, suriettive, invertibili, monotòne
  (\S\ref{sec:funzioni}, \S\ref{sec:invertibili}).
  \item Le funzioni elementari: lineari, valore assoluto, potenze, esponenziali,
  logaritmi, trigonometriche (\S\ref{sec:lineari}--\S\ref{sec:trigonometriche}).
  \item Il principio di induzione (\S\ref{sec:induzione}).
  \item Esercizi di riepilogo (\S\ref{sec:esercizi-cap1}).
  \item Punti fondamentali del capitolo (\S\ref{sec:sintesi-cap01}).
\end{itemize}

%% ----------------------------------------------------------------------------
%%  §1.1  Gli assiomi dei numeri reali
%% ----------------------------------------------------------------------------
\section{Gli assiomi dei numeri reali}
\label{sec:assiomi}

Il nostro punto di partenza è quello di \emph{assumere, come postulato, che esista il
sistema dei numeri reali}. Assumiamo cioè che esista un insieme, che chiamiamo insieme
dei numeri reali e che indichiamo con $\R$, i cui elementi si possono sommare e
moltiplicare e si possono confrontare tra loro, e che tali operazioni e tale confronto
obbediscano a un certo numero di regole. Sono le regole che il lettore ha sempre usato,
in modo naturale, da quando ha imparato a far di conto: per tutti è ovvio che $2+3$ sia
uguale a $3+2$, o che $2\cdot3=3\cdot2$. Esprimiamo questo fatto in generale, per due
numeri reali $a$ e $b$, dicendo che valgono le \emph{proprietà commutative} della
somma e del prodotto:
\begin{equation}\label{eq:commutativa}
  a+b=b+a, \qquad a\cdot b=b\cdot a .
\end{equation}
Il lettore conosce sicuramente (provi a scriverle da solo, prima di proseguire) anche
la proprietà associativa, sia rispetto alla somma sia rispetto al prodotto, e la
proprietà distributiva. Elenchiamo ora in modo completo le proprietà che assumiamo
valere \emph{per assioma}, dividendole in tre gruppi: le proprietà relative alle
operazioni, le proprietà relative all'ordinamento e l'assioma di completezza.

\begin{notastorica}[Che cosa sono i numeri reali?]
Per secoli i matematici hanno usato i numeri reali senza una definizione precisa: la
«retta dei numeri» era un'immagine geometrica sufficiente per Newton e per Eulero. Fu
solo nella seconda metà dell'Ottocento, con Richard Dedekind (1831--1916), Georg Cantor
(1845--1918) e Karl Weierstrass (1815--1897), che si sentì il bisogno di \emph{costruire}
$\R$ a partire dai numeri razionali, per dare un fondamento rigoroso al calcolo
infinitesimale. Dedekind, in \emph{Stetigkeit und irrationale Zahlen} (1872), racconta di
aver avuto l'idea delle sue «sezioni» il 24 novembre 1858, mentre preparava le lezioni
di calcolo al Politecnico di Zurigo, insoddisfatto di dover ricorrere all'intuizione
geometrica per dimostrare che una successione crescente e limitata ha un limite.
Nello stesso 1872 Cantor, e già nel 1869 il francese Charles Méray (1835--1911),
costruirono i reali per un'altra via, a partire dalle successioni di razionali che «si
stringono» (le successioni di Cauchy del \cref{sec:bw-cauchy}); Weierstrass ne esponeva
una sua nelle lezioni berlinesi. Il programma comune --- fondare l'intera Analisi sui
numeri naturali, attraverso i razionali --- fu chiamato da Felix Klein, nel 1895,
\emph{aritmetizzazione dell'analisi}. La frase di Kronecker in epigrafe ne riassume lo
spirito, anche se Kronecker stesso rifiutava le costruzioni di Dedekind e di Cantor,
che fanno uso di insiemi infiniti.
L'approccio di questo libro è diverso e più economico: non costruiamo i numeri reali,
ma \emph{postuliamo} che esistano con le proprietà elencate qui sotto. È l'approccio
\emph{assiomatico}, reso rigoroso da David Hilbert (1862--1943) nel 1900. I due approcci
sono equivalenti: si può dimostrare che esiste, a meno di «cambiare nome» agli elementi,
un solo sistema che soddisfa tutti gli assiomi che stiamo per enunciare (un modello
costruito con le sezioni di Dedekind è descritto nel \S\ref{subsec:modello-reali}).
\end{notastorica}

\subsection{Assiomi relativi alle operazioni}

Sono definite le operazioni di \emph{addizione} ($+$) e di \emph{moltiplicazione}
($\cdot$) tra coppie di numeri reali, con le seguenti proprietà ($a$, $b$, $c$ indicano
numeri reali generici).

\begin{assioma}[Proprietà delle operazioni]\label{ass:operazioni}
\leavevmode
\begin{enumerate}[label=\textup{(A\arabic*)}, ref=(A\arabic*), leftmargin=3em]
  \item\label{ass:associativa} \textbf{Proprietà associativa:}
    $(a+b)+c=a+(b+c)$ \ e \ $(a\cdot b)\cdot c=a\cdot(b\cdot c)$.
  \item\label{ass:commutativa} \textbf{Proprietà commutativa:}
    $a+b=b+a$ \ e \ $a\cdot b=b\cdot a$.
  \item\label{ass:distributiva} \textbf{Proprietà distributiva:}
    $a\cdot(b+c)=a\cdot b+a\cdot c$.
  \item\label{ass:neutri} \textbf{Esistenza degli elementi neutri:} esistono in $\R$ due
    numeri \emph{distinti}, indicati con $0$ e $1$, tali che $a+0=a$ e $a\cdot1=a$ per
    ogni $a\in\R$.
  \item\label{ass:opposto} \textbf{Esistenza degli opposti:} per ogni numero reale $a$
    esiste un numero reale, indicato con $-a$, tale che $a+(-a)=0$.
  \item\label{ass:inverso} \textbf{Esistenza degli inversi:} per ogni numero reale
    $a\neq0$ esiste un numero reale, indicato con $a^{-1}$, tale che $a\cdot a^{-1}=1$.
\end{enumerate}
\end{assioma}

\begin{osservazione}
Le proprietà (A1)--(A6) dicono, nel linguaggio dell'algebra, che $\R$ è un
\emph{campo}\index{campo}. Si noti che la sottrazione e la divisione non compaiono tra le
operazioni: esse sono \emph{definite} a partire dall'addizione e dalla moltiplicazione
ponendo
\[
  a-b \coloneqq a+(-b), \qquad \frac{a}{b}\coloneqq a\cdot b^{-1}\quad(b\neq0).
\]
Il simbolo $\coloneqq$ si legge «è definito uguale a». Si noti inoltre che nell'assioma
\ref{ass:inverso} si richiede $a\neq0$: vedremo nel \cref{sec:conseguenze} che questa
restrizione non è un capriccio, ma una necessità logica.
\end{osservazione}

\subsection{Assiomi relativi all'ordinamento}

È definita la relazione di \emph{minore o uguale} ($\le$) tra coppie di numeri reali,
con le seguenti proprietà.

\begin{assioma}[Proprietà dell'ordinamento]\label{ass:ordine}
\leavevmode
\begin{enumerate}[label=\textup{(O\arabic*)}, ref=(O\arabic*), leftmargin=3em]
  \item\label{ass:dicotomia} \textbf{Dicotomia:} per ogni coppia di numeri reali $a$, $b$ si
    ha $a\le b$ oppure $b\le a$.
  \item\label{ass:antisimmetrica} \textbf{Proprietà antisimmetrica:} se valgono
    contemporaneamente $a\le b$ e $b\le a$, allora $a=b$.
  \item\label{ass:ordine-somma} \textbf{Compatibilità con la somma:} se $a\le b$ allora
    $a+c\le b+c$ per ogni $c\in\R$.
  \item\label{ass:ordine-prodotto} \textbf{Compatibilità con il prodotto:} se $0\le a$ e
    $0\le b$ allora $0\le a+b$ e $0\le a\cdot b$.
\end{enumerate}
\end{assioma}

\begin{osservazione}
La proprietà di dicotomia \ref{ass:dicotomia}, applicata con $b=a$, dice in particolare
che $a\le a$ per ogni $a$ (\emph{proprietà riflessiva}). Tra gli assiomi non abbiamo
invece elencato la proprietà \emph{transitiva} («se $a\le b$ e $b\le c$ allora
$a\le c$»): la dimostreremo nel prossimo paragrafo come conseguenza degli altri assiomi.
Le relazioni $\ge$, $<$ e $>$ sono definite a partire da $\le$:
\begin{gather*}
  a\ge b \sse b\le a, \\
  a<b \sse (a\le b \text{ e } a\neq b), \qquad
  a>b \sse (a\ge b \text{ e } a\neq b).
\end{gather*}
Le relazioni $<$ e $>$ sono dette di \emph{minore stretto} e \emph{maggiore stretto}.
\end{osservazione}

\subsection{L'assioma di completezza}

\begin{assioma}[Completezza]\label{ass:completezza}\index{assioma di completezza}
Siano $A$ e $B$ due insiemi non vuoti di numeri reali con la proprietà che
\[
  a\le b \qquad\text{comunque si scelgano } a\in A \text{ e } b\in B.
\]
Allora esiste almeno un numero reale $c$ tale che
\[
  a\le c\le b \qquad\text{qualunque siano } a\in A \text{ e } b\in B.
\]
\end{assioma}

\begin{marginfigure}
\centering
\begin{tikzpicture}[scale=0.9]
  \draw[-Latex, grigio] (-0.3,0) -- (5.2,0);
  \foreach \x in {0.2,0.6,1.0,1.5,1.9,2.2} \fill[blu] (\x,0) circle (1.6pt);
  \foreach \x in {3.0,3.3,3.8,4.3,4.8} \fill[rame] (\x,0) circle (1.6pt);
  \draw[blu, thick, decorate, decoration={brace, amplitude=4pt, mirror}] (0.1,-0.25) -- (2.3,-0.25)
    node[midway, below=5pt, font=\scriptsize] {$A$};
  \draw[rame, thick, decorate, decoration={brace, amplitude=4pt, mirror}] (2.9,-0.25) -- (4.9,-0.25)
    node[midway, below=5pt, font=\scriptsize] {$B$};
  \draw[verde, thick] (2.6,-0.15) -- (2.6,0.15) node[above, font=\scriptsize] {$c$};
\end{tikzpicture}
\caption{L'assioma di completezza: se ogni elemento di $A$ precede ogni elemento di $B$,
esiste un «separatore» $c$.}
\label{fig:completezza}
\end{marginfigure}

A prima vista l'assioma di completezza può sembrare ovvio, e il lettore potrebbe essere
tentato di «dimostrarlo» così: «se tutti i numeri dell'insieme $A$ sono minori o uguali
di tutti i numeri dell'insieme $B$, allora esisterà certamente un numero $c$
intermedio tra $A$ e $B$; basterà infatti scegliere come $c$ il più grande elemento di
$A$, oppure il più piccolo elemento di $B$».

\begin{attenzione}[La frase precedente è sbagliata!]
Non tutti gli insiemi numerici hanno un più grande o un più piccolo elemento. Ad
esempio, l'insieme
\[
  B=\set{1,\ \tfrac12,\ \tfrac13,\ \tfrac14,\ \dots,\ \tfrac1n,\ \dots},
\]
rappresentato sulla retta nella \cref{fig:insieme-1-su-n}, ha il più grande elemento,
che è $1$, ma \emph{non} ha il più piccolo elemento. Potremmo essere tentati di dire
che lo zero è il più piccolo elemento di $B$; ma lo zero \emph{non è un elemento di
$B$}! Infatti $0\neq 1/n$ qualunque sia $n$, perché una frazione è nulla se e soltanto
se il suo numeratore è nullo. Gli elementi di $B$ si accumulano verso lo zero senza
mai raggiungerlo, e per ogni elemento $1/n$ di $B$ ce n'è uno più piccolo,
$1/(n+1)$.
\end{attenzione}

\begin{figure}[htb]
\centering
\begin{tikzpicture}[scale=1.0]
  \draw[-Latex, grigio] (-0.5,0) -- (11.3,0) node[below left, font=\scriptsize] {};
  \foreach \n in {1,...,60} {\fill[blu] ({10/\n},0) circle (1.3pt);}
  \draw[rosso, thick] (0,0) circle (2.2pt);
  \node[below, font=\scriptsize] at (0,-0.1) {$0$};
  \foreach \n/\lab in {1/1, 2/\tfrac12, 3/\tfrac13, 4/\tfrac14, 5/\tfrac15, 8/\tfrac18} {
    \draw[grigio] ({10/\n},0.08) -- ({10/\n},-0.08);
    \node[below, font=\scriptsize] at ({10/\n},-0.1) {$\lab$};}
\end{tikzpicture}
\caption{L'insieme $B=\{1/n : n\in\N\}$ ha massimo $1$ ma non ha minimo: i suoi
elementi si addensano verso $0$, che però (cerchietto rosso) non appartiene a $B$.}
\label{fig:insieme-1-su-n}
\end{figure}

L'assioma di completezza è, in effetti, un assioma molto più profondo di quanto possa
sembrare a prima vista. Come mostreremo nel \cref{sec:nzq}, soltanto tramite l'assioma di
completezza è possibile distinguere l'insieme dei numeri rappresentabili sotto forma di
frazione, i \emph{numeri razionali}, dall'insieme dei numeri reali: i razionali
soddisfano tutti gli assiomi (A1)--(A6) e (O1)--(O4), ma \emph{non} l'assioma di
completezza. In termini geometrici, l'assioma dice che la retta reale \emph{non ha
buchi}: se la tagliamo in due parti, una a sinistra e una a destra, il punto di taglio
è un numero reale.

\begin{idea}[Un'immagine fisica della completezza]
Si pensi a un'asta graduata e a un cursore che scorre su di essa. Se spingiamo il
cursore verso destra fin dove possiamo, con la regola «puoi avanzare finché resti a
sinistra di tutti i punti di $B$», il cursore si fermerà in una posizione precisa $c$:
non può «cadere in un buco» della retta. L'assioma di completezza garantisce che tale
posizione esista. Sui numeri razionali, invece, il cursore potrebbe trovarsi a dover
scegliere tra posizioni razionali sempre più vicine a un punto che razionale non è,
come $\sqrt2$, e non avrebbe dove fermarsi.
\end{idea}

\begin{python}[Verifica numerica delle proprietà di campo]
# Gli assiomi (A1)-(A6) si possono "toccare con mano" con numeri
# esatti: Fraction rappresenta i razionali senza arrotondamenti.
from fractions import Fraction as F

a, b, c = F(3, 4), F(-5, 6), F(7, 2)
print((a + b) + c == a + (b + c))      # (A1) associativa (somma)
print((a * b) * c == a * (b * c))      # (A1) associativa (prodotto)
print(a + b == b + a, a * b == b * a)  # (A2) commutativa
print(a * (b + c) == a * b + a * c)    # (A3) distributiva
print(a + 0 == a, a * 1 == a)          # (A4) elementi neutri
print(a + (-a) == 0)                   # (A5) opposto
print(a * (1 / a) == 1)                # (A6) inverso (a != 0)
\end{python}
\begin{risultato}
True
True
True True
True
True True
True
True
\end{risultato}

\begin{attenzione}[I numeri in virgola mobile non sono numeri reali]
Se nel codice precedente si sostituisce \texttt{Fraction} con i normali numeri
decimali (\texttt{float}), alcune uguaglianze possono fallire: ad esempio in Python
\texttt{0.1 + 0.2 == 0.3} restituisce \texttt{False}, e la proprietà associativa
\texttt{(0.1 + 0.2) + 0.3 == 0.1 + (0.2 + 0.3)} è anch'essa \texttt{False}. I
\texttt{float} sono un insieme \emph{finito} di numeri che approssima $\R$: ottimi per
calcolare, ma non obbediscono agli assiomi. Questa distinzione tra il modello matematico
(i numeri reali) e la sua rappresentazione al calcolatore è una delle prime cose che un
ingegnere deve avere chiara.
\end{attenzione}

%% ----------------------------------------------------------------------------
%%  §1.2  Alcune conseguenze degli assiomi dei numeri reali
%% ----------------------------------------------------------------------------
\section{Alcune conseguenze degli assiomi}
\label{sec:conseguenze}

Nel paragrafo precedente sono state elencate le proprietà dei numeri reali che vengono
assunte come assiomi. \emph{Tutte} le altre proprietà e tutti i teoremi esposti in
questo libro discendono dagli assiomi. Sono conseguenze degli assiomi anche quelle
proprietà elementari che fanno parte del «bagaglio matematico» di ogni studente, come il
fatto che \emph{un prodotto è nullo quando almeno uno dei due fattori è nullo}, oppure la
regola dei segni per il prodotto (che dagli studenti delle scuole elementari è talvolta
accettata come imposizione, perché incompresa) che schematicamente si enuncia:
\emph{meno per meno fa più}; oppure la norma, di frequente applicazione nel risolvere le
disequazioni, secondo cui \emph{moltiplicando entrambi i membri per una quantità
negativa il verso della disequazione cambia}.

Di seguito dimostriamo alcune di queste proprietà. Il lettore non si lasci ingannare
dalla loro apparente banalità: il valore di questo paragrafo sta nel \emph{metodo}. Ogni
passaggio deve essere giustificato da un assioma o da una proprietà già dimostrata, e
nient'altro; è lo stesso rigore che useremo per i teoremi più difficili del libro, e
conviene allenarlo su enunciati di cui si conosce già la risposta.

\subsection{Conseguenze degli assiomi delle operazioni}

\begin{proposizione}[Legge di semplificazione per la somma]\label{prop:semplificazione-somma}
Siano $a,b,c\in\R$. Se $a+b=a+c$, allora $b=c$.
\end{proposizione}

\begin{schema}
\Ipotesi{$a,b,c\in\R$ e $a+b=a+c$.}
\Tesi{$b=c$.}
\end{schema}

\begin{strategia}
Vogliamo «far sparire» $a$ dall'uguaglianza $a+b=a+c$. L'unico strumento che gli assiomi
ci offrono per cancellare un addendo è l'opposto: per \ref{ass:opposto} esiste $-a$ con
$a+(-a)=0$. L'idea è dunque di partire da $b$, scriverlo come $0+b$, sostituire $0$ con
$(-a)+a$ e poi, usando la proprietà associativa, far comparire il blocco $a+b$, che per
ipotesi vale $a+c$. A quel punto si ripercorre la strada a ritroso e si arriva a $c$.
\end{strategia}

\begin{dimostrazione}
Utilizziamo gli assiomi \ref{ass:neutri} e \ref{ass:opposto} (esistenza dello zero e
dell'opposto $-a$) e le proprietà commutativa \ref{ass:commutativa} e associativa
\ref{ass:associativa}:
\[
  b = 0+b = \bigl[a+(-a)\bigr]+b = \bigl[(-a)+a\bigr]+b = (-a)+(a+b).
\]
Essendo per ipotesi $a+b=a+c$, si ottiene
\begin{align*}
  b &= (-a)+(a+b) = (-a)+(a+c) = \bigl[(-a)+a\bigr]+c \\
    &= \bigl[a+(-a)\bigr]+c = 0+c = c+0 = c,
\end{align*}
che è la tesi.
\end{dimostrazione}

\begin{proposizione}[Legge di semplificazione per il prodotto]\label{prop:semplificazione-prodotto}
Siano $a,b,c\in\R$. Se $a\cdot b=a\cdot c$ e $a\neq0$, allora $b=c$.
\end{proposizione}

\begin{schema}
\Ipotesi{$a,b,c\in\R$, \ $a\cdot b=a\cdot c$ \ e \ $a\neq0$.}
\Tesi{$b=c$.}
\end{schema}

\begin{strategia}
Si procede esattamente come nella dimostrazione precedente, scambiando la somma con il
prodotto, lo zero con l'uno e l'opposto con l'inverso. L'unica accortezza è ricordare
che l'inverso $a^{-1}$ esiste \emph{purché} $a\neq0$, per \ref{ass:inverso}: ecco
perché l'ipotesi $a\neq0$ è indispensabile.
\end{strategia}

\begin{dimostrazione}
Poiché $a\neq0$, esiste l'inverso $a^{-1}$. Allora
\begin{align*}
  b &= b\cdot1 = 1\cdot b = (a\cdot a^{-1})\cdot b = (a^{-1}\cdot a)\cdot b \\
    &= a^{-1}\cdot(a\cdot b) = a^{-1}\cdot(a\cdot c) = (a^{-1}\cdot a)\cdot c \\
    &= (a\cdot a^{-1})\cdot c = 1\cdot c = c\cdot 1 = c. \qedhere
\end{align*}
\end{dimostrazione}

\begin{proposizione}[Legge di annullamento del prodotto]\label{prop:annullamento}
\index{legge di annullamento del prodotto}
Siano $a,b\in\R$. Il prodotto $a\cdot b$ è nullo se e soltanto se almeno uno dei due
fattori è nullo.
\end{proposizione}

\begin{schema}
\Ipotesi{$a,b\in\R$.}
\Tesi{$a\cdot b=0 \ \sse\ (a=0 \text{ oppure } b=0)$.}
\end{schema}

\begin{strategia}
L'enunciato è una doppia implicazione («se e soltanto se»): dobbiamo dimostrare due
cose distinte.
\begin{itemize}
  \item \textbf{Il «se»:} se $a=0$ oppure $b=0$ allora $a\cdot b=0$. Basta provare che
  $a\cdot0=0$ per ogni $a$. Sembra ovvio, ma non è un assioma! Lo zero è definito come
  elemento neutro della \emph{somma}, e nulla dice su come si comporta nel
  \emph{prodotto}. Il trucco è scrivere $0=1+0$ e usare la proprietà distributiva.
  \item \textbf{Il «soltanto se»:} se $a\cdot b=0$ allora $a=0$ oppure $b=0$. Se $a=0$
  non c'è nulla da dimostrare; se $a\neq0$ possiamo moltiplicare per $a^{-1}$ e usare la
  prima parte.
\end{itemize}
\end{strategia}

\begin{dimostrazione}
Proviamo preliminarmente l'implicazione con il «se», cioè che $a\cdot0=0$ per ogni
numero reale $a$. Ricordiamo che lo zero è, per l'assioma \ref{ass:neutri}, l'elemento
neutro rispetto alla somma, cioè $a+0=a$ per ogni $a$; e che $1$, elemento neutro
rispetto al prodotto, soddisfa $a\cdot1=a$. In base alla proprietà distributiva
\ref{ass:distributiva} abbiamo allora
\[
  a + a\cdot0 = a\cdot1 + a\cdot0 = a\cdot(1+0) = a\cdot1 = a = a+0,
\]
da cui, semplificando entrambi i membri in base alla
\cref{prop:semplificazione-somma}, otteniamo $a\cdot0=0$. Per la proprietà commutativa
vale anche $0\cdot b=0$.

Proviamo ora l'implicazione con il «soltanto se»: supponiamo $a\cdot b=0$. Se $a=0$
la tesi è raggiunta. Altrimenti, se $a\neq0$, esiste l'inverso $a^{-1}$ e si ha
\[
  b = b\cdot1 = b\cdot(a\cdot a^{-1}) = (b\cdot a)\cdot a^{-1}
    = (a\cdot b)\cdot a^{-1} = 0\cdot a^{-1} = 0 .
\]
Si noti che nell'ultimo passaggio abbiamo utilizzato quanto già provato nella prima
parte della dimostrazione.
\end{dimostrazione}

\begin{osservazione}[Perché non si può dividere per zero]
La \cref{prop:annullamento} spiega perché nell'ambito dei numeri reali non sia
possibile la divisione per zero, cioè perché nell'assioma \ref{ass:inverso} si richieda
$a\neq0$. Infatti, se fosse $a=0$, allora $a\cdot b=0\cdot b=0$ per ogni numero reale
$b$, e perciò \emph{non esiste} alcun numero reale $0^{-1}$ tale che $0\cdot0^{-1}=1$:
se esistesse, avremmo $1=0\cdot0^{-1}=0$, contro l'assioma \ref{ass:neutri} che
richiede $0\neq1$. La divisione per zero non è «vietata» da una regola arbitraria: è
\emph{impossibile} perché contraddittoria con gli assiomi.
\end{osservazione}

\begin{proposizione}[Unicità dell'opposto e dell'inverso]\label{prop:unicita}
\leavevmode
\begin{enumerate}
  \item L'opposto di un numero reale è unico.
  \item L'inverso di un numero reale non nullo è unico.
  \item Per ogni $a\in\R$ vale $-(-a)=a$.
\end{enumerate}
\end{proposizione}

\begin{schema}
\Ipotesi{$a\in\R$; per (ii) $a\neq0$.}
\Tesi{(i) se $a+x=0$ e $a+y=0$, allora $x=y$; \ (ii) se $a\cdot x=1$ e $a\cdot y=1$,
allora $x=y$; \ (iii) $-(-a)=a$.}
\end{schema}

\begin{strategia}
Per dimostrare che un oggetto è unico si procede sempre allo stesso modo: si suppone
che ce ne siano due, $x$ e $y$, con la stessa proprietà, e si dimostra che $x=y$. Qui
la proprietà è «$a+x=0$», e la legge di semplificazione fa il resto. Il punto (iii) è
un'applicazione dell'unicità: $-(-a)$ è \emph{per definizione} l'opposto di $-a$; ma
anche $a$ è un opposto di $-a$, perché $(-a)+a=0$; dunque i due coincidono.
\end{strategia}

\begin{dimostrazione}
(i) In base all'assioma \ref{ass:opposto}, per ogni numero reale $a$ esiste l'opposto
$-a$, tale che $a+(-a)=0$. Se supponiamo che risulti anche $a+b=0$ per qualche $b$,
allora $a+(-a)=a+b$ e, per la legge di semplificazione (\cref{prop:semplificazione-somma}),
$-a=b$: l'opposto è unico.

(ii) Stessa dimostrazione, con il prodotto al posto della somma e la
\cref{prop:semplificazione-prodotto} al posto della \cref{prop:semplificazione-somma}.

(iii) Il numero $-(-a)$ è, per definizione, l'opposto di $-a$. Ma essendo
$a+(-a)=(-a)+a=0$, risulta che anche $a$ è un opposto di $-a$; per l'unicità
dell'opposto provata in (i), $a=-(-a)$.
\end{dimostrazione}

\begin{proposizione}[Regola dei segni]\label{prop:segni}\index{regola dei segni}
Per ogni coppia di numeri reali $a$, $b$ risulta
\begin{enumerate}
  \item $(-a)\cdot b = -(a\cdot b)$;
  \item $(-a)\cdot(-b) = a\cdot b$.
\end{enumerate}
\end{proposizione}

\begin{schema}
\Ipotesi{$a,b\in\R$.}
\Tesi{(i) $(-a)\cdot b=-(a\cdot b)$; \quad (ii) $(-a)\cdot(-b)=a\cdot b$.}
\end{schema}

\begin{strategia}
Che cosa significa dimostrare che $(-a)\cdot b$ è uguale a $-(a\cdot b)$? Significa
dimostrare che $(-a)\cdot b$ è \emph{l'opposto} di $a\cdot b$, cioè che sommato ad
$a\cdot b$ dà zero; per l'unicità dell'opposto (\cref{prop:unicita}) ciò basta. E la
somma $(-a)\cdot b+a\cdot b$ si calcola con la proprietà distributiva, raccogliendo
$b$. Il punto (ii), il celebre «meno per meno fa più», si ottiene applicando due
volte il punto (i) e poi la regola $-(-x)=x$.
\end{strategia}

\begin{dimostrazione}
(i) Per la proprietà distributiva \ref{ass:distributiva} e per la
\cref{prop:annullamento} si ha
\[
  (-a)\cdot b + a\cdot b = \bigl[(-a)+a\bigr]\cdot b = 0\cdot b = 0,
\]
da cui $(-a)\cdot b$ è l'opposto di $a\cdot b$, cioè $(-a)\cdot b=-(a\cdot b)$.

(ii) Come conseguenza del punto (i), applicato due volte, e della proprietà commutativa
\ref{ass:commutativa}, abbiamo
\[
  (-a)\cdot(-b) = -\bigl[a\cdot(-b)\bigr] = -\bigl[(-b)\cdot a\bigr]
  = -\bigl[-(b\cdot a)\bigr] = -\bigl[-(a\cdot b)\bigr];
\]
la conclusione segue infine dalla \cref{prop:unicita}(iii), essendo
$-[-(a\cdot b)]=a\cdot b$.
\end{dimostrazione}

\begin{notastorica}[«Meno per meno fa più»]
La regola dei segni fu a lungo un ostacolo concettuale. Nel Cinquecento i numeri
negativi erano chiamati \emph{numeri absurdi} o \emph{ficti} (Cardano, Stifel) e ancora
nel 1770 Eulero, nella sua \emph{Vollständige Anleitung zur Algebra}, giustificava la
regola con un argomento poco convincente: il prodotto $(-a)\cdot(-b)$ può valere
soltanto $+ab$ oppure $-ab$; poiché $-ab$ è già il valore di $(-a)\cdot b$, esso «non
può che essere» $+ab$. Lo scrittore
Stendhal racconta nella \emph{Vita di Henry Brulard} di aver sofferto da ragazzo per
questa regola, che nessuno dei suoi insegnanti sapeva spiegare. La dimostrazione che
abbiamo appena dato mostra che la regola non è una convenzione: è l'\emph{unico} modo
di definire il prodotto di due numeri negativi se si vuole che continui a valere la
proprietà distributiva.
\end{notastorica}

\subsection{Conseguenze degli assiomi dell'ordinamento}

Gli assiomi del \cref{sec:assiomi} relativi all'ordinamento si riferiscono alla
relazione di minore o uguale ($\le$). Ricordiamo che le relazioni $\ge$, $<$, $>$ sono
state definite a partire da essa: $a\ge b$ significa $b\le a$, e $a<b$ significa
«$a\le b$ e $a\neq b$». Pertanto la relazione $\ge$ gode di proprietà analoghe a quelle
di $\le$. Nelle dimostrazioni che seguono useremo in modo sistematico la seguente
riformulazione dell'ordine in termini di segno di una differenza: essa ci permetterà
di ricondurre ogni disuguaglianza a una proprietà dei numeri $\ge0$.

\begin{proposizione}\label{prop:ordine-differenza}
Per ogni $a,b\in\R$ la relazione $a\le b$ è equivalente alla relazione $b-a\ge0$.
\end{proposizione}

\begin{schema}
\Ipotesi{$a,b\in\R$.}
\Tesi{$a\le b \ \sse\ 0\le b-a$.}
\end{schema}

\begin{strategia}
L'unico assioma che lega l'ordine alla somma è \ref{ass:ordine-somma}: «se $a\le b$
allora $a+c\le b+c$». Per passare da $a\le b$ a $0\le b-a$ basta scegliere $c=-a$; per
tornare indietro, basta scegliere $c=a$.
\end{strategia}

\begin{dimostrazione}
Se $a\le b$, allora per \ref{ass:ordine-somma} con $c=-a$:
\[
  a+(-a)\le b+(-a), \qquad\text{cioè}\qquad 0\le b-a .
\]
Viceversa, se $0\le b-a$, sempre per \ref{ass:ordine-somma} con $c=a$, e per la
proprietà associativa dell'addizione,
\[
  a = 0+a \le (b-a)+a = b+\bigl[(-a)+a\bigr] = b+0 = b. \qedhere
\]
\end{dimostrazione}

\begin{proposizione}[Proprietà transitiva]\label{prop:transitiva}
Se $a\le b$ e $b\le c$, allora $a\le c$.
\end{proposizione}

\begin{schema}
\Ipotesi{$a,b,c\in\R$, \ $a\le b$ \ e \ $b\le c$.}
\Tesi{$a\le c$.}
\end{schema}

\begin{strategia}
Traduciamo tutto in termini di numeri non negativi con la
\cref{prop:ordine-differenza}: le ipotesi dicono che $b-a\ge0$ e $c-b\ge0$. L'assioma
\ref{ass:ordine-prodotto} ci dice che la somma di due numeri non negativi è non
negativa; e la somma di $b-a$ e $c-b$ è proprio $c-a$.
\end{strategia}

\begin{dimostrazione}
Supponiamo $a\le b$ e $b\le c$; per la \cref{prop:ordine-differenza} risulta
$0\le b-a$ e $0\le c-b$. Dall'assioma \ref{ass:ordine-prodotto} si ottiene
\[
  0 \le (b-a)+(c-b) = c-a,
\]
che, ancora per la \cref{prop:ordine-differenza}, equivale ad $a\le c$.
\end{dimostrazione}

\begin{proposizione}\label{prop:segno-opposto}
Per ogni $a\in\R$ risulta $a\ge0$ se e soltanto se $-a\le0$, e $a\le0$ se e soltanto
se $-a\ge0$. Le stesse equivalenze valgono con le disuguaglianze strette.
\end{proposizione}

\begin{schema}
\Ipotesi{$a\in\R$.}
\Tesi{$a\ge0\sse -a\le0$; \quad $a\le0\sse -a\ge0$; \quad analogamente con $>$ e
$<$.}
\end{schema}

\begin{strategia}
Sommare uno stesso numero ai due membri conserva la disuguaglianza
(\ref{ass:ordine-somma}): per passare da $0\le a$ a $-a\le0$ si somma $-a$, per tornare
indietro si somma $a$. Per le disuguaglianze strette basta osservare in più che $a=0$ se
e soltanto se $-a=0$.
\end{strategia}

\begin{dimostrazione}
Per \ref{ass:ordine-somma}, se $0\le a$ allora $0+(-a)\le a+(-a)$, cioè $-a\le0$.
Viceversa, se $-a\le0$ allora, sommando $a$ ad ambo i membri, $(-a)+a\le0+a$, cioè
$0\le a$. Allo stesso modo, se $a\le0$, sommando $-a$ si ottiene $0\le-a$; e se
$0\le-a$, sommando $a$ si ottiene $a\le0$. Infine $-0=0$, perché $0+0=0$ e l'opposto è
unico (\cref{prop:unicita}); e se $-a=0$, allora $a=-(-a)=-0=0$. Dunque $a=0$ se e
soltanto se $-a=0$, e le equivalenze valgono anche con le disuguaglianze strette.
\end{dimostrazione}

\begin{proposizione}[Prodotto di una disuguaglianza per un numero]\label{prop:ordine-prodotto}
\leavevmode
\begin{enumerate}
  \item Se $a\le b$ e $c\ge0$, allora $a\cdot c\le b\cdot c$.
  \item Se $a\le b$ e $c\le0$, allora $a\cdot c\ge b\cdot c$.
\end{enumerate}
\end{proposizione}

\begin{schema}
\Ipotesi{$a,b,c\in\R$ e $a\le b$; nel caso (i) $c\ge0$, nel caso (ii) $c\le0$.}
\Tesi{(i) $ac\le bc$; \quad (ii) $ac\ge bc$.}
\end{schema}

\begin{strategia}
Ancora una volta riduciamo tutto a numeri non negativi: $a\le b$ significa $b-a\ge0$.
Nel caso (i) anche $c\ge0$, quindi per \ref{ass:ordine-prodotto} il prodotto
$(b-a)\cdot c$ è $\ge0$; sviluppandolo con la proprietà distributiva si trova
$bc-ac\ge0$, che è la tesi. Nel caso (ii) è $-c\ge0$ (\cref{prop:segno-opposto}), e si
ragiona allo stesso modo con $-c$ al posto di $c$: è qui che la disuguaglianza
\emph{cambia verso}.
\end{strategia}

\begin{dimostrazione}
(i) Se $a\le b$ allora per la \cref{prop:ordine-differenza} è $0\le b-a$, da cui, per
\ref{ass:ordine-prodotto} e per la proprietà distributiva (insieme alla regola dei segni,
\cref{prop:segni}),
\[
  0\le (b-a)\cdot c = b\cdot c - a\cdot c,
\]
cioè, ancora per la \cref{prop:ordine-differenza}, $a\cdot c\le b\cdot c$.

(ii) Per le ipotesi e per la \cref{prop:segno-opposto} si ha $0\le b-a$ e $0\le -c$.
Da \ref{ass:ordine-prodotto} si ottiene
\[
  0\le (b-a)\cdot(-c) = -b\cdot c + a\cdot c,
\]
da cui, per la \cref{prop:ordine-differenza}, $b\cdot c\le a\cdot c$.
\end{dimostrazione}

\begin{esempio}[Risolvere una disequazione \emph{citando gli assiomi}]
Risolviamo la disequazione $3-2x\ge 7$ giustificando ogni passaggio.
\begin{align*}
  3-2x\ge7 &\ \sse\ (3-2x)+(-3)\ge 7+(-3) \\
           &\ \sse\ -2x\ge4 \\
           &\ \sse\ (-2x)\cdot\bigl(-\tfrac12\bigr)\le 4\cdot\bigl(-\tfrac12\bigr) \\
           &\ \sse\ x\le-2 .
\end{align*}
Il primo passaggio è l'assioma \ref{ass:ordine-somma} con $c=-3$ (sommare lo stesso
numero ai due membri conserva la disuguaglianza); il terzo è la
\cref{prop:ordine-prodotto}(ii) con $c=-\tfrac12\le0$: il «cambio di verso», che a
scuola si impara come regola mnemonica, è esattamente quella proposizione.
Naturalmente, nel seguito del libro non citeremo più gli assiomi a ogni passaggio: ma è bene sapere che si potrebbe.
\end{esempio}

\begin{esercizio}\label{es:zero-minore-uno}
Dimostrare, a partire dagli assiomi, che $0<1$.
\begin{soluzione}
Per \ref{ass:dicotomia} vale $0\le1$ oppure $1\le0$. Se fosse $1\le0$, per la
\cref{prop:segno-opposto} sarebbe $-1\ge0$, e per \ref{ass:ordine-prodotto}
$(-1)\cdot(-1)\ge0$; ma per la regola dei segni $(-1)\cdot(-1)=1\cdot1=1$, quindi
$1\ge0$. Insieme a $1\le0$, per la proprietà antisimmetrica \ref{ass:antisimmetrica} si
avrebbe $1=0$, contro l'assioma \ref{ass:neutri}. Dunque $0\le1$, e poiché $0\neq1$,
$0<1$. \emph{Scelta logica:} la dicotomia lascia due sole possibilità, e quella sbagliata
si esclude per assurdo; l'idea chiave è che $1$ è un quadrato, $1=(-1)\cdot(-1)$, e un
prodotto di due fattori dello stesso segno non è mai negativo.
\end{soluzione}
\end{esercizio}

\begin{esercizio}\label{es:inverso-positivo}
Dimostrare che se $a>0$ allora $a^{-1}>0$. (\emph{Suggerimento:} se fosse $a^{-1}\le0$,
che cosa si potrebbe dire del prodotto $a\cdot a^{-1}$?)
\begin{soluzione}
Sia $a>0$ e supponiamo per assurdo che sia $a^{-1}\le0$. Allora $-a^{-1}\ge0$
(\cref{prop:segno-opposto}) e, poiché anche $a\ge0$, per \ref{ass:ordine-prodotto} e per
la regola dei segni (\cref{prop:segni})
\[
  0\le a\cdot(-a^{-1})=-(a\cdot a^{-1})=-1 ,
\]
cioè $1\le0$ (ancora per la \cref{prop:segno-opposto}, essendo $-(-1)=1$), contro il
risultato $0<1$ dell'\cref{es:zero-minore-uno}. Dunque non è $a^{-1}\le0$: per la
dicotomia \ref{ass:dicotomia} è $a^{-1}\ge0$ e $a^{-1}\neq0$, cioè $a^{-1}>0$.
\emph{Scelta logica:} la tesi è una disuguaglianza stretta, e la sua negazione
$a^{-1}\le0$ si può moltiplicare per $a\ge0$ grazie a \ref{ass:ordine-prodotto}; il
prodotto $a\cdot a^{-1}=1$ è l'unica informazione che lega $a^{-1}$ ad $a$, e porta a
contraddire $0<1$.
\end{soluzione}
\end{esercizio}

%% ----------------------------------------------------------------------------
%%  §1.3  Cenni di teoria degli insiemi
%% ----------------------------------------------------------------------------
\section{Cenni di teoria degli insiemi}
\label{sec:insiemi}

Introduciamo alcune notazioni e definizioni tratte dalla teoria degli insiemi: sono il
linguaggio con cui, da qui in avanti, scriveremo tutte le definizioni e tutti i
teoremi. Sia $S$ un insieme di natura qualsiasi. Per indicare che $x$ è un elemento di
$S$ scriveremo
\[
  x\in S \qquad (x \text{ \emph{appartiene} a } S);
\]
per indicare, invece, che $y$ non è un elemento di $S$, scriveremo
\[
  y\notin S \qquad (y \text{ \emph{non appartiene} a } S).
\]
Un insieme può essere descritto \emph{per elencazione}, scrivendone gli elementi tra
parentesi graffe, come $\{1,2,3\}$, oppure \emph{per proprietà caratteristica}, scrivendo
$\{x\in S : \mathcal{P}(x)\}$ per indicare «l'insieme degli $x$ di $S$ tali che vale la
proprietà $\mathcal{P}$». I due punti si leggono «tale che».

\begin{definizione}[Sottoinsiemi, insieme vuoto]
Se $A$ è un insieme i cui elementi sono anche elementi di $S$, diremo che $A$ è un
\emph{sottoinsieme}\index{sottoinsieme} o \emph{parte} di $S$ e scriveremo $A\subseteq S$.
In simboli,
\[
  (A\subseteq B) \ \sse\ (a\in A \implica a\in B).
\]
Tra i sottoinsiemi di $S$ si considera anche l'\emph{insieme vuoto}\index{insieme vuoto},
cioè l'insieme privo di elementi, che si indica con $\emptyset$; si conviene che
$\emptyset$ sia contenuto in ogni insieme. Se $A\subseteq B$ e $A\neq B$, si dice che
$A$ è una \emph{parte propria} di $B$ e si scrive $A\subsetneq B$. Due insiemi sono
uguali quando hanno gli stessi elementi, cioè $A=B$ se e solo se $A\subseteq B$ e
$B\subseteq A$.
\end{definizione}

\begin{osservazione}
Il simbolo $\implica$ si legge «implica», il simbolo $\sse$ si legge «se e solo se» o
«equivale». Per dimostrare che due insiemi $A$ e $B$ sono uguali si dimostra quasi
sempre la \emph{doppia inclusione}: si prende un generico $x\in A$ e si prova che
$x\in B$, poi si prende un generico $x\in B$ e si prova che $x\in A$. È una tecnica
che useremo molte volte.
\end{osservazione}

\begin{definizione}[Operazioni tra insiemi]\index{intersezione}\index{unione}\index{complementare}
Se $A$ e $B$ sono due sottoinsiemi di $S$ (\cref{fig:venn}):
\begin{itemize}
  \item l'\emph{intersezione} $A\cap B$ è l'insieme degli elementi di $S$ comuni ad $A$ e
  a $B$:
  \[ A\cap B = \{x\in S : x\in A \text{ e } x\in B\}; \]
  \item l'\emph{unione} $A\cup B$ è l'insieme degli elementi di $S$ che appartengono ad
  almeno uno dei due insiemi:
  \[ A\cup B = \{x\in S : x\in A \text{ oppure } x\in B\}; \]
  \item la \emph{differenza} (o \emph{complemento di $B$ rispetto ad $A$}) $A\setminus B$
  è l'insieme degli elementi di $A$ che non appartengono a $B$:
  \[ A\setminus B = \{x\in S : x\in A \text{ e } x\notin B\}. \]
\end{itemize}
In particolare, per $A=S$, l'insieme $S\setminus B$ si chiama \emph{complementare} di
$B$ e si indica con $B^{c}$. Evidentemente si ha
\begin{equation}\label{eq:inclusione-complementari}
  A\subseteq B \ \sse\ B^{c}\subseteq A^{c}.
\end{equation}
Due insiemi $A$ e $B$ con $A\cap B=\emptyset$ si dicono \emph{disgiunti}. L'insieme di
tutti i sottoinsiemi di $S$ si indica con $\parti{S}$ e si chiama \emph{insieme delle
parti}\index{insieme delle parti} di $S$.
\end{definizione}

\begin{figure}[htb]
\centering
\begin{tikzpicture}[scale=0.62, every node/.style={font=\footnotesize}]
  \tikzset{insA/.style={rotate=25}, insB/.style={rotate=-10}}
  \foreach \k/\titolo in {0/{$A\cap B$}, 1/{$A\cup B$}, 2/{$A\setminus B$}} {
  \begin{scope}[shift={(\k*6.4,0)}]
    \draw[grigio!60, thin] (-2.9,-2.1) rectangle (2.9,2.1);
    \node[anchor=north west, grigio] at (-2.85,2.05) {$S$};
    \begin{scope}
      \ifnum\k=0
        \clip[rotate=25] (-0.9,0) ellipse (1.8 and 1.2);
        \fill[blu!35, rotate=-10] (0.9,0) ellipse (1.8 and 1.2);
      \fi
      \ifnum\k=1
        \fill[blu!35, rotate=25] (-0.9,0) ellipse (1.8 and 1.2);
        \fill[blu!35, rotate=-10] (0.9,0) ellipse (1.8 and 1.2);
      \fi
      \ifnum\k=2
        \fill[blu!35, rotate=25] (-0.9,0) ellipse (1.8 and 1.2);
        \fill[white, rotate=-10] (0.9,0) ellipse (1.8 and 1.2);
      \fi
    \end{scope}
    \draw[blu, thick, rotate=25] (-0.9,0) ellipse (1.8 and 1.2);
    \draw[rame, thick, rotate=-10] (0.9,0) ellipse (1.8 and 1.2);
    \node[blu] at (-1.9,0.3) {$A$};
    \node[rame] at (1.9,-0.3) {$B$};
    \node at (0,-2.6) {\titolo};
  \end{scope}}
\end{tikzpicture}
\caption{Diagrammi di Venn: intersezione, unione e differenza di due sottoinsiemi $A$ e
$B$ di $S$. La regione colorata è l'insieme indicato sotto ciascun diagramma.}
\label{fig:venn}
\end{figure}

\begin{esempio}
Sia $S=\N$, $A=\{n\in\N : n \text{ è pari}\}$ e $B=\{n\in\N : n\le 7\}$. Allora
$A\cap B=\{2,4,6\}$, $A\setminus B=\{8,10,12,\dots\}$, $B\setminus A=\{1,3,5,7\}$ e
$A^{c}$ è l'insieme dei numeri dispari. L'insieme delle parti di $\{1,2\}$ è
$\parti{\{1,2\}}=\{\emptyset,\{1\},\{2\},\{1,2\}\}$: in generale, un insieme con $n$
elementi ha $2^{n}$ sottoinsiemi (lo dimostreremo nel \cref{sec:induzione}).
\end{esempio}

\begin{python}[Operazioni tra insiemi con il tipo \texttt{set}]
# Python ha un tipo nativo per gli insiemi finiti: set.
S = set(range(1, 13))               # S = {1, 2, ..., 12}
A = {n for n in S if n % 2 == 0}    # i pari, "per proprietà"
B = {n for n in S if n <= 7}
print("A ∩ B =", A & B)             # intersezione
print("A ∪ B =", A | B)             # unione
print("A \\ B =", A - B)            # differenza
print("B ⊆ A ?", B <= A)            # inclusione
print("A^c   =", S - A)             # complementare rispetto a S
\end{python}
\begin{risultato}
A ∩ B = {2, 4, 6}
A ∪ B = {1, 2, 3, 4, 5, 6, 7, 8, 10, 12}
A \ B = {8, 10, 12}
B ⊆ A ? False
A^c   = {1, 3, 5, 7, 9, 11}
\end{risultato}

\subsection{Prodotto cartesiano e relazioni}

\begin{definizione}[Prodotto cartesiano]\index{prodotto cartesiano}
Siano $A$ e $B$ due insiemi. Si chiama \emph{prodotto cartesiano} di $A$ e $B$, e si
indica con $A\times B$, l'insieme di tutte le \emph{coppie ordinate} $(a,b)$ con la
prima coordinata $a$ appartenente ad $A$ e la seconda coordinata $b$ appartenente a $B$.
Le coppie ordinate sono caratterizzate dalla seguente proprietà:
\begin{equation}\label{eq:coppie}
  (a,b)=(a',b') \quad\text{se e solo se}\quad a=a' \ \text{e}\ b=b'.
\end{equation}
\end{definizione}

\begin{marginfigure}
\centering
\begin{tikzpicture}[scale=0.55]
  \draw[-Latex, grigio] (-0.5,0) -- (4.5,0) node[below, font=\scriptsize] {$A$};
  \draw[-Latex, grigio] (0,-0.5) -- (0,3.5) node[left, font=\scriptsize] {$B$};
  \foreach \x in {1,2,3,4} \foreach \y in {1,2,3} \fill[blu] (\x,\y) circle (2pt);
  \fill[rame] (3,2) circle (3pt) node[above right, font=\scriptsize] {$(a,b)$};
  \draw[tratteggio] (3,0) node[below, font=\scriptsize] {$a$} -- (3,2) -- (0,2) node[left, font=\scriptsize] {$b$};
\end{tikzpicture}
\caption{Il prodotto cartesiano $A\times B$ di due insiemi finiti: una «griglia» di
coppie ordinate.}
\end{marginfigure}

L'aggettivo \emph{ordinate} è essenziale: $(1,2)\neq(2,1)$, mentre $\{1,2\}=\{2,1\}$.
Il prodotto cartesiano $\R\times\R$, che si indica anche con $\R^{2}$, è l'insieme delle
coppie di numeri reali: nel \cref{sec:funzioni} vedremo che esso si identifica con il
piano cartesiano, e il nome «cartesiano» viene proprio da lì.

\begin{definizione}[Relazioni di equivalenza e d'ordine]\index{relazione di equivalenza}\index{relazione d'ordine}
Nel caso particolare in cui sia $A=B$, un sottoinsieme $\mathcal{R}$ di $A\times A$ si
chiama \emph{relazione binaria} in $A$; se $(a,b)\in\mathcal{R}$ si scrive anche
$a\,\mathcal{R}\,b$ e si dice che «$a$ è in relazione con $b$». Una relazione binaria
$\mathcal{R}$ in $A$ si chiama \emph{relazione di equivalenza} se gode delle seguenti
proprietà:
\begin{enumerate}[label=\arabic*)]
  \item \emph{riflessiva}: per ogni $a\in A$ si ha $(a,a)\in\mathcal{R}$;
  \item \emph{simmetrica}: $(a,b)\in\mathcal{R}$ implica $(b,a)\in\mathcal{R}$;
  \item \emph{transitiva}: se $(a,b)\in\mathcal{R}$ e $(b,c)\in\mathcal{R}$, allora
  $(a,c)\in\mathcal{R}$.
\end{enumerate}
Se $(a,b)\in\mathcal{R}$ si scrive $a\sim b$ e si dice che $a$ e $b$ sono
\emph{equivalenti}. Si chiama invece \emph{relazione d'ordine} una relazione che gode
delle proprietà riflessiva, transitiva e
\begin{enumerate}[label=\arabic*'),start=2]
  \item \emph{antisimmetrica}: se $(a,b)\in\mathcal{R}$ e $(b,a)\in\mathcal{R}$, allora
  $a=b$.
\end{enumerate}
\end{definizione}

La relazione di minore o uguale tra numeri reali, considerata nei
\cref{sec:assiomi,sec:conseguenze}, è una relazione d'ordine: la proprietà
riflessiva segue dalla dicotomia \ref{ass:dicotomia}, quella antisimmetrica è l'assioma
\ref{ass:antisimmetrica} e quella transitiva è la \cref{prop:transitiva}. È inoltre un
ordine \emph{totale}, perché due elementi qualsiasi sono sempre confrontabili (ancora
per la dicotomia); l'inclusione $\subseteq$ tra i sottoinsiemi di un insieme è invece
un ordine \emph{parziale}: $\{1\}$ e $\{2\}$ non sono confrontabili.

\begin{definizione}[Classi di equivalenza e insieme quoziente]\index{classe di equivalenza}\index{insieme quoziente}
Data una relazione di equivalenza $\sim$ in $A$, indichiamo con $[a]$ la \emph{classe di
equivalenza} di $a\in A$, cioè l'insieme degli elementi equivalenti ad $a$:
$[a]=\{x\in A : x\sim a\}$. L'insieme delle classi di equivalenza si chiama
\emph{insieme quoziente} e si indica con $A/\!\sim$:
\[
  A/\!\sim\ = \{[a] : a\in A\}.
\]
\end{definizione}

\begin{proposizione}\label{prop:classi}
Due classi di equivalenza o coincidono o sono prive di elementi comuni. Di conseguenza
le classi di equivalenza formano una \emph{partizione} di $A$: ogni elemento di $A$
appartiene a una e una sola classe.
\end{proposizione}

\begin{schema}
\Ipotesi{$\sim$ relazione di equivalenza in $A$; $a,b\in A$.}
\Tesi{$[a]=[b]$ oppure $[a]\cap[b]=\emptyset$; ogni elemento di $A$ appartiene a una e
una sola classe.}
\end{schema}

\begin{strategia}
Dobbiamo provare un'alternativa: «o questo, o quello». La strada naturale è supporre
che \emph{non} valga la seconda possibilità, cioè che le due classi abbiano un elemento
in comune $c$, e dedurne la prima, cioè $[a]=[b]$. L'elemento $c$ fa da «ponte»: è
equivalente sia ad $a$ sia a $b$, e la transitività permette di passare da $a$ a $b$
attraverso $c$. Per l'uguaglianza di insiemi useremo la doppia inclusione.
\end{strategia}

\begin{dimostrazione}
Supponiamo che $[a]\cap[b]\neq\emptyset$ e sia $c$ un elemento comune: $c\sim a$ e
$c\sim b$. Per la proprietà simmetrica $a\sim c$; per la transitiva, da $a\sim c$ e
$c\sim b$ segue $a\sim b$. Proviamo ora che $[a]\subseteq[b]$: se $x\in[a]$, cioè
$x\sim a$, allora da $x\sim a$ e $a\sim b$ segue $x\sim b$, cioè $x\in[b]$. Scambiando
i ruoli di $a$ e $b$ (la situazione è simmetrica) si ottiene $[b]\subseteq[a]$, e
quindi $[a]=[b]$. Infine ogni $a\in A$ appartiene alla propria classe $[a]$ per la
proprietà riflessiva.
\end{dimostrazione}

\begin{esempio}[Le frazioni come classi di equivalenza]
Consideriamo l'insieme delle coppie $(m,n)$ con $m,n\in\Z$ (gli interi, che
richiameremo nel \cref{sec:nzq}) e $n\neq0$, e diciamo che
$(m,n)\sim(m',n')$ se $m\,n'=m'\,n$. Si verifica che è una relazione di equivalenza; la
classe di $(1,2)$ contiene $(2,4)$, $(3,6)$, $(-1,-2)$, \dots{}: è esattamente ciò che
intendiamo quando diciamo che $\tfrac12$, $\tfrac24$, $\tfrac36$ sono «la stessa
frazione». L'insieme quoziente è il modo rigoroso di definire i numeri razionali a
partire dagli interi; noi, avendo postulato l'esistenza di $\R$, non ne avremo bisogno,
ma è utile sapere che la costruzione è possibile.
\end{esempio}

\begin{esercizio}
Stabilire quali tra le seguenti relazioni in $\Z$ sono di equivalenza, quali d'ordine e
quali né l'una né l'altra: (a) $a\,\mathcal{R}\,b$ se $a-b$ è pari; (b)
$a\,\mathcal{R}\,b$ se $a$ divide $b$; (c) $a\,\mathcal{R}\,b$ se $a<b$.
\begin{soluzione}
(a) È di equivalenza: $a-a=0$ è pari; se $a-b$ è pari lo è anche $b-a$; se $a-b$ e
$b-c$ sono pari lo è la loro somma $a-c$. Le classi sono due: i pari e i dispari.
(b) Non è né d'ordine né di equivalenza in $\Z$: $2$ divide $-2$ e $-2$ divide $2$, ma
$2\neq-2$ (manca l'antisimmetria); e $2$ divide $4$ ma $4$ non divide $2$ (manca la
simmetria). In $\N$, invece, è una relazione d'ordine parziale. (c) Non è riflessiva
($a<a$ è falso), quindi non è né d'ordine né di equivalenza; è una relazione
d'\emph{ordine stretto}. \emph{Scelta logica:} per stabilire che una relazione
\emph{gode} di una proprietà serve una verifica valida per tutti gli elementi, mentre per
mostrare che \emph{non} ne gode basta un controesempio: per questo, nelle risposte
negative, abbiamo esibito una sola coppia.
\end{soluzione}
\end{esercizio}

%% ----------------------------------------------------------------------------
%%  §1.4  Numeri naturali, interi, razionali
%% ----------------------------------------------------------------------------
\section{Numeri naturali, interi, razionali}
\label{sec:nzq}

Abbiamo visto come tra gli assiomi dei numeri reali ci sia l'esistenza degli elementi
neutri $0$ e $1$. Quindi apparterranno ad $\R$ anche i risultati delle operazioni
eseguite a partire da $0$ e $1$: in particolare sono numeri reali $1+1=2$, $(1+1)+1=3$, e
così via.

\begin{definizione}[Numeri naturali, interi, razionali]\index{numeri naturali}\index{numeri interi}\index{numeri razionali}
Il sottoinsieme di $\R$ ottenuto a partire da $1$ sommando ripetutamente $1$ si chiama
insieme dei \emph{numeri naturali} e si indica con
\begin{equation}\label{eq:N}
  \N=\{1,2,3,\dots,n,\dots\}.
\end{equation}
Indichiamo con $\Z$ il sottoinsieme di $\R$ costituito dagli elementi di $\N$, dai loro
opposti e dallo zero; i suoi elementi si chiamano \emph{numeri interi} (o
\emph{interi relativi}):
\begin{equation}\label{eq:Z}
  \Z=\{0,\pm1,\pm2,\dots\}=\{0\}\cup\{\pm n : n\in\N\}.
\end{equation}
I risultati della divisione $m/n$ (che, con la terminologia introdotta dagli assiomi,
significa $m\cdot n^{-1}$) con $m,n\in\Z$ e $n\neq0$ si chiamano \emph{numeri razionali}
e si indicano con
\begin{equation}\label{eq:Q}
  \Q=\set{\frac{m}{n} \tc m,n\in\Z,\ n\neq0}.
\end{equation}
\end{definizione}

La frase «sommando ripetutamente $1$» è meno innocente di quanto sembri: definire $n$
come $1+1+\dots+1$, con $n$ addendi, usa proprio il numero $n$ che si vuole definire. La
renderemo rigorosa nel \cref{sec:induzione}, dove $\N$ sarà caratterizzato come il più
piccolo sottoinsieme di $\R$ che contiene $1$ e che, insieme a ogni suo elemento $n$,
contiene anche $n+1$.

Risulta $\N\subseteq\Z\subseteq\Q\subseteq\R$. Naturalmente, essendo $\N$, $\Z$, $\Q$
sottoinsiemi di $\R$, su di essi sono definite le operazioni di addizione e di
moltiplicazione e l'ordinamento indotti da $\R$. Però \emph{essi non soddisfano tutti
gli assiomi dei numeri reali}. Ad esempio $\N$ non soddisfa \ref{ass:opposto}:
nell'ambito dei soli numeri naturali non esiste l'opposto di alcun numero. $\Z$ non
soddisfa \ref{ass:inverso}: tutti i numeri interi non nulli, esclusi $\pm1$, hanno per
inverso un numero reale che non è intero.

Si può verificare che invece $\Q$ soddisfa tutte le proprietà algebriche relative alle
operazioni e all'ordine, cioè gli assiomi (A1)--(A6) e (O1)--(O4): la somma, il
prodotto, l'opposto e l'inverso di numeri razionali sono razionali. \emph{L'unico assioma
non soddisfatto da $\Q$ è l'assioma di completezza}. Questo è il risultato centrale del
paragrafo, e per dimostrarlo premettiamo un teorema celebre, che risale alla scuola
pitagorica.

\begin{teorema}[Irrazionalità di $\sqrt2$]\label{teo:radice2}\index{irrazionalità di $\sqrt2$}
Non esiste alcun numero razionale $c$ tale che $c^{2}=2$.
\end{teorema}

\begin{schema}
\Ipotesi{Nessuna, oltre agli assiomi (e alle proprietà elementari dei numeri interi,
che diamo per note: un intero è pari oppure dispari, il quadrato di un dispari è
dispari, ogni frazione si può ridurre ai minimi termini).}
\Tesi{Per ogni $c\in\Q$ risulta $c^{2}\neq2$.}
\end{schema}

\begin{strategia}
È una tesi \emph{negativa} («non esiste»), e le tesi negative si dimostrano quasi
sempre \emph{per assurdo}: supponiamo che un tale $c$ esista e facciamo discendere da
questa supposizione una contraddizione. La contraddizione sarà di tipo aritmetico:
scriveremo $c=m/n$ con la frazione ridotta ai minimi termini, così che $m$ e $n$ non
siano entrambi pari, e dalla relazione $m^{2}=2n^{2}$ dedurremo che $m$ è pari; ma
allora $m^{2}$ è divisibile per $4$, dunque $n^{2}$ è pari, dunque anche $n$ è pari.
Entrambi pari: ecco l'assurdo. L'intera dimostrazione ruota attorno a un'osservazione
semplice: \emph{il quadrato di un numero dispari è dispari}, e quindi se $m^{2}$ è pari,
$m$ è pari.
\end{strategia}

\begin{dimostrazione}
Sia per assurdo $c$ un numero razionale tale che $c^{2}=2$; possiamo supporre $c>0$
(altrimenti consideriamo $-c$, che ha lo stesso quadrato). In base alla \eqref{eq:Q}
esistono $m,n$ numeri interi, che possiamo supporre entrambi positivi, tali che
$c=m/n$. Se necessario, possiamo «semplificare» la frazione $m/n$, ottenendo $m$ e $n$
\emph{non entrambi pari}. Risulta $(m/n)^{2}=c^{2}=2$, cioè
\[
  2n^{2}=m^{2}.
\]
Essendo il primo membro $2n^{2}$ un numero intero pari, anche $m^{2}$ deve essere pari;
ma allora anche $m$ deve essere pari (se $m$ fosse dispari, $m=2h+1$, anche
$m^{2}=4h^{2}+4h+1=2(2h^{2}+2h)+1$ sarebbe dispari). Quindi $m=2k$, con $k$ intero.
Ne segue che
\begin{equation}\label{eq:radice2-passo}
  2n^{2}=m^{2}=4k^{2}, \qquad\text{cioè}\qquad n^{2}=2k^{2}.
\end{equation}
Ripetendo lo stesso ragionamento sulla \eqref{eq:radice2-passo}, risulta che anche
$n$ deve essere un numero pari, ciò che contrasta con l'ipotesi che $m$ ed $n$ siano
numeri interi non entrambi pari. L'assurdo prova la tesi.
\end{dimostrazione}

\begin{notastorica}[Ippaso e la crisi degli incommensurabili]
Il \cref{teo:radice2} è, nella sostanza, uno dei più antichi teoremi della storia
della matematica. I pitagorici (VI--V secolo a.C.) credevano che «tutto è numero», nel
senso che il rapporto tra due grandezze omogenee qualsiasi --- due segmenti, per esempio
--- fosse sempre esprimibile come rapporto di numeri interi. La scoperta che la diagonale
e il lato di un quadrato non hanno una misura comune (sono \emph{incommensurabili}: il
loro rapporto è $\sqrt2$) fu un trauma per la scuola. La tradizione, riportata da
Giamblico, attribuisce la scoperta a Ippaso di Metaponto, e racconta che i confratelli lo
avrebbero fatto annegare, o quantomeno espulso, per aver divulgato il segreto. La
dimostrazione «pari e dispari» che abbiamo dato è citata da Aristotele negli
\emph{Analitici primi} come esempio tipico di ragionamento per assurdo, ed è riportata in
appendice al libro X degli \emph{Elementi} di Euclide.
La crisi ebbe conseguenze durature. Se gli interi e le frazioni non bastano a misurare
tutte le grandezze, si può rinunciare a misurarle e confrontarne direttamente i
rapporti: è la \emph{teoria delle proporzioni}, che la tradizione attribuisce a Eudosso
di Cnido e che occupa il libro~V degli \emph{Elementi}. La definizione~5 di quel libro
dice, in linguaggio moderno, che $a:b=c:d$ quando, per ogni coppia di naturali $m,n$,
il multiplo $ma$ supera, uguaglia o è minore di $nb$ esattamente quando $mc$ supera,
uguaglia o è minore di $nd$. Poiché $ma>nb$ equivale a $\frac ab>\frac nm$, un rapporto è
individuato dal modo in cui \emph{divide} le frazioni in quelle che gli stanno sotto e
quelle che gli stanno sopra: è l'idea delle sezioni di Dedekind
(\S\ref{subsec:assioma-dedekind}), con oltre due millenni di anticipo. Mancava soltanto
l'ipotesi che \emph{ogni} taglio delle frazioni provenga da un rapporto, cioè la
completezza. Il prezzo fu una lunga separazione tra il mondo discreto dei numeri e quello
continuo delle grandezze geometriche, che si ricompose solo in età moderna.
\end{notastorica}

\begin{teorema}[$\Q$ non è completo]\label{teo:Q-non-completo}\index{completezza!di $\Q$ (mancanza)}
L'insieme $\Q$ dei numeri razionali non soddisfa l'assioma di completezza.
\end{teorema}

\begin{schema}
\Ipotesi{Gli assiomi dei numeri reali e il \cref{teo:radice2}.}
\Tesi{Esistono due sottoinsiemi non vuoti $A,B\subseteq\Q$ con $a\le b$ per ogni
$a\in A$, $b\in B$, tali che \emph{nessun} numero razionale $c$ soddisfa $a\le c\le b$
per ogni $a\in A$ e $b\in B$.}
\end{schema}

\begin{strategia}
Dobbiamo esibire un controesempio: due «mezze rette razionali» separate da un buco.
Il buco, naturalmente, sarà $\sqrt2$. Prendiamo per $A$ i razionali «a sinistra di
$\sqrt2$» e per $B$ quelli «a destra», ma --- attenzione --- senza usare $\sqrt2$, che
in $\Q$ non esiste: diremo che $a\in A$ se $a\le0$ oppure $a^{2}<2$, e che $b\in B$ se
$b>0$ e $b^{2}>2$. Per il \cref{teo:radice2} ogni razionale sta in $A$ o in $B$, e in
uno solo dei due. Se esistesse un separatore razionale $c$, dovrebbe stare in $A$
oppure in $B$. Se $c\in A$, $c$ sarebbe il massimo di $A$; ma mostreremo che $A$ non
ha massimo: per ogni $c\in A$ con $c>0$ riusciamo a trovare un razionale \emph{un po'
più grande}, $c+1/n$, che sta ancora in $A$, cioè con $(c+1/n)^{2}<2$. Il lavoro
tecnico consiste nel trovare un $n$ abbastanza grande, e lo faremo con una stima
esplicita. Il caso $c\in B$ è analogo.
\end{strategia}

\begin{dimostrazione}
Consideriamo gli insiemi
\begin{align}
  A&=\{a\in\Q : a\le0\}\cup\{a\in\Q : a>0,\ a^{2}<2\}, \label{eq:AB-radice2}\\
  B&=\{b\in\Q : b>0,\ b^{2}>2\}. \notag
\end{align}
Entrambi sono non vuoti ($0\in A$, $2\in B$). Tutti i numeri di $A$ sono minori di
tutti i numeri di $B$: se $a\le0<b$ è ovvio; se $0<a$ e fosse $a\ge b$, per la
\cref{prop:ordine-prodotto} sarebbe $a^{2}\ge ab\ge b^{2}>2$, contro $a^{2}<2$. Inoltre,
per il \cref{teo:radice2}, risulta $A\cup B=\Q$ e $A\cap B=\emptyset$: un razionale
positivo ha quadrato minore di $2$ oppure maggiore di $2$, mai uguale.

Se esistesse un numero \emph{razionale} $c$ con la proprietà che $a\le c\le b$ per ogni
$a\in A$ e $b\in B$, tale numero dovrebbe appartenere ad $A$ oppure a $B$.

\emph{Supponiamo $c\in A$.} Non potendo essere $c\le0$ (perché $c\ge1\in A$), ne segue
che $c>0$ e $c^{2}<2$. Sia $n$ un numero naturale maggiore di
$(2c+1)/(2-c^{2})$;\footnote{L'esistenza di un numero naturale maggiore di un
qualsiasi numero reale assegnato è la \emph{proprietà di Archimede}, che dimostreremo
nel prossimo capitolo come conseguenza dell'assioma di completezza (\cref{teo:archimede}). Qui, trattandosi
del numero razionale $(2c+1)/(2-c^{2})=p/q$ con $p,q\in\N$, basta prendere
$n=p+1$: infatti $p/q\le p<p+1$.} tale $n$ esiste. Allora, essendo $1/n^{2}\le1/n$,
\begin{equation}\label{eq:stima-radice2}
  \Bigl(c+\frac1n\Bigr)^{2} = c^{2}+\frac{2c}{n}+\frac{1}{n^{2}}
  \le c^{2}+\frac{2c+1}{n} < c^{2}+(2-c^{2}) = 2,
\end{equation}
dove nell'ultima disuguaglianza abbiamo usato $n>(2c+1)/(2-c^{2})$, cioè
$(2c+1)/n<2-c^{2}$. Dunque $c+\frac1n\in A$, il che è assurdo perché $c$ è maggiore o
uguale di tutti gli elementi di $A$, mentre $c+\frac1n>c$.

\emph{Supponiamo $c\in B$.} Analogamente si perviene a un assurdo: $c>0$, $c^{2}>2$, e
scelto $n$ naturale maggiore di $2c/(c^{2}-2)$ si ha
\[
  \Bigl(c-\frac1n\Bigr)^{2} = c^{2}-\frac{2c}{n}+\frac{1}{n^{2}} > c^{2}-\frac{2c}{n}
  > c^{2}-(c^{2}-2) = 2,
\]
e inoltre $c-\frac1n>0$ (perché $n>2c/(c^{2}-2)>1/c$, dato che $2c^{2}>c^{2}-2>0$);
quindi $c-\frac1n\in B$, assurdo perché $c$ è minore o uguale di tutti gli elementi di
$B$.

In conclusione nessun numero razionale separa $A$ e $B$: l'assioma di completezza non
vale in $\Q$.
\end{dimostrazione}

\begin{figure}[htb]
\centering
\begin{tikzpicture}[scale=0.93]
  \draw[-Latex, grigio] (-0.5,0) -- (11.5,0);
  % razionali in A (sotto) e in B (sopra) che si avvicinano a sqrt2
  \pgfmathsetmacro{\s}{sqrt(2)}
  \foreach \x in {0,0.25,...,1.3} \fill[blu] (\x*5,0) circle (1.2pt);
  \foreach \x in {1.5,1.75,...,2.3} \fill[rame] (\x*5,0) circle (1.2pt);
  \foreach \p/\q in {1/1, 3/2, 7/5, 17/12, 41/29, 99/70} {
    \pgfmathsetmacro{\v}{\p/\q}
    \ifdim \v pt < \s pt \fill[blu] (\v*5,0) circle (1.6pt); \else \fill[rame] (\v*5,0) circle (1.6pt); \fi}
  \foreach \x/\l in {0/0, 5/1, 10/2} \draw[grigio] (\x,0.07) -- (\x,-0.07);
  \node[below, font=\scriptsize, grigio] at (0,-0.08) {$0$};
  \node[blu, below, font=\scriptsize] at (5,-0.08) {$1$};
  \node[below, font=\scriptsize, grigio] at (10,-0.08) {$2$};
  % etichette con linee di richiamo, per non sovrapporle vicino a sqrt2
  \draw[grigio, very thin] (7.0,-0.06) -- (6.3,-0.42)
    node[blu, left, font=\scriptsize, inner sep=1pt] {$\tfrac75$};
  \draw[grigio, very thin] (7.083,0.06) -- (6.3,0.42)
    node[rame, left, font=\scriptsize, inner sep=1pt] {$\tfrac{17}{12}$};
  \node[rame, above, font=\scriptsize] at (7.5,0.08) {$\tfrac32$};
  \draw[grigio, very thin] (\s*5,-0.08) -- (8.3,-0.42)
    node[rosso, right, font=\scriptsize, inner sep=1pt] {$\sqrt2\notin\Q$};
  \draw[rosso, thick] (\s*5,0) circle (2.4pt);
  \draw[blu, thick, decorate, decoration={brace, amplitude=4pt, mirror}] (-0.3,-0.85) -- (6.95,-0.85)
    node[midway, below=5pt, font=\scriptsize] {$A=\{a\le0\}\cup\{a>0,\ a^{2}<2\}$};
  \draw[rame, thick, decorate, decoration={brace, amplitude=4pt}] (7.2,0.75) -- (11.3,0.75)
    node[midway, above=5pt, font=\scriptsize] {$B=\{b>0,\ b^{2}>2\}$};
\end{tikzpicture}
\caption{Il «buco» di $\Q$: gli insiemi $A$ e $B$ della \eqref{eq:AB-radice2} esauriscono
$\Q$, ma nessun razionale li separa. Le frazioni $1,\tfrac75,\tfrac{41}{29}\in A$ e
$\tfrac32,\tfrac{17}{12},\tfrac{99}{70}\in B$ si avvicinano quanto si vuole al punto
$\sqrt2$, che però non appartiene a $\Q$.}
\label{fig:buco-Q}
\end{figure}

\begin{osservazione}[Dove vive $\sqrt2$]
Osserviamo ora che i due insiemi $A$ e $B$, costituiti da numeri razionali, possono
essere riguardati come insiemi di numeri \emph{reali}: e quindi, per l'assioma di
completezza, esiste un numero reale $c$ con la proprietà che $a\le c\le b$ per ogni
$a\in A$, $b\in B$. Tale numero --- che si può dimostrare essere unico --- è tale che
$c^{2}=2$, è irrazionale e si denota con il simbolo $\sqrt2$. In altre parole:
\emph{l'assioma di completezza è esattamente ciò che garantisce l'esistenza di
$\sqrt2$}, e più in generale, come vedremo, l'esistenza delle radici, dei logaritmi, dei
limiti. Senza di esso i numeri reali sarebbero indistinguibili dai razionali, e la retta
sarebbe piena di buchi.
\end{osservazione}

\begin{idea}[Un riassunto in parole semplici]
Nell'insieme dei numeri naturali $\N$ si possono eseguire le operazioni di addizione e di
moltiplicazione, ma non è possibile in generale eseguire le operazioni inverse di
sottrazione e di divisione. $\Z$ è un ampliamento di $\N$ che permette di calcolare
anche le differenze, ma non i quozienti. $\Q$ è un ulteriore ampliamento: in $\Q$ è
possibile eseguire le quattro operazioni fondamentali (tranne, naturalmente, la
divisione per zero), ma non è possibile in generale eseguire altri calcoli altrettanto
utili, come ad esempio l'estrazione di radice. Come vedremo, $\R$ è invece
sufficientemente ricco per la maggior parte delle applicazioni: vi si possono
estrarre radici, calcolare logaritmi, e soprattutto \emph{passare al limite}.
\end{idea}

\begin{python}[Approssimare $\sqrt2$ con razionali, da sotto e da sopra]
# Costruiamo razionali in A (quadrato < 2) e in B (quadrato > 2)
# sempre più vicini tra loro; con Fraction restiamo sempre in Q.
from fractions import Fraction as F

a, b = F(1), F(2)              # a in A (1^2 < 2), b in B (2^2 > 2)
for passo in range(1, 31):
    m = (a + b) / 2            # punto medio: è ancora razionale
    if m * m < 2:
        a = m                  # m sta in A: alziamo a
    else:
        b = m                  # m sta in B: abbassiamo b
    if passo in (1, 2, 3, 10, 20, 30):
        print(f"passo {passo:2d}: {float(a):.12f} < √2 <",
              f"{float(b):.12f}   (b - a = {float(b - a):.2e})")
# La distanza b - a si dimezza a ogni passo e tende a zero, ma
# nessun razionale ha quadrato uguale a 2 (irrazionalità di √2):
print(any(x * x == 2 for x in (a, b)))
\end{python}
\begin{risultato}
passo  1: 1.000000000000 < √2 < 1.500000000000   (b - a = 5.00e-01)
passo  2: 1.250000000000 < √2 < 1.500000000000   (b - a = 2.50e-01)
passo  3: 1.375000000000 < √2 < 1.500000000000   (b - a = 1.25e-01)
passo 10: 1.414062500000 < √2 < 1.415039062500   (b - a = 9.77e-04)
passo 20: 1.414213180542 < √2 < 1.414214134216   (b - a = 9.54e-07)
passo 30: 1.414213561453 < √2 < 1.414213562384   (b - a = 9.31e-10)
False
\end{risultato}

Il programma illustra il contenuto del \cref{teo:Q-non-completo}: le due successioni di
razionali si «stringono» attorno a un punto, ma quel punto non è razionale. Il metodo
usato (dimezzare ripetutamente un intervallo) si chiama \emph{bisezione} e lo
ritroveremo nel capitolo sulle funzioni continue, dove sarà alla base della
dimostrazione del teorema dell'esistenza degli zeri (\cref{teo:zeri}).

\begin{esercizio}
Dimostrare che $\sqrt3$ è irrazionale, adattando la dimostrazione del
\cref{teo:radice2}. (\emph{Suggerimento:} al posto di «pari/dispari» si usi la
divisibilità per $3$: se $m^{2}$ è divisibile per $3$ allora $m$ è divisibile per $3$,
perché se $m=3h\pm1$ allora $m^{2}=3(3h^{2}\pm2h)+1$.)
\begin{soluzione}
Supponiamo per assurdo che esista $c\in\Q$, $c>0$, con $c^{2}=3$, e scriviamo
$c=m/n$ con $m,n\in\N$ non entrambi divisibili per $3$ (basta semplificare la
frazione). Allora $m^{2}=3n^{2}$, e $m^{2}$ è divisibile per $3$. Ogni intero è della
forma $3h$, $3h+1$ oppure $3h-1$; negli ultimi due casi
$m^{2}=9h^{2}\pm6h+1=3(3h^{2}\pm2h)+1$ non è divisibile per $3$. Dunque $m=3k$ con
$k\in\N$, e da $3n^{2}=9k^{2}$ segue $n^{2}=3k^{2}$: per lo stesso motivo anche $n$ è
divisibile per $3$, contro la scelta di $m$ e $n$. \emph{Scelta logica:} la
dimostrazione del \cref{teo:radice2} usa soltanto la proprietà «se $2$ divide $m^{2}$,
allora $2$ divide $m$»; sostituendo la classificazione pari/dispari con quella dei resti
della divisione per $3$, il ragionamento per assurdo si ripete parola per parola.
\end{soluzione}
\end{esercizio}

\begin{esercizio}
Dimostrare che la somma di un numero razionale e di un numero irrazionale è irrazionale.
Che cosa si può dire della somma di due irrazionali?
\begin{soluzione}
Sia $q\in\Q$ e $x\notin\Q$. Se $q+x$ fosse razionale, lo sarebbe anche
$(q+x)-q=x$, perché $\Q$ è chiuso rispetto alla differenza: assurdo. La somma di due
irrazionali può invece essere razionale: $\sqrt2+(-\sqrt2)=0$. \emph{Scelta logica:}
«irrazionale» è una proprietà negativa («non razionale»), che si dimostra per assurdo,
come nel \cref{teo:radice2}; la chiusura di $\Q$ rispetto alla differenza trasforma
l'ipotesi assurda in una contraddizione immediata. Per la seconda domanda, che chiede di
smentire una possibile regola generale, basta un esempio.
\end{soluzione}
\end{esercizio}

\subsection{Approfondimento: l'assioma di Dedekind}
\label{subsec:assioma-dedekind}

L'assioma di completezza si trova enunciato, nei vari testi, in forme diverse. Una delle
più diffuse, adottata ad esempio in \cite{giusti}, risale a Dedekind (1872) e usa la
nozione di \emph{sezione}, cioè di taglio della retta in due parti. Confrontiamola con
la nostra: il confronto è un buon esercizio sull'uso degli assiomi, e mostra che una
stessa proprietà si può esprimere in modi che a prima vista sembrano diversi.

\begin{definizione}[Sezione di $\R$]\label{def:sezione}\index{sezione!di $\R$}
Una coppia $(A,B)$ di sottoinsiemi di $\R$ si chiama \emph{sezione} di $\R$ se
\begin{enumerate}[label=\textup{(\roman*)}]
  \item $A\neq\emptyset$ e $B\neq\emptyset$;
  \item $A\cup B=\R$ e $A\cap B=\emptyset$;
  \item $a<b$ per ogni $a\in A$ e per ogni $b\in B$.
\end{enumerate}
Un numero $L$ tale che $a\le L\le b$ per ogni $a\in A$ e per ogni $b\in B$ si chiama
\emph{elemento separatore} della sezione.
\end{definizione}

L'\emph{assioma di Dedekind}\index{assioma di Dedekind} afferma:
\begin{center}
\textbf{(D)}\enspace \emph{ogni sezione di $\R$ ha uno e un solo elemento separatore.}
\end{center}
Le due formulazioni non dicono la stessa cosa. Il nostro \cref{ass:completezza} si
applica a una coppia qualsiasi di insiemi non vuoti con $a\le b$: gli insiemi non devono
esaurire $\R$ e possono perfino avere un elemento in comune, come $A=B=\{0\}$; in cambio,
afferma soltanto l'\emph{esistenza} di un separatore. L'assioma (D) si applica solo alle
sezioni, ma ne afferma anche l'\emph{unicità}, che in generale è falsa: per $A=[0,1]$ e
$B=[2,3]$ ogni $c\in[1,2]$ è un separatore. Nessuna delle due formulazioni sembra
contenere l'altra; eppure, in presenza degli altri assiomi, si equivalgono.

\begin{teorema}[Equivalenza con l'assioma di Dedekind]\label{teo:dedekind-equivalente}
Supponiamo che valgano gli assiomi \textup{(A1)--(A6)} e \textup{(O1)--(O4)}. Allora
l'assioma di completezza (\cref{ass:completezza}) e l'assioma \textup{(D)} sono
equivalenti: ciascuno dei due implica l'altro.
\end{teorema}

\begin{schema}
\Ipotesi{Gli assiomi (A1)--(A6) e (O1)--(O4).}
\Tesi{(i) completezza $\implica$ (D); \quad (ii) (D) $\implica$ completezza.}
\end{schema}

\begin{strategia}
Entrambe le implicazioni si reggono su un'unica osservazione: \emph{tra due numeri
distinti $u<v$ ce n'è sempre un terzo}, il punto medio $\frac{u+v}{2}$. (i) Una sezione è
una particolare coppia di insiemi separati, dunque il separatore esiste per il nostro
assioma; se ce ne fossero due, $c<c'$, il loro punto medio non potrebbe stare né in $A$
né in $B$, contro $A\cup B=\R$. (ii) Dati $A$ e $B$ come nel nostro assioma, li
«completiamo» a una sezione: a destra mettiamo \emph{tutti} i numeri che sono $\ge$ di
ogni elemento di $A$, a sinistra tutti gli altri. L'assioma (D) fornisce un separatore
$L$ di questa sezione; resta da controllare che $L$ separi anche $A$ e $B$, e il punto
medio servirà ancora per escludere che qualche elemento di $A$ superi $L$.
\end{strategia}

\begin{dimostrazione}
Premettiamo il fatto annunciato: \emph{se $u<v$ e $m=\frac{u+v}{2}$, allora $u<m<v$}.
Infatti $2=1+1\ge0+1=1>0$ (per \ref{ass:ordine-somma} e l'\cref{es:zero-minore-uno}),
quindi $\frac12>0$ (\cref{es:inverso-positivo}); inoltre $m-u=v-m=\frac{v-u}{2}$ è il
prodotto dei numeri positivi $v-u$ e $\frac12$, che è $\ge0$ per
\ref{ass:ordine-prodotto} e diverso da zero per la \cref{prop:annullamento}. Dunque
$m-u>0$ e $v-m>0$.

(i) Supponiamo che valga l'assioma di completezza e sia $(A,B)$ una sezione di $\R$. Gli
insiemi $A$ e $B$ sono non vuoti e $a<b$, quindi $a\le b$, per ogni $a\in A$ e $b\in B$:
per il \cref{ass:completezza} esiste un elemento separatore $c$. Se $c'$ fosse un secondo
separatore diverso da $c$, potremmo supporre (dicotomia) $c<c'$, e il punto medio $m$ di
$c$ e $c'$ verificherebbe $c<m<c'$. Poiché $A\cup B=\R$, $m$ sta in $A$ oppure in $B$.
Se $m\in A$, essendo $c$ un separatore, $m\le c$, contro $c<m$; se $m\in B$, essendo $c'$
un separatore, $c'\le m$, contro $m<c'$. Dunque il separatore è unico: vale (D).

(ii) Supponiamo che valga (D) e siano $A$, $B$ non vuoti con $a\le b$ per ogni $a\in A$ e
$b\in B$. Poniamo
\[
  B^{*}=\{x\in\R : a\le x \ \perogni a\in A\}, \qquad A^{*}=\R\setminus B^{*}.
\]
Per la dicotomia, $x\in A^{*}$ significa che esiste $a\in A$ con $x<a$. Verifichiamo che
$(A^{*},B^{*})$ è una sezione. $B^{*}$ contiene $B$, dunque non è vuoto; $A^{*}$ non è
vuoto, perché se $a_{0}\in A$ allora $a_{0}-1<a_{0}$, cioè $a_{0}-1\in A^{*}$. La
proprietà (ii) della \cref{def:sezione} vale per costruzione. Infine, se $x\in A^{*}$ e
$y\in B^{*}$, esiste $a\in A$ con $x<a$, ed è $a\le y$ perché $y\in B^{*}$: dunque $x<y$.

Per (D) la sezione $(A^{*},B^{*})$ ha un elemento separatore $L$, cioè $x\le L\le y$ per
ogni $x\in A^{*}$ e $y\in B^{*}$. In particolare $L\le b$ per ogni $b\in B$, perché
$B\subseteq B^{*}$. Resta da provare che $a\le L$ per ogni $a\in A$. Se fosse $L<a$ per
qualche $a\in A$, il punto medio $m$ di $L$ e $a$ verificherebbe $L<m<a$; ma $m<a$
significa $m\in A^{*}$, e quindi $m\le L$: assurdo. Dunque $a\le L\le b$ per ogni $a\in
A$ e $b\in B$, e vale l'assioma di completezza.
\end{dimostrazione}

\begin{marginfigure}
\centering
\begin{tikzpicture}[scale=0.85, every node/.style={font=\scriptsize}]
  \foreach \r/\col/\tit in {0/blu/{$L=\max A$}, 1/rame/{$L=\min B$}} {
  \begin{scope}[yshift=-\r*1.25cm]
    \draw[blu, very thick] (0,0) -- (2.4,0);
    \draw[rame, very thick] (2.4,0) -- (4.8,0);
    \fill[\col] (2.4,0) circle (2.4pt);
    \node[above] at (2.4,0.1) {\tit};
    \node[blu, below] at (1.2,-0.06) {$A$};
    \node[rame, below] at (3.6,-0.06) {$B$};
  \end{scope}}
  \begin{scope}[yshift=-2.5cm]
    \draw[blu, very thick, dotted] (0,0) -- (2.3,0);
    \draw[rame, very thick, dotted] (2.5,0) -- (4.8,0);
    \draw[rosso, thick] (2.4,0) circle (2.4pt);
    \node[above] at (2.4,0.1) {in $\Q$: nessun separatore};
    \node[rosso, below] at (2.4,-0.06) {$\sqrt2$};
  \end{scope}
\end{tikzpicture}
\caption{Le due forme di una sezione di $\R$: il punto di taglio $L$ sta in $A$ (in
alto) oppure in $B$ (al centro). In basso, la sezione di $\Q$ della
\eqref{eq:AB-radice2}: il taglio cade in un «buco».}
\label{fig:sezioni}
\end{marginfigure}

\begin{esempio}[Come è fatta una sezione di $\R$]\label{es:sezioni-R}
Sia $(A,B)$ una sezione di $\R$ e sia $L$ il suo elemento separatore. Se $x<L$, allora
$x\notin B$, perché ogni elemento di $B$ è $\ge L$, e dunque $x\in A$; allo stesso modo,
se $x>L$, allora $x\in B$. Il numero $L$ appartiene a uno solo dei due insiemi, e quindi
le sezioni di $\R$ sono soltanto di due tipi:
\begin{align*}
  A&=\{x\le L\},\ \ B=\{x>L\}, && \text{e allora } L=\max A;\\
  A&=\{x<L\},\ \ B=\{x\ge L\}, && \text{e allora } L=\min B
\end{align*}
(\cref{fig:sezioni}). Ogni taglio della retta reale avviene \emph{in un punto}. In $\Q$
le cose vanno diversamente: gli insiemi \eqref{eq:AB-radice2} formano una sezione di
$\Q$, come abbiamo visto nella dimostrazione del \cref{teo:Q-non-completo}, ma $A$ non ha
massimo e $B$ non ha minimo: il taglio non passa per alcun numero razionale.
\end{esempio}

\begin{osservazione}[Il separatore come «primo dei maggioranti»]
Nella parte (ii) della dimostrazione, $B^{*}$ è l'insieme dei numeri che sono $\ge$ di
tutti gli elementi di $A$: nel prossimo capitolo li chiameremo \emph{maggioranti} di $A$.
Il separatore $L$ è $\ge$ di ogni elemento di $A$, dunque appartiene a $B^{*}$, e non
supera alcun elemento di $B^{*}$: è il più piccolo dei maggioranti, cioè l'\emph{estremo
superiore} di $A$, lo strumento con cui useremo la completezza in tutto il libro
(\cref{teo:esistenza-sup}).
\end{osservazione}

\subsection{Approfondimento: un modello dei numeri reali}
\label{subsec:modello-reali}

Nel nostro approccio l'esistenza di $\R$ è un postulato. Dedekind, nel 1872, percorse la
strada opposta: partendo dai soli numeri razionali \emph{costruì} un insieme che
soddisfa tutti gli assiomi del \cref{sec:assiomi}. Ne diamo un'idea, seguendo
\cite{giusti}, e dimostriamo per intero la proprietà più importante, la completezza. In
questo sottoparagrafo dimentichiamo dunque $\R$ e supponiamo di conoscere soltanto $\Q$,
con le sue operazioni e il suo ordinamento, che soddisfano (A1)--(A6) e (O1)--(O4).
L'idea guida viene dall'\cref{es:sezioni-R}: un numero reale è individuato dal taglio
che produce sui razionali; allora \emph{definiamo} un numero reale come un taglio di
$\Q$.

\begin{definizione}[Sezione di $\Q$]\label{def:sezione-Q}\index{sezione!di $\Q$}
Una \emph{sezione} di $\Q$ è una coppia $(A,B)$ di sottoinsiemi non vuoti di $\Q$ tali
che $A\cup B=\Q$, $A\cap B=\emptyset$, $a<b$ per ogni $a\in A$ e $b\in B$, e tali che
$A$ \emph{non abbia massimo}. Indichiamo con $\widetilde{\R}$ l'insieme di tutte le
sezioni di $\Q$, e diciamo che $\alpha=(A,B)$ precede $\beta=(C,D)$, e scriviamo
$\alpha\preceq\beta$, se $A\subseteq C$.
\end{definizione}

La richiesta che $A$ non abbia massimo serve a evitare doppioni: il taglio prodotto da
un razionale $r$ si può fare mettendo $r$ a sinistra o a destra, e conveniamo di
metterlo sempre a destra. Si noti che $B=\Q\setminus A$: una sezione è determinata dalla
sua classe inferiore $A$, e per questo l'ordinamento si definisce confrontando le classi
inferiori.

\begin{esempio}[Sezioni razionali e irrazionali]\label{es:sezioni-Q}
Per ogni $r\in\Q$ la coppia
\[
  \alpha_{r}=\bigl(\{q\in\Q : q<r\},\ \{q\in\Q : q\ge r\}\bigr)
\]
è una sezione di $\Q$: la classe inferiore non ha massimo, perché se $q<r$ il punto
medio di $q$ e $r$, che è razionale, sta ancora a sinistra di $r$ ed è maggiore di $q$.
In $\alpha_{r}$ il numero $r$ è il minimo della classe superiore. Invece la coppia
$(A,B)$ della \eqref{eq:AB-radice2} è una sezione di $\Q$ in cui $A$ non ha massimo e $B$
non ha minimo (lo abbiamo visto nella dimostrazione del \cref{teo:Q-non-completo}):
non è della forma $\alpha_{r}$ per alcun razionale $r$. Nel modello, questa sezione è il
numero «nuovo» $\sqrt2$.
\end{esempio}

\begin{proposizione}[L'ordinamento del modello]\label{prop:ordine-modello}
La relazione $\preceq$ è una relazione d'ordine totale in $\widetilde{\R}$.
\end{proposizione}

\begin{schema}
\Ipotesi{$\alpha=(A,B)$, $\beta=(C,D)$, $\gamma=(E,F)$ elementi di $\widetilde{\R}$.}
\Tesi{$\alpha\preceq\alpha$; \ se $\alpha\preceq\beta$ e $\beta\preceq\alpha$, allora
$\alpha=\beta$; \ se $\alpha\preceq\beta$ e $\beta\preceq\gamma$, allora
$\alpha\preceq\gamma$; \ $\alpha\preceq\beta$ oppure $\beta\preceq\alpha$.}
\end{schema}

\begin{strategia}
Le proprietà riflessiva, antisimmetrica e transitiva sono quelle dell'inclusione tra
insiemi. L'unica che richiede un'idea è la dicotomia: due sottoinsiemi di $\Q$ in
generale non sono confrontabili per inclusione (si pensi a $\{1\}$ e $\{2\}$). Le classi
inferiori, però, sono «chiuse verso sinistra»: con ogni loro elemento contengono tutti i
razionali più piccoli. Se allora $A\not\subseteq C$ e $a\in A\setminus C$, l'intera
classe $C$ sta a sinistra di $a$, e quindi dentro $A$.
\end{strategia}

\begin{dimostrazione}
La riflessività e la transitività seguono da quelle dell'inclusione. Se $A\subseteq C$ e
$C\subseteq A$, allora $A=C$ e quindi anche $B=\Q\setminus A=\Q\setminus C=D$: dunque
$\alpha=\beta$.

Proviamo preliminarmente che, se $a\in A$ e $q\in\Q$ con $q<a$, allora $q\in A$.
Infatti, se fosse $q\in B$, sarebbe $a<q$ per la definizione di sezione, contro $q<a$; e
poiché $A\cup B=\Q$, deve essere $q\in A$.

Per la dicotomia, supponiamo che non sia $A\subseteq C$ e sia $a\in A$ con $a\notin C$,
cioè $a\in D$. Per ogni $c\in C$ si ha $c<a$, perché $c\in C$ e $a\in D$; per quanto
appena provato, $c\in A$. Dunque $C\subseteq A$, cioè $\beta\preceq\alpha$.
\end{dimostrazione}

\begin{teorema}[Completezza del modello]\label{teo:completezza-modello}
Siano $X$ e $Y$ due sottoinsiemi non vuoti di $\widetilde{\R}$ tali che
$\alpha\preceq\beta$ per ogni $\alpha\in X$ e per ogni $\beta\in Y$. Allora esiste
$\lambda\in\widetilde{\R}$ tale che $\alpha\preceq\lambda\preceq\beta$ per ogni
$\alpha\in X$ e per ogni $\beta\in Y$.
\end{teorema}

\begin{schema}
\Ipotesi{$X,Y\subseteq\widetilde{\R}$ non vuoti; \ $\alpha\preceq\beta$ per ogni
$\alpha\in X$, $\beta\in Y$.}
\Tesi{Esiste $\lambda\in\widetilde{\R}$ con $\alpha\preceq\lambda\preceq\beta$ per ogni
$\alpha\in X$, $\beta\in Y$.}
\end{schema}

\begin{strategia}
Poiché l'ordine è l'inclusione delle classi inferiori, il candidato naturale per la classe
inferiore di $\lambda$ è l'\emph{unione} $A_{0}$ delle classi inferiori degli elementi
di $X$: contiene ciascuna di esse, quindi $\lambda$ seguirà ogni $\alpha\in X$; ed è
contenuta nella classe inferiore di ogni $\beta\in Y$, perché questa contiene già tutte
le classi che stiamo unendo. Bisogna soltanto verificare che $A_{0}$ e il suo
complementare formino una sezione di $\Q$: ciascuna proprietà di $(A_{0},\Q\setminus
A_{0})$ discende dalla proprietà corrispondente delle sezioni di $X$ e di $Y$.
\end{strategia}

\begin{dimostrazione}
Poniamo
\[
  A_{0}=\bigcup_{(A,B)\in X} A, \qquad B_{0}=\Q\setminus A_{0},
\]
cioè $q\in A_{0}$ se e solo se $q\in A$ per almeno una sezione $(A,B)\in X$. Fissiamo
$\beta=(C,D)\in Y$: per ipotesi $A\subseteq C$ per ogni $(A,B)\in X$, e quindi
$A_{0}\subseteq C$. Verifichiamo che $\lambda=(A_{0},B_{0})$ è una sezione di $\Q$.
\begin{itemize}
  \item $A_{0}\neq\emptyset$, perché $X$ non è vuoto e le classi inferiori non sono vuote;
  $B_{0}\neq\emptyset$, perché $A_{0}\subseteq C$ implica $D=\Q\setminus C\subseteq
  B_{0}$, e $D\neq\emptyset$.
  \item $A_{0}\cup B_{0}=\Q$ e $A_{0}\cap B_{0}=\emptyset$ per costruzione.
  \item Siano $a\in A_{0}$ e $b\in B_{0}$. Esiste $(A,B)\in X$ con $a\in A$; poiché
  $b\notin A_{0}$, è $b\notin A$, cioè $b\in B$, e dunque $a<b$.
  \item $A_{0}$ non ha massimo: se $a\in A_{0}$, allora $a\in A$ per qualche $(A,B)\in
  X$; $A$ non ha massimo, dunque esiste $a'\in A$ con $a'>a$, e $a'\in A_{0}$.
\end{itemize}
Infine, $A\subseteq A_{0}$ per ogni $(A,B)\in X$, cioè $\alpha\preceq\lambda$ per ogni
$\alpha\in X$; e $A_{0}\subseteq C$ per ogni $(C,D)\in Y$, cioè $\lambda\preceq\beta$ per
ogni $\beta\in Y$.
\end{dimostrazione}

\begin{osservazione}[Le operazioni nel modello]
Per completare la costruzione bisogna definire in $\widetilde{\R}$ la somma e il prodotto
e verificare gli assiomi (A1)--(A6) e (O1)--(O4). La somma si definisce in modo naturale:
se $\alpha=(A,B)$ e $\beta=(C,D)$, si pone $\alpha+\beta=(A+C,\ \Q\setminus(A+C))$, dove
$A+C=\{a+c : a\in A,\ c\in C\}$, e lo zero è $\alpha_{0}$. Il prodotto è più delicato a
causa della regola dei segni: per due sezioni «positive» si moltiplicano gli elementi
delle classi superiori, non quelli delle classi inferiori, che contengono numeri negativi
grandi in valore assoluto, con prodotti positivi e grandissimi. Le verifiche sono
elementari ma lunghe, e le omettiamo. Si prova infine che
$\alpha_{r}+\alpha_{s}=\alpha_{r+s}$, che $\alpha_{r}\cdot\alpha_{s}=\alpha_{rs}$ e che
$\alpha_{r}\preceq\alpha_{s}$ se e solo se $r\le s$: la corrispondenza $r\mapsto\alpha_{r}$
«immerge» $\Q$ nel modello senza alterarne operazioni e ordinamento, e la sezione
dell'\cref{es:sezioni-Q} risulta avere quadrato $\alpha_{2}$.
\end{osservazione}

\begin{osservazione}[Costruire o postulare?]
Che cosa si guadagna costruendo $\widetilde{\R}$? Weierstrass e i matematici della sua
generazione giudicavano la costruzione più soddisfacente del metodo assiomatico, perché
riconduce l'esistenza dei reali a quella, ritenuta più sicura, dei razionali e quindi
dei naturali (si veda la nota storica del \cref{sec:assiomi}). Ma in matematica
«esistere» significa, in ultima analisi, «non essere contraddittorio», e la costruzione
fa largo uso della teoria degli insiemi: ogni sezione è una coppia di insiemi infiniti di
razionali, e $\widetilde{\R}$ è un insieme di tali coppie. Ciò che si ottiene è dunque una
garanzia \emph{relativa}: se non sono contraddittori gli assiomi dei naturali e quelli
della teoria degli insiemi, non lo sono neppure gli assiomi di $\R$. Vale anche il
viceversa, perché dentro $\R$ c'è un insieme, $\N$, che soddisfa gli assiomi di Peano
(\cref{oss:peano}). Una garanzia assoluta non è raggiungibile: nel 1931 Kurt Gödel
(1906--1978) dimostrò che la non contraddittorietà dell'aritmetica non si può provare con
i soli mezzi dell'aritmetica stessa. La scelta del punto di partenza è quindi, in larga
misura, una questione di convenienza; per l'Analisi, che studia il continuo, è naturale
partire dagli assiomi dei numeri reali \cite{giusti}.
\end{osservazione}

%% ----------------------------------------------------------------------------
%%  §1.5  Funzioni e rappresentazione cartesiana
%% ----------------------------------------------------------------------------
\section{Funzioni e rappresentazione cartesiana}
\label{sec:funzioni}

\begin{definizione}[Funzione]\index{funzione}\index{dominio}\index{immagine}
Siano $A$ e $B$ due insiemi. Una \emph{funzione} di $A$ in $B$ è una legge che a ogni
elemento di $A$ fa corrispondere \emph{uno e un solo} elemento di $B$. Se indichiamo con
$f$ tale funzione, scriveremo
\[
  f\colon A\to B, \qquad\text{oppure}\qquad y=f(x),
\]
intendendo che a ogni elemento $x\in A$ corrisponde, tramite la funzione $f$, l'elemento
$y=f(x)\in B$. Si dice che $A$ è il \emph{dominio} o \emph{insieme di definizione} di
$f$; il valore $f(x)$ si chiama anche \emph{immagine} di $x$ mediante $f$, e $x$ si
chiama \emph{argomento} o \emph{variabile indipendente}.
\end{definizione}

\begin{marginfigure}
\centering
\begin{tikzpicture}[scale=0.8]
  \draw[blu, thick] (0,0) ellipse (0.9 and 1.4) node[above=1.4cm, font=\scriptsize] {$A$};
  \draw[rame, thick] (3,0) ellipse (0.9 and 1.4) node[above=1.4cm, font=\scriptsize] {$B$};
  \foreach \y/\n in {0.8/1, 0.2/2, -0.4/3, -1.0/4} \fill[blu] (0,\y) circle (1.6pt) node[left, font=\scriptsize] {$x_{\n}$};
  \foreach \y/\n in {0.9/1, 0.3/2, -0.3/3, -0.9/4} \fill[rame] (3,\y) circle (1.6pt) node[right, font=\scriptsize] {$y_{\n}$};
  \draw[-Latex, grigio] (0.1,0.8) -- (2.9,0.3);
  \draw[-Latex, grigio] (0.1,0.2) -- (2.9,0.3);
  \draw[-Latex, grigio] (0.1,-0.4) -- (2.9,-0.9);
  \draw[-Latex, grigio] (0.1,-1.0) -- (2.9,0.9);
  \node[font=\scriptsize, grigio] at (1.5,1.7) {$f$};
\end{tikzpicture}
\caption{Una funzione $f\colon A\to B$: da ogni elemento di $A$ parte \emph{una sola}
freccia. Due frecce possono arrivare sullo stesso elemento di $B$, e qualche elemento di
$B$ può non essere raggiunto.}
\label{fig:frecce}
\end{marginfigure}

Nel linguaggio degli insiemi del paragrafo precedente, una funzione può essere definita
in modo del tutto rigoroso come un sottoinsieme $G$ del prodotto cartesiano $A\times B$
con la proprietà che per ogni $x\in A$ esiste \emph{una e una sola} coppia $(x,y)\in G$:
quella coppia definisce $y=f(x)$. L'insieme $G$ si chiama \emph{grafico} della
funzione, e la definizione «per frecce» della \cref{fig:frecce} e quella «per grafico»
sono equivalenti.

Il simbolo $f(\ )$ indica un complesso di operazioni che devono effettuarsi su $x$ per
ottenere $y$ (il \emph{valore} di $f$), come negli esempi seguenti:
\begin{align}
  f(x)&=2x+1, \label{eq:f1}\\
  f(x)&=1/x, \label{eq:f2}\\
  f(x)&=\sqrt{x}, \label{eq:f3}\\
  f(x)&=\begin{cases}0 & \text{se } x\in\Z,\\ 1 & \text{altrimenti.}\end{cases}\label{eq:f4}
\end{align}
Nella \eqref{eq:f1} occorre moltiplicare $x$ per $2$ e sommare $1$; nella \eqref{eq:f2}
occorre calcolare l'inverso di $x$; nella \eqref{eq:f3} si deve calcolare la radice
quadrata di $x$; nella \eqref{eq:f4} occorre riconoscere se $x$ è intero oppure no, e di
conseguenza assegnare a $f(x)$ il valore $0$ oppure $1$.
La funzione \eqref{eq:f1} è definita per ogni $x$ reale; in altre parole il suo dominio è
tutto $\R$. La funzione \eqref{eq:f2} è definita per $x\neq0$, quindi il suo dominio è
$A=\{x\in\R : x\neq0\}$. Il dominio della \eqref{eq:f3} è $A=\{x\in\R : x\ge0\}$,
mentre la \eqref{eq:f4} è definita su tutto $\R$. In questo libro, quando una funzione è
assegnata mediante una formula senza che sia specificato il dominio, si intende che il
dominio è il più grande sottoinsieme di $\R$ in cui la formula ha senso (il cosiddetto
\emph{dominio naturale}).

\begin{attenzione}[Una funzione non è una formula]
La \eqref{eq:f4} mostra che una funzione può essere definita «a pezzi», senza una
formula unica; esistono funzioni definite in modo ancora più indiretto, come la funzione
definita per $n\in\N$ da «$f(n)$ è la $n$-esima cifra decimale di $\pi$». Ciò che conta è che a ogni $x$ del dominio
corrisponda \emph{uno e un solo} valore. Per lo stesso motivo la «legge» $y^{2}=x$ non
definisce una funzione di $x$: a $x=4$ corrisponderebbero due valori, $y=2$ e $y=-2$.
\end{attenzione}

\begin{notastorica}[Da Leibniz a Dirichlet: l'idea di funzione]
Nella \emph{Géométrie} (1637) di Cartesio una curva è descritta da un'equazione
$P(x,y)=0$ tra le coordinate dei suoi punti: una \emph{relazione} tra due variabili, non
ancora la dipendenza dell'una dall'altra. La parola \emph{funzione} compare in un
manoscritto di Gottfried Wilhelm Leibniz (1646--1716) del 1673, e poi a stampa negli
anni Novanta del Seicento: per Leibniz erano «funzioni» i segmenti associati ai punti di
una curva, come l'ordinata, la tangente o la normale. Nella sua corrispondenza con
Johann Bernoulli (1667--1748) la parola cominciò a indicare una quantità formata a
partire da una variabile, e nel 1718 Bernoulli ne diede la prima definizione esplicita:
una quantità composta in qualsiasi modo da una grandezza variabile e da costanti.
Leonhard Euler (1707--1783), che rese d'uso comune la notazione $f(x)$, fece della
funzione il concetto centrale dell'Analisi con l'\emph{Introductio in analysin
infinitorum} (1748); per lui, però, una funzione era un'\emph{espressione analitica},
cioè una formula costruita a partire dalla variabile con le operazioni algebriche e le
funzioni elementari (sulle conseguenze di questa idea si veda la nota storica del
\cref{sec:continue}). Nelle \emph{Institutiones calculi differentialis} (1755) Eulero
allargò la definizione a ogni modo in cui una quantità dipende da un'altra, e il distacco
definitivo dalla formula venne con Dirichlet (\cref{oss:dirichlet}): la sua idea di
corrispondenza arbitraria è la «legge» della nostra definizione. La formulazione come
sottoinsieme del prodotto cartesiano, cioè come grafico, si è imposta nel Novecento con
la teoria degli insiemi.
\end{notastorica}

\subsection{Rappresentazione cartesiana}

Ricordiamo brevemente come effettuare la \emph{rappresentazione cartesiana}\index{piano cartesiano}
di una funzione. Riferendoci alla \cref{fig:piano-cartesiano}, consideriamo due rette
perpendicolari che si intersecano in un punto $O$, \emph{origine} degli assi. Fissiamo
sulle due rette una direzione positiva e un'unità di misura. Chiamiamo \emph{asse delle
ascisse}, o asse $x$, una delle due rette, e \emph{asse delle ordinate}, o asse $y$,
l'altra.

A ogni numero reale $x$ corrisponde in modo biunivoco un punto $P_{1}$ dell'asse delle
$x$, scelto in modo che il segmento $OP_{1}$ abbia lunghezza uguale a $\abs{x}$ e stia
dalla parte positiva se $x>0$. Analogamente a ogni numero reale $y$ corrisponde un
punto $P_{2}$ dell'asse $y$. Tracciamo due rette parallele agli assi e passanti
rispettivamente per $P_{1}$ e $P_{2}$. Il punto $P$ di incontro delle due rette
corrisponde in modo biunivoco alla coppia di numeri reali $(x,y)$: il piano «è»
l'insieme $\R\times\R=\R^{2}$ delle coppie ordinate.

\begin{figure}[htb]
\centering
\begin{tikzpicture}[scale=0.9]
  \draw[-Latex, grigio] (-0.8,0) -- (5.5,0) node[below, font=\footnotesize] {$x$};
  \draw[-Latex, grigio] (0,-0.8) -- (0,3.5) node[left, font=\footnotesize] {$y$};
  \node[below left, font=\footnotesize] at (0,0) {$O$};
  \fill[blu] (3.8,2.4) circle (2pt) node[above right, font=\footnotesize] {$P=(x,y)$};
  \draw[tratteggio] (3.8,0) -- (3.8,2.4) -- (0,2.4);
  \fill[rame] (3.8,0) circle (1.6pt) node[below, font=\footnotesize] {$P_{1}$};
  \fill[rame] (0,2.4) circle (1.6pt) node[left, font=\footnotesize] {$P_{2}$};
  \draw[decorate, decoration={brace, amplitude=3pt, mirror}, grigio] (0,-0.35) -- (3.8,-0.35) node[midway, below=3pt, font=\scriptsize] {$x$};
  \draw[decorate, decoration={brace, amplitude=3pt}, grigio] (-0.35,0) -- (-0.35,2.4) node[midway, left=3pt, font=\scriptsize] {$y$};
\end{tikzpicture}
\caption{Il piano cartesiano: al punto $P$ corrisponde la coppia ordinata $(x,y)$ delle
sue coordinate, e viceversa.}
\label{fig:piano-cartesiano}
\end{figure}

\begin{definizione}[Grafico]\index{grafico}
Se è assegnata una funzione $f\colon A\to\R$ con $A\subseteq\R$, l'insieme dei punti
del piano di coordinate $(x,f(x))$, al variare di $x\in A$,
\[
  G_{f}=\{(x,y)\in\R^{2} : x\in A,\ y=f(x)\},
\]
si chiama \emph{grafico} della funzione.
\end{definizione}

Un sottoinsieme $G$ di $A\times\R$ è il grafico di una funzione definita in $A$ se e
solo se ogni retta verticale $x=x_{0}$, con $x_{0}\in A$, lo interseca in
\emph{esattamente un} punto:
è la traduzione geometrica dell'«uno e un solo» della definizione di funzione. Nella
\cref{fig:quattro-grafici} abbiamo riportato i grafici delle quattro funzioni
\eqref{eq:f1}--\eqref{eq:f4}; il lettore verifichi i disegni \emph{per punti}, cioè
assegnando alla variabile indipendente $x$ dei valori su cui calcolare semplicemente la
funzione: ad esempio, nel caso \eqref{eq:f3}, $x=0,1,4,9,\dots$

\begin{figure}[htb]
\centering
\begin{tikzpicture}
\begin{axis}[width=5.6cm, height=4.6cm, xmin=-2.2, xmax=2.2, ymin=-2.5, ymax=4,
  xtick={-1,1}, ytick={1,3}, title={\footnotesize $y=2x+1$}, title style={at={(0.5,0.9)}}]
  \addplot[curva, domain=-2:1.4] {2*x+1};
\end{axis}
\end{tikzpicture}\hfill
\begin{tikzpicture}
\begin{axis}[width=5.6cm, height=4.6cm, xmin=-3, xmax=3, ymin=-3, ymax=3,
  xtick={1}, ytick={1}, title={\footnotesize $y=1/x$}, title style={at={(0.5,0.9)}}]
  \addplot[curva, domain=0.34:3] {1/x};
  \addplot[curva, domain=-3:-0.34] {1/x};
  \addplot[tratteggio] coordinates {(1,0) (1,1) (0,1)};
\end{axis}
\end{tikzpicture}\\[2mm]
\begin{tikzpicture}
\begin{axis}[width=5.6cm, height=4.6cm, xmin=-0.6, xmax=5.4, ymin=-0.6, ymax=2.8,
  xtick={1,2,3,4,5}, ytick={1,2}, title={\footnotesize $y=\sqrt{x}$}, title style={at={(0.5,0.9)}}]
  \addplot[curva, domain=0:5.2, samples=100] {sqrt(x)};
  \foreach \p in {0,1,4} {\addplot[only marks, mark=*, mark size=1.3pt, blu] coordinates {(\p,{sqrt(\p)})};}
\end{axis}
\end{tikzpicture}\hfill
\begin{tikzpicture}
\begin{axis}[width=5.6cm, height=4.6cm, xmin=-3.3, xmax=3.3, ymin=-0.6, ymax=1.8,
  xtick={-2,-1,1,2}, ytick={1}, title={\footnotesize $y=0$ se $x\in\Z$, $y=1$ altrimenti}, title style={at={(0.5,0.9)}}]
  \addplot[curva] coordinates {(-3.3,1) (3.3,1)};
  \foreach \k in {-3,...,3} {
    \addplot[only marks, mark=*, mark size=1.3pt, blu] coordinates {(\k,0)};
    \addplot[only marks, mark=o, mark size=1.3pt, blu, fill=white] coordinates {(\k,1)};}
\end{axis}
\end{tikzpicture}
\caption{I grafici delle quattro funzioni \eqref{eq:f1}--\eqref{eq:f4}. Nell'ultimo, i
cerchietti vuoti indicano punti che \emph{non} appartengono al grafico: sugli interi la
funzione vale $0$, non $1$.}
\label{fig:quattro-grafici}
\end{figure}

\begin{python}[Disegnare un grafico per punti con \texttt{matplotlib}]
# Il calcolatore disegna i grafici proprio "per punti": sceglie
# molti valori di x, calcola f(x) e unisce i punti con segmenti.
import numpy as np
import matplotlib.pyplot as plt

x = np.linspace(0, 5, 200)          # 200 valori di x tra 0 e 5
y = np.sqrt(x)                      # f(x) = radice di x, su tutti

fig, ax = plt.subplots(figsize=(4, 3))
ax.plot(x, y, color="#1F4E79", lw=1.5, label=r"$y=\sqrt{x}$")
ax.plot([0, 1, 4], [0, 1, 2], "o", color="#B4662F")  # punti facili
ax.set_xlabel("x"); ax.set_ylabel("y"); ax.legend()
ax.grid(alpha=0.3)
fig.tight_layout()
fig.savefig("figure/radice.png", dpi=220)
\end{python}

\begin{figure}[htb]
\centering
\includegraphics[width=0.55\textwidth]{radice}
\caption{Il grafico di $y=\sqrt{x}$ prodotto dal programma precedente
(\texttt{python/fig\_radice.py}). I tre punti evidenziati sono quelli calcolabili «a
mente»: $(0,0)$, $(1,1)$, $(4,2)$.}
\label{fig:radice-python}
\end{figure}

\begin{attenzione}[Dove il calcolatore sbaglia]
Se si prova a disegnare con lo stesso metodo la funzione $y=1/x$ su un intervallo che
contiene $0$, il programma unirà con un segmento l'ultimo punto a sinistra di $0$ (con
$y$ molto negativo) e il primo a destra (con $y$ molto positivo), tracciando una linea
verticale che \emph{non fa parte del grafico}. Allo stesso modo, la funzione
\eqref{eq:f4} disegnata «per punti» appare come la retta $y=1$: con $200$ punti
equispaziati è facile non cadere mai esattamente su un intero. Il grafico al calcolatore
è uno strumento prezioso, ma va sempre interpretato alla luce della definizione.
\end{attenzione}

%% ----------------------------------------------------------------------------
%%  §1.6  Funzioni invertibili. Funzioni monotòne
%% ----------------------------------------------------------------------------
\section{Funzioni invertibili. Funzioni monotòne}
\label{sec:invertibili}

\begin{marginfigure}
\centering
\begin{tikzpicture}[scale=0.62, every node/.style={font=\scriptsize}]
  \foreach \k/\tit in {0/{iniettiva, non suriettiva}, 1/{suriettiva, non iniettiva}, 2/{biunivoca}} {
  \begin{scope}[shift={(0,-\k*3.6)}]
    \draw[blu, thick] (0,0) ellipse (0.7 and 1.3);
    \draw[rame, thick] (3,0) ellipse (0.7 and 1.3);
    \node at (1.5,-1.75) {\tit};
    \ifnum\k=0
      \foreach \y in {0.7,0,-0.7} \fill[blu] (0,\y) circle (1.4pt);
      \foreach \y in {0.9,0.3,-0.3,-0.9} \fill[rame] (3,\y) circle (1.4pt);
      \draw[-Latex,grigio] (0.1,0.7) -- (2.9,0.9);
      \draw[-Latex,grigio] (0.1,0) -- (2.9,-0.3);
      \draw[-Latex,grigio] (0.1,-0.7) -- (2.9,-0.9);
    \fi
    \ifnum\k=1
      \foreach \y in {0.9,0.3,-0.3,-0.9} \fill[blu] (0,\y) circle (1.4pt);
      \foreach \y in {0.7,0,-0.7} \fill[rame] (3,\y) circle (1.4pt);
      \draw[-Latex,grigio] (0.1,0.9) -- (2.9,0.7);
      \draw[-Latex,grigio] (0.1,0.3) -- (2.9,0.7);
      \draw[-Latex,grigio] (0.1,-0.3) -- (2.9,0);
      \draw[-Latex,grigio] (0.1,-0.9) -- (2.9,-0.7);
    \fi
    \ifnum\k=2
      \foreach \y in {0.7,0,-0.7} \fill[blu] (0,\y) circle (1.4pt);
      \foreach \y in {0.7,0,-0.7} \fill[rame] (3,\y) circle (1.4pt);
      \draw[-Latex,grigio] (0.1,0.7) -- (2.9,0);
      \draw[-Latex,grigio] (0.1,0) -- (2.9,-0.7);
      \draw[-Latex,grigio] (0.1,-0.7) -- (2.9,0.7);
    \fi
  \end{scope}}
\end{tikzpicture}
\caption{Iniettiva: frecce che non si incontrano mai. Suriettiva: ogni elemento di $B$
è raggiunto. Biunivoca: entrambe le cose, e le frecce si possono «invertire».}
\label{fig:iniettiva-suriettiva}
\end{marginfigure}

\begin{definizione}[Funzioni iniettive, suriettive, biunivoche]\index{funzione!iniettiva}\index{funzione!suriettiva}\index{funzione!biunivoca}
Sia $f\colon A\to B$ una funzione.
\begin{itemize}
  \item $f$ si dice \emph{iniettiva} se elementi distinti hanno immagini distinte, cioè,
  equivalentemente, se per ogni $x_{1},x_{2}\in A$ sussiste l'implicazione
  \begin{equation}\label{eq:iniettiva}
    f(x_{1})=f(x_{2}) \ \implica\ x_{1}=x_{2}.
  \end{equation}
  \item $f$ si dice \emph{suriettiva} (da $A$ \emph{verso} $B$) se per ogni $y\in B$
  esiste almeno un $x\in A$ tale che $y=f(x)$.
  \item $f$ si dice \emph{biunivoca} (o \emph{biiettiva}) se è contemporaneamente
  iniettiva e suriettiva. Ciò vuol dire che $f$ non solo fa corrispondere a ogni $x\in A$
  uno e un solo valore $y\in B$, ma anche che per ogni $y\in B$ esiste \emph{un solo}
  $x\in A$ tale che $y=f(x)$.
\end{itemize}
\end{definizione}


\begin{definizione}[Funzione inversa]\index{funzione!inversa}
Se $f\colon A\to B$ è biunivoca diciamo che $f$ è \emph{invertibile}. La funzione da $B$
ad $A$ che a ogni $y\in B$ fa corrispondere l'unico $x\in A$ per cui $f(x)=y$ si chiama
\emph{funzione inversa} e si indica con $f^{-1}$. Si ha (il simbolo $\perogni$ si legge
\emph{per ogni})
\begin{equation}\label{eq:inversa}
  f^{-1}(f(x))=x \quad\perogni x\in A; \qquad\qquad f(f^{-1}(y))=y \quad\perogni y\in B.
\end{equation}
\end{definizione}

Da notare che spesso si cambiano le notazioni e, invece di usare $x=f^{-1}(y)$, si
preferisce --- con un puro scambio di simboli --- scrivere $y=f^{-1}(x)$, per mettere in
luce che la variabile indipendente è quella dell'argomento di $f^{-1}$. Con questa
convenzione il grafico di $f^{-1}$ si ottiene da quello di $f$ scambiando il ruolo degli
assi, cioè con una simmetria rispetto alla bisettrice $y=x$ (\cref{fig:inversa}).

\begin{figure}[htb]
\centering
\begin{tikzpicture}
\begin{axis}[width=6.2cm, height=6.2cm, xmin=-0.3, xmax=3.2, ymin=-0.3, ymax=3.2,
  xtick=\empty, ytick=\empty, title={\footnotesize $y=f(x)$}]
  \addplot[tratteggio, domain=-0.3:3.2] {x};
  \addplot[curva, domain=0.5:2.5, samples=80] {0.9+0.5*(x-0.5)+0.275*(x-0.5)^2};
  \addplot[tratteggio] coordinates {(0.5,0) (0.5,0.9) (0,0.9)};
  \addplot[tratteggio] coordinates {(2.5,0) (2.5,3.0) (0,3.0)};
  \draw[blu, thick, decorate, decoration={brace, amplitude=3pt, mirror}] (axis cs:0.5,-0.15) -- (axis cs:2.5,-0.15) node[midway, below=3pt, font=\scriptsize] {$A$};
  \draw[rame, thick, decorate, decoration={brace, amplitude=3pt}] (axis cs:-0.15,0.9) -- (axis cs:-0.15,3.0) node[midway, left=3pt, font=\scriptsize] {$B$};
\end{axis}
\end{tikzpicture}\hfill
\begin{tikzpicture}
\begin{axis}[width=6.2cm, height=6.2cm, xmin=-0.3, xmax=3.2, ymin=-0.3, ymax=3.2,
  xtick=\empty, ytick=\empty, title={\footnotesize $y=f^{-1}(x)$}]
  \addplot[tratteggio, domain=-0.3:3.2] {x};
  % inversa: scambio delle coordinate, disegnata in forma parametrica
  \addplot[curva2, domain=0.5:2.5, samples=80] ({0.9+0.5*(x-0.5)+0.275*(x-0.5)^2},{x});
  \addplot[tratteggio] coordinates {(0.9,0) (0.9,0.5) (0,0.5)};
  \addplot[tratteggio] coordinates {(3.0,0) (3.0,2.5) (0,2.5)};
  \draw[rame, thick, decorate, decoration={brace, amplitude=3pt, mirror}] (axis cs:0.9,-0.15) -- (axis cs:3.0,-0.15) node[midway, below=3pt, font=\scriptsize] {$B$};
  \draw[blu, thick, decorate, decoration={brace, amplitude=3pt}] (axis cs:-0.15,0.5) -- (axis cs:-0.15,2.5) node[midway, left=3pt, font=\scriptsize] {$A$};
\end{axis}
\end{tikzpicture}
\caption{Il grafico di una funzione invertibile $f\colon A\to B$ e quello della sua
inversa $f^{-1}\colon B\to A$: l'uno è il simmetrico dell'altro rispetto alla
bisettrice del primo e terzo quadrante (tratteggiata).}
\label{fig:inversa}
\end{figure}

\begin{esempio}[Le funzioni \eqref{eq:f1}--\eqref{eq:f4} sono invertibili?]
La funzione \eqref{eq:f1}, $f(x)=2x+1$, è invertibile da $\R$ in $\R$: fissato
$y=f(x)$, risulta $x=(y-1)/2$, e questo $x$ è unico. La funzione inversa è quindi
\begin{equation}\label{eq:inversa-f1}
  f^{-1}(x)=\frac{x-1}{2}.
\end{equation}
Anche la funzione \eqref{eq:f2}, $f(x)=1/x$, è invertibile da $\R\setminus\{0\}$ in
$\R\setminus\{0\}$ e risulta $f^{-1}(x)=1/x$: è solo un caso che $f$ e $f^{-1}$
coincidano! La funzione \eqref{eq:f3} è invertibile da $[0,+\infty)$ in $[0,+\infty)$
e $f^{-1}(x)=x^{2}$, per $x\ge0$. La \eqref{eq:f4} non è invertibile, perché non
stabilisce una corrispondenza biunivoca tra $\R$ e l'insieme $B=\{0,1\}$ dei valori:
tutti gli interi, per esempio, hanno la stessa immagine $0$.
\end{esempio}

\begin{attenzione}[$f^{-1}$ non è $1/f$]
Il simbolo $f^{-1}$ indica la funzione inversa e \emph{non} il reciproco $1/f$. Per
$f(x)=2x+1$ si ha $f^{-1}(x)=(x-1)/2$, mentre $1/f(x)=1/(2x+1)$: funzioni del tutto
diverse. La notazione è infelice ma universale, e vale la pena di abituarsi subito.
\end{attenzione}

\begin{notastorica}[Cantor: tanti punti in un segmento quanti in un quadrato]
Le corrispondenze biunivoche sono lo strumento con cui Georg Cantor, a partire dal
1874, imparò a confrontare gli insiemi infiniti: due insiemi hanno la stessa
\emph{potenza} --- per gli insiemi finiti, lo stesso numero di elementi --- quando
esiste una funzione biunivoca dall'uno all'altro. Nel 1877 Cantor si chiese se il
quadrato $[0,1]\times[0,1]$ abbia «più punti» del segmento $[0,1]$, come tutti davano
per scontato, e trovò invece una corrispondenza biunivoca tra i due. L'idea, nella forma
più semplice, è di intrecciare le cifre decimali: al punto
$(0{,}a_{1}a_{2}a_{3}\dots,\ 0{,}b_{1}b_{2}b_{3}\dots)$ si associa il numero
$0{,}a_{1}b_{1}a_{2}b_{2}a_{3}b_{3}\dots$. Dedekind, a cui Cantor aveva sottoposto la
dimostrazione, osservò che, per via dei numeri con due scritture decimali (come
$0{,}1=0{,}0999\dots$), una corrispondenza costruita così è iniettiva ma non suriettiva;
Cantor rimediò con un argomento più elaborato e il 29 giugno 1877 gli scrisse: «Je le
vois, mais je ne le crois pas!» («Lo vedo, ma non ci credo!»). Il risultato, pubblicato
nel 1878, sembrava togliere significato alla nozione di dimensione; la questione fu
chiarita solo nel 1911, quando L.~E.~J. Brouwer dimostrò che tra spazi di dimensione
diversa non esistono corrispondenze biunivoche continue insieme con la loro inversa.
\end{notastorica}

\subsection{Funzioni monotòne}

\begin{definizione}[Funzioni monotòne]\index{funzione!monotòna}\index{funzione!crescente}\index{funzione!decrescente}
Diremo che una funzione $f$ è \emph{monotòna} in un insieme $A$ se verifica una delle
condizioni seguenti, per ogni $x_{1},x_{2}\in A$:
\begin{align}
  &f \text{ \emph{strettamente crescente}:}   & x_{1}<x_{2} &\implica f(x_{1})<f(x_{2}),\label{eq:str-cresc}\\
  &f \text{ \emph{crescente}:}                & x_{1}<x_{2} &\implica f(x_{1})\le f(x_{2}),\label{eq:cresc}\\
  &f \text{ \emph{strettamente decrescente}:} & x_{1}<x_{2} &\implica f(x_{1})>f(x_{2}),\label{eq:str-decr}\\
  &f \text{ \emph{decrescente}:}              & x_{1}<x_{2} &\implica f(x_{1})\ge f(x_{2}).\label{eq:decr}
\end{align}
Una funzione che verifica la \eqref{eq:str-cresc} oppure la \eqref{eq:str-decr} si
dice \emph{strettamente monotòna}.
\end{definizione}

\begin{marginfigure}
\centering
\begin{tikzpicture}
\begin{axis}[width=5cm, height=4cm, xmin=0, xmax=4, ymin=0, ymax=3, xtick=\empty, ytick=\empty]
  \addplot[curva, domain=0.2:3.8, samples=60] {0.5+0.9*(x-0.2)^0.6};
  \addplot[curva2] coordinates {(0.2,0.6) (1,1.4) (2,1.4) (3,2.2) (3.8,2.2)};
\end{axis}
\end{tikzpicture}
\caption{In blu una funzione strettamente crescente; in rame una funzione crescente
ma non strettamente: ha dei «pianerottoli» su cui $f(x_1)=f(x_2)$ con $x_1<x_2$.}
\end{marginfigure}

Ad esempio, la funzione \eqref{eq:f1} è strettamente crescente su $\R$; la funzione
\eqref{eq:f2} è strettamente decrescente, \emph{separatamente} negli insiemi $\{x>0\}$ e
$\{x<0\}$ (ma non su tutto il suo dominio: $-1<1$ e $f(-1)=-1<1=f(1)$). La funzione
\eqref{eq:f3} è strettamente crescente. La funzione \eqref{eq:f4} non è monotòna su
$\R$. Una funzione costante è contemporaneamente crescente e decrescente, ma non
strettamente.

Ci sarà utile nel seguito un \emph{criterio} per riconoscere se una data funzione è
invertibile. Il risultato che segue è semplice, ma è uno dei più usati del libro: è
grazie a esso che potremo definire le radici, i logaritmi e le funzioni trigonometriche
inverse.

\begin{teorema}[Criterio di invertibilità]\label{teo:monotona-invertibile}\index{criterio di invertibilità}
Sia $f\colon A\to B$ una funzione suriettiva. Se $f$ è strettamente monotòna in $A$,
allora $f$ è invertibile, e l'inversa $f^{-1}\colon B\to A$ è strettamente monotòna
dello stesso tipo (crescente se $f$ è crescente, decrescente se $f$ è decrescente).
\end{teorema}

\begin{schema}
\Ipotesi{$f\colon A\to B$ suriettiva e strettamente monotòna.}
\Tesi{$f$ è iniettiva (dunque biunivoca, dunque invertibile) e $f^{-1}$ è strettamente
monotòna dello stesso tipo.}
\end{schema}

\begin{strategia}
Poiché $f$ è suriettiva per ipotesi, per avere l'invertibilità manca solo
l'iniettività: dobbiamo mostrare che due punti distinti $x_{1}\neq x_{2}$ hanno immagini
distinte. Ma se $x_{1}\neq x_{2}$, uno dei due è più piccolo, diciamo $x_{1}<x_{2}$
(dicotomia); e allora la stretta monotonia dice che $f(x_{1})<f(x_{2})$ oppure
$f(x_{1})>f(x_{2})$: in ogni caso $f(x_{1})\neq f(x_{2})$. Per la monotonia
dell'inversa ragioniamo per assurdo: se $y_{1}<y_{2}$ ma $f^{-1}(y_{1})\ge f^{-1}(y_{2})$,
applicando $f$ (crescente) otterremmo $y_{1}\ge y_{2}$.
\end{strategia}

\begin{dimostrazione}
Supponiamo, per fissare le idee, che $f$ sia strettamente crescente (il caso
decrescente è identico, con le disuguaglianze invertite). Siano $x_{1},x_{2}\in A$ con
$x_{1}\neq x_{2}$. Per la dicotomia \ref{ass:dicotomia} è $x_{1}<x_{2}$ oppure
$x_{2}<x_{1}$; nel primo caso $f(x_{1})<f(x_{2})$, nel secondo $f(x_{2})<f(x_{1})$. In
entrambi i casi $f(x_{1})\neq f(x_{2})$: $f$ è iniettiva. Essendo anche suriettiva, è
biunivoca e quindi invertibile.

Proviamo che $f^{-1}$ è strettamente crescente. Siano $y_{1},y_{2}\in B$ con
$y_{1}<y_{2}$ e poniamo $x_{1}=f^{-1}(y_{1})$, $x_{2}=f^{-1}(y_{2})$, così che
$f(x_{1})=y_{1}$ e $f(x_{2})=y_{2}$. Se fosse $x_{1}\ge x_{2}$, allora per la stretta
crescenza di $f$ (e banalmente, nel caso $x_{1}=x_{2}$) sarebbe $f(x_{1})\ge
f(x_{2})$, cioè $y_{1}\ge y_{2}$, contro l'ipotesi. Dunque $x_{1}<x_{2}$, cioè
$f^{-1}(y_{1})<f^{-1}(y_{2})$.
\end{dimostrazione}

\begin{attenzione}[Il viceversa è falso]
Una funzione invertibile non è necessariamente monotòna. La funzione
$f\colon[0,2]\to[0,2]$ definita da $f(x)=x$ per $x\in[0,1)$ e $f(x)=3-x$ per
$x\in[1,2]$ è biunivoca (il lettore ne disegni il grafico) ma non è monotòna. Per le
funzioni \emph{continue su un intervallo}, invece, invertibilità e stretta monotonia sono
equivalenti: nel capitolo sulle funzioni continue dimostreremo che una funzione continua
e strettamente monotòna su un intervallo è invertibile (\cref{teo:invertibilita}); anche
il viceversa è vero, come conseguenza del teorema dei valori intermedi.
\end{attenzione}

\subsection{Immagini, immagini inverse, funzioni composte}

Concludiamo il paragrafo con alcune utili definizioni e notazioni.

\begin{definizione}[Immagine e immagine inversa di un insieme]\index{immagine!di un insieme}\index{codominio}\index{immagine inversa}
Sia $f\colon A\to B$ una funzione. Se $X$ è un sottoinsieme di $A$, l'\emph{immagine}
di $X$ mediante $f$, indicata con $f(X)$, è il sottoinsieme di $B$ definito da
\begin{equation}\label{eq:immagine}
  f(X)=\{y\in B : \esiste x\in X : y=f(x)\}
\end{equation}
(il simbolo $\esiste$ si legge \emph{esiste}). L'insieme $f(A)$, cioè l'immagine di $A$
mediante $f$, si chiama \emph{codominio} o \emph{immagine} di $f$. Evidentemente la
funzione $f$ è suriettiva da $A$ verso $B$ se e solo se $f(A)=B$. Se $Y$ è un
sottoinsieme di $B$, l'\emph{immagine inversa} (o \emph{controimmagine}) di $Y$ mediante
$f$, indicata con $f^{-1}(Y)$, è il sottoinsieme di $A$ definito da
\begin{equation}\label{eq:controimmagine}
  f^{-1}(Y)=\{x\in A : f(x)\in Y\}.
\end{equation}
\end{definizione}

\begin{osservazione}
La notazione $f^{-1}(Y)$ ha senso \emph{anche se $f$ non è invertibile}: $f^{-1}(Y)$ è
un insieme, non il valore di una funzione. Ad esempio, per $f(x)=x^{2}$ si ha
$f^{-1}(\{4\})=\{-2,2\}$ e $f^{-1}(\{-1\})=\emptyset$. Allo stesso modo, nel linguaggio
degli ingegneri, l'immagine inversa di un intervallo di valori «accettabili» in uscita è
l'insieme degli ingressi che producono quelle uscite, indipendentemente dal fatto che la
legge ingresso--uscita sia invertibile.
\end{osservazione}

\begin{definizione}[Funzione composta]\index{funzione!composta}
Siano $X,Y,Z$ tre insiemi e siano $g\colon X\to Y$ e $f\colon Y\to Z$ due funzioni, tali
che l'insieme $Y$ contenente i valori della prima coincida con il dominio della seconda.
Allora si può considerare la \emph{funzione composta} $h\colon X\to Z$, definita da
$h(x)=f(g(x))$ per $x\in X$. In tal caso si usa la notazione $h=f\circ g$; in altre
parole, si pone
\[
  (f\circ g)(x)=f(g(x)) \qquad\perogni x\in X.
\]
\end{definizione}

\begin{marginfigure}
\centering
\begin{tikzpicture}[scale=0.8, every node/.style={font=\scriptsize}]
  \node[draw, blu, rounded corners=2pt, minimum width=1cm] (X) at (0,0) {$X$};
  \node[draw, blu, rounded corners=2pt, minimum width=1cm] (Y) at (2,0) {$Y$};
  \node[draw, blu, rounded corners=2pt, minimum width=1cm] (Z) at (4,0) {$Z$};
  \draw[-Latex] (X) -- node[above] {$g$} (Y);
  \draw[-Latex] (Y) -- node[above] {$f$} (Z);
  \draw[-Latex, rame] (X) to[bend right=40] node[below] {$f\circ g$} (Z);
\end{tikzpicture}
\caption{La composizione $f\circ g$: prima si applica $g$, poi $f$. Si legge «$f$
composto $g$» e si calcola da destra verso sinistra.}
\end{marginfigure}

\begin{esempio}
Siano $g(x)=2x+1$ e $f(y)=\sqrt{y}$. Allora $(f\circ g)(x)=\sqrt{2x+1}$, definita per
$x\ge-\tfrac12$, mentre $(g\circ f)(x)=2\sqrt{x}+1$, definita per $x\ge0$: la
composizione \emph{non è commutativa}. Con la notazione della composizione, le
relazioni \eqref{eq:inversa} che caratterizzano la funzione inversa si scrivono in modo
compatto come $f^{-1}\circ f=\mathrm{id}_{A}$ e $f\circ f^{-1}=\mathrm{id}_{B}$, dove
$\mathrm{id}$ indica la funzione \emph{identità} $x\mapsto x$.
\end{esempio}

\begin{python}[Comporre e invertire funzioni con \texttt{sympy}]
# sympy manipola le espressioni simbolicamente, come faremmo a mano.
import sympy as sp
x, y = sp.symbols("x y", real=True)

g = 2*x + 1
f = sp.sqrt(x)
print("f(g(x)) =", f.subs(x, g))         # composizione f o g
print("g(f(x)) =", g.subs(x, f))         # composizione g o f

# Inversa di g: risolviamo y = g(x) rispetto a x
inversa = sp.solve(sp.Eq(y, g), x)[0]
print("g^-1(y) =", inversa)
# Verifica delle relazioni g^-1(g(x)) = x e g(g^-1(y)) = y
print(sp.simplify(inversa.subs(y, g)),
      sp.simplify(g.subs(x, inversa)))
\end{python}
\begin{risultato}
f(g(x)) = sqrt(2*x + 1)
g(f(x)) = 2*sqrt(x) + 1
g^-1(y) = y/2 - 1/2
x y
\end{risultato}

\begin{esercizio}
Sia $f\colon\R\to\R$, $f(x)=x^{3}$. Dimostrare, usando solo la definizione, che $f$ è
strettamente crescente, e dedurne che è iniettiva. (La suriettività, cioè l'esistenza
della radice cubica di ogni reale, richiede l'estremo superiore e sarà provata nel
\cref{teo:radice-n}.)
\begin{soluzione}
Siano $x_{1}<x_{2}$. Allora $x_{2}-x_{1}>0$ e
$x_{2}^{3}-x_{1}^{3}=(x_{2}-x_{1})(x_{2}^{2}+x_{1}x_{2}+x_{1}^{2})$. Il secondo fattore
è $\bigl(x_{2}+\tfrac12x_{1}\bigr)^{2}+\tfrac34x_{1}^{2}\ge0$, e si annulla solo se
$x_{1}=x_{2}=0$, escluso. Dunque il prodotto è $>0$, cioè $x_{1}^{3}<x_{2}^{3}$.
L'iniettività segue dalla stretta crescenza, come nella prima parte della dimostrazione
del \cref{teo:monotona-invertibile}: se $x_{1}\neq x_{2}$, uno dei due è minore
dell'altro, e le immagini sono diverse. \emph{Scelta logica:} per confrontare
$x_{1}^{3}$ e $x_{2}^{3}$ si studia il segno della differenza
(\cref{prop:ordine-differenza}); la fattorizzazione di $x_{2}^{3}-x_{1}^{3}$ riduce il
problema al segno di un trinomio, che il completamento del quadrato rende evidente.
\end{soluzione}
\end{esercizio}

%% ----------------------------------------------------------------------------
%%  §1.7  Funzioni lineari. Funzione valore assoluto
%% ----------------------------------------------------------------------------
\section{Funzioni lineari. Funzione valore assoluto}
\label{sec:lineari}

\begin{definizione}[Funzione lineare]\index{funzione!lineare}\index{coefficiente angolare}
Si chiama \emph{funzione lineare} (o \emph{funzione affine}) una funzione del tipo
\begin{equation}\label{eq:lineare}
  y=mx+q,
\end{equation}
dove $m,q$ sono numeri reali fissati. Il grafico di una tale funzione è una retta, di cui
il parametro $m$ è detto \emph{coefficiente angolare} e $q$ \emph{ordinata all'origine}
(è il valore $f(0)$).
\end{definizione}

\begin{proposizione}\label{prop:lineare-monotona}
Ogni funzione lineare $f(x)=mx+q$ è monotòna su $\R$; più precisamente è strettamente
crescente se $m>0$, strettamente decrescente se $m<0$ e costante se $m=0$. Se $m\neq0$
essa è invertibile da $\R$ in $\R$, con inversa
\begin{equation}\label{eq:inversa-lineare}
  f^{-1}(y)=\frac{y-q}{m}.
\end{equation}
\end{proposizione}

\begin{schema}
\Ipotesi{$f(x)=mx+q$, con $m,q\in\R$.}
\Tesi{$m>0\implica f$ strettamente crescente; $m<0\implica f$ strettamente decrescente;
$m=0\implica f$ costante; se $m\neq0$, $f$ è invertibile e $f^{-1}(y)=(y-q)/m$.}
\end{schema}

\begin{strategia}
È un'applicazione diretta della \cref{prop:ordine-prodotto}: moltiplicare la
disuguaglianza $x_{1}<x_{2}$ per $m$ la conserva se $m>0$ e la rovescia se $m<0$;
sommare $q$ non cambia nulla. Per l'inversa basta risolvere l'equazione $y=mx+q$
rispetto a $x$, cosa possibile se e solo se $m\neq0$.
\end{strategia}

\begin{dimostrazione}
Consideriamo $x_{1}<x_{2}$; allora
\[
  f(x_{1})=mx_{1}+q, \qquad f(x_{2})=mx_{2}+q .
\]
Se il coefficiente angolare $m$ è positivo, per la \cref{prop:ordine-prodotto}(i) da
$x_{1}<x_{2}$ segue $mx_{1}<mx_{2}$ (la disuguaglianza resta stretta perché $m\neq0$ e
la \cref{prop:semplificazione-prodotto} esclude l'uguaglianza) e quindi, sommando $q$
(assioma \ref{ass:ordine-somma}), $f(x_{1})<f(x_{2})$: $f$ è strettamente crescente. Se
$m$ è negativo, per la \cref{prop:ordine-prodotto}(ii) da $x_{1}<x_{2}$ segue
$mx_{1}>mx_{2}$ e quindi $f(x_{1})>f(x_{2})$: $f$ è strettamente decrescente. Infine, se
$m=0$, risulta $f(x)=q$ per ogni $x$: la funzione è costante, dunque
contemporaneamente crescente e decrescente, e il suo grafico è una retta parallela
all'asse $x$.

Se $m\neq0$, $f$ è strettamente monotòna; è anche suriettiva su $\R$, perché per ogni
$y\in\R$ l'equazione $mx+q=y$ ha la soluzione $x=(y-q)/m$. Per il
\cref{teo:monotona-invertibile} $f$ è invertibile, e la \eqref{eq:inversa-lineare} è
di verifica immediata: $f\bigl((y-q)/m\bigr)=m\cdot(y-q)/m+q=y$. Al contrario, se $m=0$
la funzione costante $f(x)=q$ non stabilisce una corrispondenza biunivoca tra $\R$ e
l'insieme $\{q\}$, e non è invertibile.
\end{dimostrazione}

Dato che per due punti distinti del piano passa una e una sola retta, per disegnare il
grafico di una funzione lineare è sufficiente calcolare l'ordinata $y$ in
\eqref{eq:lineare} in corrispondenza a due valori distinti della variabile $x$. Ad
esempio, nel caso della funzione lineare $y=3x-1$, a $x=0$ corrisponde $y=-1$ e a $x=1$
corrisponde $y=2$; il grafico è la retta passante per i punti $(0,-1)$ e $(1,2)$
(\cref{fig:rette}).

\begin{figure}[htb]
\centering
\begin{tikzpicture}
\begin{axis}[width=5.8cm, height=5.2cm, xmin=-1.5, xmax=2.5, ymin=-2.5, ymax=4,
  xtick={1}, ytick={-1,2}, title={\footnotesize $y=3x-1$}, title style={at={(0.3,0.9)}}]
  \addplot[curva, domain=-0.5:1.6] {3*x-1};
  \addplot[tratteggio] coordinates {(1,0) (1,2) (0,2)};
  \addplot[only marks, mark=*, mark size=1.4pt, rame] coordinates {(0,-1) (1,2)};
\end{axis}
\end{tikzpicture}\hfill
\begin{tikzpicture}
\begin{axis}[width=5.8cm, height=5.2cm, xmin=-1, xmax=3.5, ymin=-1.5, ymax=4,
  xtick={1,2,3}, ytick={1,2,3}, title={\footnotesize $y=3-2x$}, title style={at={(0.75,0.9)}}]
  \addplot[curva, domain=-0.5:2.2] {3-2*x};
  \addplot[only marks, mark=*, mark size=1.4pt, rame] coordinates {(0,3) (1.5,0)};
\end{axis}
\end{tikzpicture}
\caption{Grafici di due funzioni lineari: $y=3x-1$ (coefficiente angolare $m=3>0$,
strettamente crescente) e $y=3-2x$ ($m=-2<0$, strettamente decrescente). In rame i
due punti usati per tracciare ciascuna retta.}
\label{fig:rette}
\end{figure}

\begin{idea}[Il coefficiente angolare come «velocità»]
Il coefficiente angolare $m$ misura di quanto cresce $y$ quando $x$ cresce di $1$:
$f(x+1)-f(x)=m$. Se $x$ è il tempo e $y$ la posizione di un carrello lungo una guida,
$m$ è la velocità: positiva se il carrello avanza, negativa se arretra, nulla se è fermo.
La \cref{prop:lineare-monotona} dice allora qualcosa di ovvio dal punto di vista
fisico: un carrello con velocità costante positiva non torna mai indietro. Il calcolo
differenziale, nei capitoli successivi, estenderà questa idea a funzioni qualsiasi,
definendo una «velocità istantanea» --- la derivata --- che varia da punto a punto.
\end{idea}

\subsection{Il valore assoluto}

\begin{definizione}[Valore assoluto]\index{valore assoluto}\index{modulo}
Il \emph{valore assoluto} (o \emph{modulo}) di $x\in\R$, indicato con il simbolo
$\abs{x}$, è definito da
\begin{equation}\label{eq:modulo}
  \abs{x}=\begin{cases} x & \text{se } x\ge0,\\ -x & \text{se } x<0. \end{cases}
\end{equation}
\end{definizione}

\begin{marginfigure}
\centering
\begin{tikzpicture}
\begin{axis}[width=5cm, height=4cm, xmin=-2.5, xmax=2.5, ymin=-0.5, ymax=2.5,
  xtick={-1,1}, ytick={1}, title={\footnotesize $y=\abs{x}$}, title style={at={(0.5,0.9)}}]
  \addplot[curva, domain=-2.3:0] {-x};
  \addplot[curva, domain=0:2.3] {x};
\end{axis}
\end{tikzpicture}
\caption{Il grafico della funzione valore assoluto è l'unione della semiretta $y=x$,
per $x\ge0$, e della semiretta $y=-x$, per $x<0$.}
\label{fig:modulo}
\end{marginfigure}

Il grafico della funzione valore assoluto $f(x)=\abs{x}$ è composto da due semirette
uscenti dall'origine, di equazione rispettivamente $y=x$ e $y=-x$, come nella
\cref{fig:modulo}. Geometricamente, $\abs{x}$ è la \emph{distanza} del punto $x$
dall'origine sulla retta reale, e $\abs{x-y}$ è la distanza tra i punti $x$ e $y$: è
questa interpretazione che rende il valore assoluto onnipresente in Analisi, dove
«vicino» significa sempre «a distanza piccola».

\begin{proposizione}[Proprietà elementari del valore assoluto]\label{prop:modulo-proprieta}
Per ogni $x,x_{1},x_{2}\in\R$:
\begin{align}
  &\abs{x}\ge0; \label{eq:mod1}\\
  &\abs{x}=0 \ \sse\ x=0; \label{eq:mod2}\\
  &\abs{-x}=\abs{x}; \label{eq:mod3}\\
  &\abs{x_{1}\cdot x_{2}}=\abs{x_{1}}\cdot\abs{x_{2}}; \label{eq:mod4}\\
  &\abs{x_{1}/x_{2}}=\abs{x_{1}}/\abs{x_{2}} \quad (x_{2}\neq0). \label{eq:mod5}
\end{align}
\end{proposizione}

\begin{schema}
\Ipotesi{$x,x_{1},x_{2}\in\R$; nella \eqref{eq:mod5}, $x_{2}\neq0$.}
\Tesi{Le relazioni \eqref{eq:mod1}--\eqref{eq:mod5}.}
\end{schema}

\begin{strategia}
Il valore assoluto è definito \emph{per casi}, e le sue proprietà si dimostrano quindi
\emph{distinguendo i casi}: $x\ge0$ e $x<0$ (e, per le proprietà con due variabili, le
quattro combinazioni di segno). È un metodo noioso ma infallibile; conviene farlo una
volta per tutte e poi non pensarci più.
\end{strategia}

\begin{dimostrazione}
Dimostriamo ad esempio la \eqref{eq:mod1}: se $x\ge0$ allora $\abs{x}=x\ge0$; se invece
$x<0$, allora $\abs{x}=-x$ e, per la \cref{prop:segno-opposto}, $-x>0$; in definitiva
$\abs{x}\ge0$ per ogni $x\in\R$. La \eqref{eq:mod2} è immediata dalla definizione:
$\abs{x}=0$ significa $x=0$ (caso $x\ge0$) oppure $-x=0$ (caso $x<0$, impossibile).
Proviamo la \eqref{eq:mod3}: se $x>0$ allora $\abs{x}=x$ e $\abs{-x}=-(-x)=x$, dato che
$-x<0$; se $x<0$ risulta $\abs{x}=-x$ e $\abs{-x}=-x$, essendo $-x>0$; infine, se $x=0$,
si ha $-0=0$ e pertanto $\abs{-0}=\abs{0}$. La \eqref{eq:mod4} è diretta conseguenza
della regola dei segni (\cref{prop:segni}): ad esempio, se $x_{1}<0$ e $x_{2}\ge0$, il
prodotto $x_{1}x_{2}$ è $\le0$, dunque
$\abs{x_{1}x_{2}}=-(x_{1}x_{2})=(-x_{1})x_{2}=\abs{x_{1}}\abs{x_{2}}$; gli altri casi
sono analoghi. Infine la \eqref{eq:mod5} segue dalla \eqref{eq:mod4}:
$\abs{x_{1}}=\abs{(x_{1}/x_{2})\,x_{2}}=\abs{x_{1}/x_{2}}\,\abs{x_{2}}$, e basta dividere
per $\abs{x_{2}}$, che è non nullo per la \eqref{eq:mod2}.
\end{dimostrazione}

Esaminiamo ora due proprietà la cui interpretazione geometrica è schematizzata nella
\cref{fig:modulo-r}: esse permettono di tradurre una disuguaglianza con il valore
assoluto in una \emph{doppia} disuguaglianza senza valore assoluto, ed è così che le
useremo sistematicamente nello studio dei limiti.

\begin{proposizione}\label{prop:modulo-r}
Per ogni numero reale $r\ge0$ valgono le equivalenze:
\begin{align}
  \abs{x}\le r &\ \sse\ -r\le x\le r; \label{eq:mod-le}\\
  \abs{x}<r    &\ \sse\ -r<x<r. \label{eq:mod-lt}
\end{align}
\end{proposizione}

\begin{schema}
\Ipotesi{$r\ge0$, $x\in\R$.}
\Tesi{$\abs{x}\le r$ se e solo se $x\le r$ \emph{e} $x\ge-r$ (e analogamente con le
disuguaglianze strette).}
\end{schema}

\begin{strategia}
Ancora una distinzione di casi: la relazione $\abs{x}\le r$ significa «$x\le r$ se
$x\ge0$» oppure «$-x\le r$ se $x<0$». Ciascun caso dà un pezzo dell'intervallo
$[-r,r]$: il primo la metà destra $[0,r]$, il secondo la metà sinistra $[-r,0)$.
Riunendo i due pezzi si ottiene tutto l'intervallo.
\end{strategia}

\begin{dimostrazione}
La verifica delle due equivalenze procede allo stesso modo; perciò ci limitiamo a
dimostrare la \eqref{eq:mod-le}. In base alla definizione \eqref{eq:modulo} del valore
assoluto, la relazione $\abs{x}\le r$ equivale a uno dei due sistemi
\begin{equation}\label{eq:mod-sistemi}
  \begin{cases} x\ge0\\ x\le r\end{cases} \qquad\text{oppure}\qquad
  \begin{cases} x<0\\ -x\le r.\end{cases}
\end{equation}
Il primo sistema si riscrive nella forma $0\le x\le r$, mentre il secondo equivale (per
la \cref{prop:ordine-prodotto}(ii), moltiplicando per $-1$) a $-r\le x<0$. Unendo i due
risultati si trova infine $-r\le x\le r$. Viceversa, se $-r\le x\le r$, allora
$x\le r$ e $-x\le r$; poiché $\abs{x}$ è uguale a $x$ oppure a $-x$, in ogni caso
$\abs{x}\le r$.
\end{dimostrazione}

\begin{figure}[htb]
\centering
\begin{tikzpicture}
\begin{axis}[width=8cm, height=5cm, xmin=-2.8, xmax=2.8, ymin=-0.4, ymax=2.6,
  xtick={-1.6,1.6}, xticklabels={$-r$,$r$}, ytick={1.6}, yticklabels={$r$}]
  \addplot[curva, domain=-2.5:0] {-x};
  \addplot[curva, domain=0:2.5] {x};
  \addplot[tratteggio] coordinates {(-1.6,0) (-1.6,1.6) (1.6,1.6) (1.6,0)};
  \addplot[rame, line width=2.2pt] coordinates {(-1.6,0) (1.6,0)};
  \addplot[verde, line width=2.2pt] coordinates {(0,0) (0,1.6)};
  \node[font=\scriptsize, rame, below] at (axis cs:0.85,-0.02) {$-r\le x\le r$};
  \node[font=\scriptsize, verde, right] at (axis cs:0.05,0.9) {$y\le r$};
  \node[font=\scriptsize, blu, anchor=east] at (axis cs:2.05,2.35) {$y=\abs{x}$};
\end{axis}
\end{tikzpicture}
\caption{Interpretazione geometrica della \cref{prop:modulo-r}: i punti del grafico di
$y=\abs{x}$ con ordinata $y\le r$ (segmento verde) sono esattamente quelli con ascissa
nell'intervallo $[-r,r]$ (segmento color rame).}
\label{fig:modulo-r}
\end{figure}

La seguente proprietà del valore assoluto, detta \emph{disuguaglianza triangolare}, è
di grande importanza. La spiegazione intuitiva è semplice: a secondo membro della
\eqref{eq:triangolare} compare la somma di due numeri positivi o nulli (i moduli dei
numeri $x_{1}$ e $x_{2}$), mentre a primo membro compare il modulo della \emph{somma
algebrica} di $x_{1}$ e $x_{2}$; se i segni di $x_{1}$ e $x_{2}$ sono discordi, il
primo membro è minore del secondo, mentre se i segni sono concordi i due membri sono
uguali.

\begin{teorema}[Disuguaglianza triangolare]\label{teo:triangolare}\index{disuguaglianza triangolare}
Per ogni coppia di numeri reali $x_{1},x_{2}$ vale la disuguaglianza
\begin{equation}\label{eq:triangolare}
  \abs{x_{1}+x_{2}}\le\abs{x_{1}}+\abs{x_{2}}.
\end{equation}
\end{teorema}

\begin{schema}
\Ipotesi{$x_{1},x_{2}\in\R$.}
\Tesi{$\abs{x_{1}+x_{2}}\le\abs{x_{1}}+\abs{x_{2}}$.}
\end{schema}

\begin{strategia}
Potremmo distinguere i quattro casi di segno, ma c'è una strada più elegante che
sfrutta la \cref{prop:modulo-r} \emph{due volte}: una volta «in uscita», per
trasformare le ovvie relazioni $\abs{x_{i}}\le\abs{x_{i}}$ nelle doppie disuguaglianze
$-\abs{x_{i}}\le x_{i}\le\abs{x_{i}}$; poi sommiamo membro a membro; infine usiamo la
\cref{prop:modulo-r} «in entrata», per ricompattare la doppia disuguaglianza ottenuta
in una disuguaglianza con il modulo. È un esempio tipico di come una
\emph{equivalenza} si possa usare nei due versi.
\end{strategia}

\begin{dimostrazione}
Per ogni numero reale $x$ la relazione $\abs{x}\le\abs{x}$ è ovvia (anzi, vale con il
segno di uguaglianza); indicando con $r$ il secondo membro di tale ovvia relazione, per
l'equivalenza \eqref{eq:mod-le} abbiamo anche $-\abs{x}\le x\le\abs{x}$; in particolare,
per $x=x_{1}$ e $x=x_{2}$:
\begin{equation}\label{eq:triang-passo1}
  -\abs{x_{1}}\le x_{1}\le\abs{x_{1}}, \qquad -\abs{x_{2}}\le x_{2}\le\abs{x_{2}}.
\end{equation}
Sommando membro a membro (è lecito: per \ref{ass:ordine-somma} e la proprietà
transitiva, da $a\le b$ e $c\le d$ segue $a+c\le b+c\le b+d$):
\begin{equation}\label{eq:triang-passo2}
  -\bigl(\abs{x_{1}}+\abs{x_{2}}\bigr)\le x_{1}+x_{2}\le\abs{x_{1}}+\abs{x_{2}}.
\end{equation}
La conclusione segue da una nuova applicazione dell'equivalenza \eqref{eq:mod-le}, con
$r=\abs{x_{1}}+\abs{x_{2}}$ e $x=x_{1}+x_{2}$.
\end{dimostrazione}

\begin{notastorica}[Perché «triangolare»?]
Il nome viene dalla geometria: in un triangolo, ogni lato è minore della somma degli
altri due (proposizione I.20 degli \emph{Elementi} di Euclide, che il filosofo Proclo
riferisce fosse derisa dagli epicurei perché «nota anche a un asino», che per
raggiungere il fieno va dritto e non lungo i due lati). Interpretando $x_{1}$ e $x_{2}$
come vettori nel piano, $\abs{x_{1}+x_{2}}\le\abs{x_{1}}+\abs{x_{2}}$ è esattamente
quella proposizione; sulla retta i «triangoli» sono degeneri, ma il nome è rimasto. La
disuguaglianza triangolare è la proprietà fondamentale di ogni nozione di
\emph{distanza}: è uno degli assiomi con cui Maurice Fréchet, nel 1906, introdusse gli
spazi che oggi chiamiamo \emph{metrici}.
\end{notastorica}

\begin{corollario}\label{cor:triangolare-inversa}
Per ogni $x_{1},x_{2}\in\R$ risulta
\begin{equation}\label{eq:triangolare-inversa}
  \bigl|\abs{x_{1}}-\abs{x_{2}}\bigr|\le\abs{x_{1}-x_{2}}.
\end{equation}
\end{corollario}

\begin{schema}
\Ipotesi{$x_{1},x_{2}\in\R$.}
\Tesi{$\bigl|\abs{x_{1}}-\abs{x_{2}}\bigr|\le\abs{x_{1}-x_{2}}$.}
\end{schema}

\begin{strategia}
Vogliamo una stima di $\abs{x_{1}}-\abs{x_{2}}$: il trucco è scrivere $x_{1}$ come
$(x_{1}-x_{2})+x_{2}$ e applicare la disuguaglianza triangolare a questa somma. Si
ottiene $\abs{x_{1}}-\abs{x_{2}}\le\abs{x_{1}-x_{2}}$; scambiando i ruoli si ottiene
la stima con il segno opposto, e la \cref{prop:modulo-r} ricompone le due.
\end{strategia}

\begin{dimostrazione}
Per il \cref{teo:triangolare}, $\abs{x_{1}}=\abs{(x_{1}-x_{2})+x_{2}}\le\abs{x_{1}-x_{2}}+\abs{x_{2}}$,
da cui $\abs{x_{1}}-\abs{x_{2}}\le\abs{x_{1}-x_{2}}$. Scambiando $x_{1}$ con $x_{2}$ e
usando la \eqref{eq:mod3}: $\abs{x_{2}}-\abs{x_{1}}\le\abs{x_{2}-x_{1}}=\abs{x_{1}-x_{2}}$,
cioè $-\abs{x_{1}-x_{2}}\le\abs{x_{1}}-\abs{x_{2}}$. Le due stime, per la
\eqref{eq:mod-le} con $r=\abs{x_{1}-x_{2}}$, danno la tesi.
\end{dimostrazione}

\begin{python}[Verifica sperimentale della disuguaglianza triangolare]
# Un esperimento numerico non dimostra nulla, ma aiuta a "vedere"
# quando la disuguaglianza è stretta e quando è un'uguaglianza.
import random
random.seed(1)
for _ in range(6):
    x1, x2 = random.uniform(-5, 5), random.uniform(-5, 5)
    sinistra, destra = abs(x1 + x2), abs(x1) + abs(x2)
    segni = "concordi" if x1 * x2 >= 0 else "discordi"
    rel = "=" if abs(sinistra - destra) < 1e-12 else "<"
    print(f"x1={x1:+.2f} x2={x2:+.2f} ({segni}): "
          f"|x1+x2|={sinistra:.2f} {rel} |x1|+|x2|={destra:.2f}")
\end{python}
\begin{risultato}
x1=-3.66 x2=+3.47 (discordi): |x1+x2|=0.18 < |x1|+|x2|=7.13
x1=+2.64 x2=-2.45 (discordi): |x1+x2|=0.19 < |x1|+|x2|=5.09
x1=-0.05 x2=-0.51 (concordi): |x1+x2|=0.55 = |x1|+|x2|=0.55
x1=+1.52 x2=+2.89 (concordi): |x1+x2|=4.40 = |x1|+|x2|=4.40
x1=-4.06 x2=-4.72 (concordi): |x1+x2|=8.78 = |x1|+|x2|=8.78
x1=+3.36 x2=-0.67 (discordi): |x1+x2|=2.69 < |x1|+|x2|=4.03
\end{risultato}

\begin{esercizio}
Risolvere la disequazione $\abs{2x-3}<5$.
\begin{soluzione}
Per la \eqref{eq:mod-lt} con $r=5$: $-5<2x-3<5$, cioè $-2<2x<8$, cioè $-1<x<4$.
\emph{Scelta logica:} la \cref{prop:modulo-r} trasforma la disequazione con il modulo
in una doppia disequazione senza modulo, ed evita la discussione dei due casi
$2x-3\ge0$ e $2x-3<0$.
\end{soluzione}
\end{esercizio}

\begin{esercizio}
Dimostrare che per ogni $x,y\in\R$ si ha $\max\{x,y\}=\dfrac{x+y+\abs{x-y}}{2}$ e
$\min\{x,y\}=\dfrac{x+y-\abs{x-y}}{2}$.
\begin{soluzione}
I secondi membri non cambiano scambiando $x$ con $y$, perché $\abs{x-y}=\abs{y-x}$
(\eqref{eq:mod3}); lo stesso vale per $\max\{x,y\}$ e $\min\{x,y\}$. Possiamo quindi
supporre $x\ge y$. Allora $x-y\ge0$, $\abs{x-y}=x-y$, e
\[
\begin{aligned}
  \frac{x+y+\abs{x-y}}{2}&=\frac{x+y+(x-y)}{2}=x=\max\{x,y\},\\
  \frac{x+y-\abs{x-y}}{2}&=\frac{x+y-(x-y)}{2}=y=\min\{x,y\}.
\end{aligned}
\]
\emph{Scelta logica:} il modulo è definito per casi, e il caso da trattare dipende dal
segno del suo argomento $x-y$; la simmetria delle formule rispetto allo scambio di $x$
con $y$ permette di trattarne uno solo, «senza perdita di generalità».
\end{soluzione}
\end{esercizio}

%% ----------------------------------------------------------------------------
%%  §1.8  Le funzioni potenza, esponenziale, logaritmo
%% ----------------------------------------------------------------------------
\section{Le funzioni potenza, esponenziale, logaritmo}
\label{sec:potenze}

\subsection{Potenze con esponente naturale e radici}

Consideriamo la \emph{funzione potenza con esponente $n\in\N$}:\index{funzione!potenza}
\begin{equation}\label{eq:potenza-n}
  f(x)=x^{n},
\end{equation}
che è definita, per ogni $x\in\R$, moltiplicando il numero $x$ per se stesso $n$ volte
(più precisamente: $x^{1}=x$ e $x^{n+1}=x^{n}\cdot x$; si pone inoltre $x^{0}=1$).

\begin{proposizione}\label{prop:potenza-crescente}
La funzione $f(x)=x^{n}$ è strettamente crescente per $x\ge0$, cioè
\begin{equation}\label{eq:potenza-crescente}
  0\le x_{1}<x_{2} \ \implica\ x_{1}^{n}<x_{2}^{n} \qquad\perogni n\in\N.
\end{equation}
\end{proposizione}

\begin{schema}
\Ipotesi{$n\in\N$; \ $0\le x_{1}<x_{2}$.}
\Tesi{$x_{1}^{n}<x_{2}^{n}$.}
\end{schema}

\begin{strategia}
Per $n=2$ il ragionamento è questo: da $x_{1}<x_{2}$, moltiplicando per $x_{1}\ge0$ si
ottiene $x_{1}^{2}\le x_{1}x_{2}$ (\cref{prop:ordine-prodotto}; uguaglianza possibile
solo se $x_{1}=0$); moltiplicando invece per $x_{2}>0$ si ottiene $x_{1}x_{2}<x_{2}^{2}$.
Mettendo insieme le due, $x_{1}^{2}<x_{2}^{2}$. Per $n=3$ si riparte da
$x_{1}^{2}<x_{2}^{2}$ e si ripete lo stesso gioco. Il «e così via» che ne segue è
proprio il contenuto del \emph{principio di induzione}, che enunceremo e useremo
rigorosamente nel \cref{sec:induzione}: qui ci accontentiamo dei primi passi.
\end{strategia}

\begin{dimostrazione}
Per $n=1$ la tesi coincide con l'ipotesi. Se $n=2$ e $0\le x_{1}<x_{2}$, allora
moltiplicando prima per $x_{1}\ge0$, poi per $x_{2}>0$, si ottiene
$x_{1}^{2}\le x_{1}x_{2}$ e $x_{1}x_{2}<x_{2}^{2}$ (la seconda disuguaglianza è stretta
perché $x_{2}\neq0$ e $x_{1}\neq x_{2}$, per la \cref{prop:semplificazione-prodotto});
cioè $x_{1}^{2}<x_{2}^{2}$.
Se $n=3$, partendo da $x_{1}^{2}<x_{2}^{2}$ si ottiene
$x_{1}^{3}=x_{1}x_{1}^{2}\le x_{1}x_{2}^{2}<x_{2}x_{2}^{2}=x_{2}^{3}$, e così via per
$n>3$ (si veda la \cref{prop:induzione-potenza} per la dimostrazione completa per
induzione).
\end{dimostrazione}

Per mezzo dell'estremo superiore dimostreremo nel prossimo capitolo
(\cref{teo:radice-n}) che a ogni $y\ge0$ corrisponde almeno un numero reale $x\ge0$ per cui
$f(x)=x^{n}=y$: cioè che $f$ è suriettiva da $[0,+\infty)$ su $[0,+\infty)$. La
condizione di stretta monotonia \eqref{eq:potenza-crescente} implica allora, per il
criterio di invertibilità (\cref{teo:monotona-invertibile}), che la funzione è
invertibile. Perciò è definita la funzione inversa di $f(x)=x^{n}$ ($x\ge0$), che si
chiama funzione \emph{radice $n$-esima}\index{radice $n$-esima} e si indica con
\begin{equation}\label{eq:radice}
  f^{-1}(x)=\sqrt[n]{x}=x^{1/n} \qquad (x\ge0).
\end{equation}
I grafici delle funzioni \eqref{eq:potenza-n}, per $x\ge0$, e \eqref{eq:radice} sono
riportati nella \cref{fig:potenze-radici}: in accordo con la \cref{fig:inversa}, sono
simmetrici rispetto alla bisettrice.

\begin{figure}[htb]
\centering
\begin{tikzpicture}
\begin{axis}[width=6cm, height=5.6cm, xmin=-0.1, xmax=1.8, ymin=-0.1, ymax=1.8,
  xtick={1}, ytick={1}, title={\footnotesize $y=x^{n}$ \ ($x\ge0$)}, title style={at={(0.4,0.92)}},
  legend entries={$n=1$,$n=2$,$n=3$}, legend pos=south east, legend cell align=left,
  legend style={font=\scriptsize, draw=none, fill=none}]
  \addplot[curva, domain=0:1.75] {x};
  \addplot[curva2, domain=0:1.32] {x^2};
  \addplot[curva3, domain=0:1.2] {x^3};
\end{axis}
\end{tikzpicture}\hfill
\begin{tikzpicture}
\begin{axis}[width=6cm, height=5.6cm, xmin=-0.1, xmax=1.8, ymin=-0.1, ymax=1.8,
  xtick={1}, ytick={1}, title={\footnotesize $y=\sqrt[n]{x}$}, title style={at={(0.4,0.92)}},
  legend entries={$n=1$,$n=2$,$n=3$}, legend pos=south east, legend cell align=left,
  legend style={font=\scriptsize, draw=none, fill=none}]
  \addplot[curva, domain=0:1.75] {x};
  \addplot[curva2, domain=0:1.75, samples=100] {sqrt(x)};
  \addplot[curva3, domain=0:1.75, samples=100] {x^(1/3)};
\end{axis}
\end{tikzpicture}
\caption{Le funzioni potenza $y=x^{n}$ (a sinistra) e le loro inverse, le radici
$y=\sqrt[n]{x}$ (a destra), per $n=1,2,3$. Tutte passano per $(0,0)$ e $(1,1)$; per
$0<x<1$ le potenze decrescono al crescere di $n$, per $x>1$ crescono.}
\label{fig:potenze-radici}
\end{figure}

\subsection{Potenze con esponente razionale e reale}

Per mezzo delle funzioni \eqref{eq:potenza-n} e \eqref{eq:radice} si può definire
l'elevazione a esponente \emph{razionale} ($m,n\in\N$, $x\in\R$, $x>0$):
\begin{equation}\label{eq:potenza-razionale}
  x^{m/n}=\sqrt[n]{x^{m}}, \qquad x^{-m/n}=\frac{1}{\sqrt[n]{x^{m}}}, \qquad x^{0}=1.
\end{equation}
Si verifica che la definizione è ben posta, cioè che il valore non dipende dalla
rappresentazione della frazione: $x^{2/4}=x^{1/2}$ (è un buon esercizio, che usa
l'iniettività della potenza $n$-esima). A questo punto è stato definito il significato
di $a^{b}$, con $a$ numero reale positivo e $b$ numero razionale. Utilizzando
l'assioma di completezza è possibile estendere la definizione di $a^{b}$ anche al caso
in cui l'esponente $b$ sia un numero reale non razionale, come vedremo nel prossimo
capitolo (\cref{sec:esponenziale-reale}): l'idea è di definire, per $a>1$, $a^{\sqrt2}$
come l'unico numero reale che separa
$\{a^{q}: q\in\Q,\ q<\sqrt2\}$ da $\{a^{q}: q\in\Q,\ q>\sqrt2\}$ --- ancora una volta
la completezza «riempie il buco». Un'altra definizione equivalente, tramite la funzione
esponenziale, sarà proposta nel capitolo sulle serie.

Elenchiamo le proprietà fondamentali delle potenze, valide per ogni coppia di basi
$a,b>0$ e ogni coppia di esponenti $s,t\in\R$ (per esponenti interi sono conseguenze
immediate della definizione; per esponenti razionali si deducono da quelle con esponente
intero; per esponenti reali saranno dimostrate nei prossimi capitoli,
\cref{teo:esponenziale} e \cref{teo:potenze-reali}):
\begin{gather}
  a^{s}\cdot a^{t}=a^{s+t}, \qquad (a^{s})^{t}=a^{s\cdot t}, \qquad (a\cdot b)^{s}=a^{s}\cdot b^{s}; \label{eq:pot1}\\
  a^{s}>0; \label{eq:pot2}\\
  a<b,\ s>0 \ \implica\ a^{s}<b^{s}; \label{eq:pot3}\\
  a<b,\ s<0 \ \implica\ a^{s}>b^{s}; \label{eq:pot4}\\
  a>1,\ s<t \ \implica\ a^{s}<a^{t}; \label{eq:pot5}\\
  a<1,\ s<t \ \implica\ a^{s}>a^{t}. \label{eq:pot6}
\end{gather}

\subsection{Funzione potenza e funzione esponenziale}

Dall'espressione $a^{b}$ derivano due diversi tipi di funzione, a seconda che si faccia
variare la base $a$ o l'esponente $b$. Nel primo caso consideriamo la \emph{funzione
potenza} $f(x)=x^{b}$, con $b\in\R$ fissato. Nel secondo caso abbiamo la \emph{funzione
esponenziale}\index{funzione!esponenziale} $f(x)=a^{x}$, con $a$ numero reale positivo
fissato.

Casi particolari della funzione potenza $f(x)=x^{b}$ sono quelli con $b=n\in\N$, oppure
$b=1/n$, già esaminati in precedenza. La funzione potenza, per $x>0$, in base alla
\eqref{eq:pot2}, è positiva. Inoltre $x^{b}$ \emph{è una funzione strettamente crescente
se $b>0$ e strettamente decrescente se $b<0$}, in base alle \eqref{eq:pot3},
\eqref{eq:pot4}. Esempi di grafici nei vari casi sono riportati nella
\cref{fig:potenze-reali}.

\begin{figure}[htb]
\centering
\begin{tikzpicture}
\begin{axis}[width=4.3cm, height=4.6cm, xmin=-0.1, xmax=2.2, ymin=-0.1, ymax=3.2,
  xtick={1}, ytick={1}, title={\footnotesize $x^{b}$, $b>1$}, title style={at={(0.4,0.92)}}]
  \addplot[curva, domain=0:1.75, samples=100] {x^2};
\end{axis}
\end{tikzpicture}\hfill
\begin{tikzpicture}
\begin{axis}[width=4.3cm, height=4.6cm, xmin=-0.1, xmax=2.2, ymin=-0.1, ymax=3.2,
  xtick={1}, ytick={1}, title={\footnotesize $x^{b}$, $0<b<1$}, title style={at={(0.45,0.92)}}]
  \addplot[curva, domain=0:2.1, samples=100] {x^0.5};
\end{axis}
\end{tikzpicture}\hfill
\begin{tikzpicture}
\begin{axis}[width=4.3cm, height=4.6cm, xmin=-0.1, xmax=2.2, ymin=-0.1, ymax=3.2,
  xtick={1}, ytick={1}, title={\footnotesize $x^{b}$, $b<0$}, title style={at={(0.5,0.92)}}]
  \addplot[curva, domain=0.56:2.1, samples=100] {x^(-2)};
  \addplot[curva2, domain=0.1:2.1, samples=100] {x^(-0.5)};
  \node[blu, font=\scriptsize, anchor=west] at (axis cs:0.75,2.2) {$b=-2$};
  \node[rame, font=\scriptsize, anchor=south] at (axis cs:1.6,0.95) {$b=-\tfrac12$};
\end{axis}
\end{tikzpicture}
\caption{La funzione potenza $y=x^{b}$ per $x>0$: crescente se $b>0$ (convessa se
$b>1$, concava se $0<b<1$), decrescente se $b<0$.}
\label{fig:potenze-reali}
\end{figure}

La \emph{funzione esponenziale} $f(x)=a^{x}$, con $a$ numero reale positivo, è una
funzione positiva, \emph{strettamente crescente se $a>1$ e strettamente decrescente se
$a<1$}, in base alle \eqref{eq:pot2}, \eqref{eq:pot5}, \eqref{eq:pot6}. Esempi di
grafici sono riportati nella \cref{fig:esponenziali}.

\begin{figure}[htb]
\centering
\begin{tikzpicture}
\begin{axis}[width=6cm, height=5cm, xmin=-2.5, xmax=2.2, ymin=-0.3, ymax=4.5,
  xtick={1}, ytick={1}, title={\footnotesize $a^{x}$, $a>1$}, title style={at={(0.3,0.9)}}]
  \addplot[curva, domain=-2.4:2.1, samples=100] {2^x};
  \addplot[tratteggio] coordinates {(1,0) (1,2) (0,2)};
\end{axis}
\end{tikzpicture}\hfill
\begin{tikzpicture}
\begin{axis}[width=6cm, height=5cm, xmin=-2.2, xmax=2.5, ymin=-0.3, ymax=4.5,
  xtick={1}, ytick={1}, title={\footnotesize $a^{x}$, $0<a<1$}, title style={at={(0.7,0.9)}}]
  \addplot[curva, domain=-2.1:2.4, samples=100] {0.5^x};
\end{axis}
\end{tikzpicture}
\caption{La funzione esponenziale $y=a^{x}$: $a=2$ (a sinistra) e $a=\tfrac12$ (a
destra). Entrambe passano per $(0,1)$, sono sempre positive, e l'una è la riflessa
dell'altra rispetto all'asse $y$, perché $(1/2)^{x}=2^{-x}$.}
\label{fig:esponenziali}
\end{figure}

Un caso notevolmente importante, e ciò diventerà chiaro nello studio delle derivate, si
ha quando la base è uguale al \emph{numero di Nepero}\index{numero di Nepero}
$e=2{,}718\,281\,8\dots$, che definiremo nel capitolo sulle successioni come limite di
$(1+1/n)^{n}$. In tal caso ovviamente si indica la funzione esponenziale con
$f(x)=e^{x}$; dato che $e>1$, la funzione $e^{x}$ è strettamente crescente.

\begin{notastorica}[Il numero $e$ e gli interessi composti]
Il numero $e$ compare per la prima volta, senza nome, nel 1618, in una tavola di
logaritmi naturali pubblicata in appendice a un'edizione dell'opera di Nepero (John
Napier, 1550--1617), probabilmente redatta da William Oughtred; fu Jakob Bernoulli, nel
1683, a incontrarlo studiando il
problema dell'\emph{interesse composto}: un capitale $C$ investito al tasso annuo del
$100\%$, con interessi capitalizzati $n$ volte l'anno, dopo un anno diventa
$C(1+1/n)^{n}$, e Bernoulli osservò che al crescere di $n$ questa quantità non cresce
all'infinito ma si avvicina a un numero compreso tra $2$ e $3$. La lettera $e$ compare
per la prima volta in una lettera di Eulero a Christian Goldbach del 1731, e sempre
Eulero dimostrò nel 1737 che $e$ è irrazionale. La trascendenza di $e$ --- il fatto che non sia radice di alcun
polinomio a coefficienti interi --- fu provata da Charles Hermite nel 1873.
\end{notastorica}

\subsection{La funzione logaritmo}

Se $a=1$, la funzione $a^{x}$ è identicamente uguale a $1$: è costante e, naturalmente,
non invertibile. Se invece $a\neq1$, la funzione esponenziale $a^{x}$ è strettamente
monotòna e quindi invertibile; si dimostra inoltre (\cref{teo:esponenziale}) che la sua
immagine è costituita da \emph{tutti} i numeri reali positivi.
La funzione inversa è quindi definita sui numeri reali positivi.

\begin{definizione}[Logaritmo]\index{logaritmo}
Sia $a>0$, $a\neq1$. La funzione inversa dell'esponenziale $x\mapsto a^{x}$, definita
per $x>0$ e a valori in $\R$, si chiama \emph{funzione logaritmo in base $a$} e si
scrive $f(x)=\log_{a}x$. Essa è dunque caratterizzata da
\begin{equation}\label{eq:log-def}
  y=\log_{a}x \ \sse\ a^{y}=x \qquad (x>0,\ y\in\R).
\end{equation}
Si suole omettere l'indicazione esplicita della base se tale base è il numero $e$:
\begin{equation}\label{eq:ln-def}
  y=\log x \ \sse\ e^{y}=x .
\end{equation}
\end{definizione}

In parole: $\log_{a}x$ è \emph{l'esponente da dare ad $a$ per ottenere $x$}. Le formule
\eqref{eq:log-def}, \eqref{eq:ln-def} sono molto usate anche nella forma seguente, che
non è altro che la \eqref{eq:inversa} scritta per questa coppia di funzioni:
\begin{equation}\label{eq:log-inversa}
\begin{aligned}
  a^{\log_{a}x}&=x, & e^{\log x}&=x & &(x>0); \\
  \log_{a}(a^{y})&=y, & \log(e^{y})&=y & &(y\in\R).
\end{aligned}
\end{equation}

\begin{proposizione}\label{prop:log-monotono}
Se la base $a$ è maggiore di $1$, il logaritmo $\log_{a}x$ è una funzione strettamente
crescente; se $0<a<1$ è strettamente decrescente.
\end{proposizione}

\begin{schema}
\Ipotesi{$a>1$ (oppure $0<a<1$); \ $0<x_{1}<x_{2}$.}
\Tesi{$\log_{a}x_{1}<\log_{a}x_{2}$ (rispettivamente $\log_{a}x_{1}>\log_{a}x_{2}$).}
\end{schema}

\begin{strategia}
È un caso particolare del \cref{teo:monotona-invertibile}: l'inversa di una funzione
strettamente crescente è strettamente crescente. Ne ripetiamo la dimostrazione per
assurdo in questo caso concreto, perché è istruttivo vedere come la monotonia
dell'esponenziale si «trasferisce» al logaritmo.
\end{strategia}

\begin{dimostrazione}
Siano $x_{1}<x_{2}$ e $y_{1}=\log_{a}x_{1}$, $y_{2}=\log_{a}x_{2}$. Se fosse
$y_{2}<y_{1}$, per la \eqref{eq:pot5} avremmo $x_{2}=a^{y_{2}}<a^{y_{1}}=x_{1}$,
contrariamente alle ipotesi. Analogamente è assurdo che $y_{1}=y_{2}$, perché avremmo
$x_{1}=x_{2}$. Quindi deve risultare $y_{1}<y_{2}$. Con dimostrazione analoga (usando
la \eqref{eq:pot6}) si verifica che il logaritmo è strettamente decrescente se la base è
minore di $1$.
\end{dimostrazione}

\begin{figure}[htb]
\centering
\begin{tikzpicture}
\begin{axis}[width=6cm, height=5cm, xmin=-0.5, xmax=4.5, ymin=-2.5, ymax=2.2,
  xtick={1}, ytick={1}, title={\footnotesize $\log_{a}x$, $a>1$}, title style={at={(0.7,0.9)}}]
  \addplot[curva, domain=0.12:4.4, samples=150] {ln(x)/ln(2)};
\end{axis}
\end{tikzpicture}\hfill
\begin{tikzpicture}
\begin{axis}[width=6cm, height=5cm, xmin=-0.5, xmax=4.5, ymin=-2.2, ymax=2.5,
  xtick={1}, ytick={1}, title={\footnotesize $\log_{a}x$, $0<a<1$}, title style={at={(0.7,0.9)}}]
  \addplot[curva, domain=0.12:4.4, samples=150] {-ln(x)/ln(2)};
\end{axis}
\end{tikzpicture}
\caption{La funzione logaritmo $y=\log_{a}x$: $a=2$ (a sinistra) e $a=\tfrac12$ (a
destra). Sono definite solo per $x>0$, passano per $(1,0)$ e sono le simmetriche
rispetto alla bisettrice delle esponenziali della \cref{fig:esponenziali}.}
\label{fig:logaritmi}
\end{figure}

\needspace{8\baselineskip}
\begin{proposizione}[Proprietà dei logaritmi]\label{prop:log-proprieta}
Sia $a>0$, $a\neq1$. Per ogni $x,x_{1},x_{2}>0$ e ogni $b\in\R$, e per ogni altra base
$c>0$, $c\neq1$, valgono:
\begin{align}
  \log_{a}(x_{1}\cdot x_{2}) &= \log_{a}x_{1}+\log_{a}x_{2}; \label{eq:log1}\\
  \log_{a}(x_{1}/x_{2})      &= \log_{a}x_{1}-\log_{a}x_{2}; \label{eq:log2}\\
  \log_{a}x^{b}              &= b\log_{a}x; \label{eq:log3}\\
  \log_{c}x                  &= \frac{\log_{a}x}{\log_{a}c} \qquad\text{(cambiamento di base).} \label{eq:log4}
\end{align}
\end{proposizione}

\begin{schema}
\Ipotesi{$a,c>0$, \ $a,c\neq1$; \ $x,x_{1},x_{2}>0$; \ $b\in\R$.}
\Tesi{Le relazioni \eqref{eq:log1}--\eqref{eq:log4}.}
\end{schema}

\begin{strategia}
Ogni proprietà del logaritmo è la «traduzione» di una proprietà delle potenze
\eqref{eq:pot1}, letta attraverso la definizione \eqref{eq:log-def}. Il metodo è sempre
lo stesso: si dà un nome agli esponenti, $y_{1}=\log_{a}x_{1}$ e $y_{2}=\log_{a}x_{2}$;
si riscrive tutto in forma esponenziale, $x_{1}=a^{y_{1}}$, $x_{2}=a^{y_{2}}$; si
applica la proprietà delle potenze; si torna ai logaritmi. Lo facciamo per la
\eqref{eq:log1} e lasciamo le altre al lettore.
\end{strategia}

\begin{dimostrazione}
Per provare la \eqref{eq:log1} poniamo
\[
  \log_{a}x_{1}=y_{1}, \qquad \log_{a}x_{2}=y_{2};
\]
in base alla definizione \eqref{eq:log-def} risulta $a^{y_{1}}=x_{1}$ e
$a^{y_{2}}=x_{2}$. Pertanto, per la prima delle \eqref{eq:pot1},
\[
  x_{1}\cdot x_{2}=a^{y_{1}}\cdot a^{y_{2}}=a^{y_{1}+y_{2}},
\]
che, di nuovo in base alla definizione \eqref{eq:log-def}, equivale a
\[
  y_{1}+y_{2}=\log_{a}(x_{1}\cdot x_{2}).
\]
Ricordando i simboli introdotti all'inizio, quest'ultima è la tesi \eqref{eq:log1}. Le
altre relazioni si dimostrano in modo analogo: per la \eqref{eq:log3} si parte da
$x^{b}=(a^{y})^{b}=a^{by}$; per la \eqref{eq:log4} si pone $y=\log_{c}x$, cioè
$x=c^{y}$, e si prende il logaritmo in base $a$ di entrambi i membri usando la
\eqref{eq:log3}: $\log_{a}x=y\log_{a}c$. Il lettore è invitato a eseguire i calcoli
esplicitamente.
\end{dimostrazione}

\begin{notastorica}[I logaritmi: trasformare prodotti in somme]
Le proprietà \eqref{eq:log1}--\eqref{eq:log3} non sono curiosità: sono il motivo per cui
i logaritmi furono inventati. All'inizio del Seicento astronomi e navigatori dovevano
eseguire a mano moltiplicazioni tra numeri di molte cifre; Nepero (1614) e, poco dopo,
Henry Briggs (1617, con la base 10) pubblicarono tavole che permettevano di sostituire
ogni prodotto con una somma: si cercano i logaritmi dei due fattori, si sommano, e si
cerca nella tavola il numero che ha quel logaritmo. Laplace disse che i logaritmi,
«abbreviando le fatiche, raddoppiarono la vita degli astronomi». Il regolo calcolatore,
usato dagli ingegneri fino agli anni Settanta del Novecento (anche per progettare
aerei e missioni spaziali), è una tavola di logaritmi incisa su due righelli
scorrevoli: sommare lunghezze proporzionali a $\log x_{1}$ e $\log x_{2}$ dà una
lunghezza proporzionale a $\log(x_{1}x_{2})$.
\end{notastorica}

\begin{python}[Esponenziali, logaritmi e le loro proprietà]
import math
# La definizione: log_a(x) è l'esponente y tale che a**y == x
a, x = 2.0, 10.0
y = math.log(x, a)                     # logaritmo in base a
print(f"log_2(10) = {y:.6f},   2**{y:.6f} = {a**y:.6f}")

# Proprietà (log1), (log3), (log4) verificate numericamente
x1, x2, b = 3.0, 7.0, 2.5
print(math.isclose(math.log(x1*x2), math.log(x1) + math.log(x2)))
print(math.isclose(math.log(x1**b), b*math.log(x1)))
print(math.isclose(math.log(x1, 10), math.log(x1)/math.log(10)))

# Il numero di Nepero come limite di (1 + 1/n)**n
for n in (1, 10, 100, 1000, 10**6):
    print(f"n = {n:>8d}   (1+1/n)^n = {(1 + 1/n)**n:.8f}")
print(f"math.e         = {math.e:.8f}")
\end{python}
\begin{risultato}
log_2(10) = 3.321928,   2**3.321928 = 10.000000
True
True
True
n =        1   (1+1/n)^n = 2.00000000
n =       10   (1+1/n)^n = 2.59374246
n =      100   (1+1/n)^n = 2.70481383
n =     1000   (1+1/n)^n = 2.71692393
n =  1000000   (1+1/n)^n = 2.71828047
math.e         = 2.71828183
\end{risultato}

\begin{esercizio}
Risolvere l'equazione $2^{x+1}=3^{x}$.
\begin{soluzione}
Prendendo il logaritmo (in qualsiasi base, per esempio naturale) di ambo i membri, che
sono positivi, e usando le \eqref{eq:log3} e \eqref{eq:log2}:
$(x+1)\log2=x\log3$, da cui $x(\log3-\log2)=\log2$ e
$x=\log2/(\log3-\log2)=\log2/\log(3/2)\approx1{,}71$. Il passaggio al logaritmo è
un'equivalenza: da $u=v$ segue $\log u=\log v$ perché $\log$ è una funzione, e da
$\log u=\log v$ segue $u=v$ \emph{perché il logaritmo è iniettivo}
(\cref{prop:log-monotono}). \emph{Scelta logica:} l'incognita compare all'esponente, e
il logaritmo, che trasforma le potenze in prodotti \eqref{eq:log3}, la «porta giù»,
riducendo il problema a un'equazione di primo grado.
\end{soluzione}
\end{esercizio}

\begin{esercizio}
Dimostrare che $\log_{a}(1/x)=-\log_{a}x$ e che $\log_{1/a}x=-\log_{a}x$.
\begin{soluzione}
Poiché $a^{0}=1$, è $\log_{a}1=0$; per la \eqref{eq:log2} con $x_{1}=1$ e $x_{2}=x$ si
ha allora $\log_{a}(1/x)=\log_{a}1-\log_{a}x=-\log_{a}x$. In particolare
$\log_{a}(1/a)=-\log_{a}a=-1$. Per la seconda identità usiamo il cambiamento di base
\eqref{eq:log4} con $c=1/a$, lecito perché $1/a>0$ e $1/a\neq1$:
$\log_{1/a}x=\log_{a}x/\log_{a}(1/a)=-\log_{a}x$. \emph{Scelta logica:} entrambe le
identità sono casi particolari delle proprietà generali \eqref{eq:log2} e
\eqref{eq:log4}: conviene riconoscere la struttura (un quoziente, un cambio di base)
invece di ripartire ogni volta dalla definizione.
\end{soluzione}
\end{esercizio}

%% ----------------------------------------------------------------------------
%%  §1.9  Le funzioni trigonometriche
%% ----------------------------------------------------------------------------
\section{Le funzioni trigonometriche}
\label{sec:trigonometriche}

Riportiamo in questo paragrafo un breve riassunto di nozioni di trigonometria, in
genere già note al lettore; lo scopo è fissare le notazioni e, soprattutto, mettere in
evidenza le poche proprietà che useremo davvero nel seguito.

\subsection{La misura degli angoli in radianti}

Il lettore sa misurare gli angoli in gradi: un angolo retto misura $90^{\circ}$, un
angolo piatto $180^{\circ}$, un angolo giro $360^{\circ}$. Per studiare le funzioni
trigonometriche è opportuno adottare una diversa unità di misura.

\begin{definizione}[Radiante]\index{radiante}
La \emph{misura in radianti} di un angolo piano è la \emph{lunghezza dell'arco di
circonferenza di raggio $1$ e centro nel vertice dell'angolo intercettato dalle due
semirette che individuano l'angolo} (\cref{fig:radianti}).
\end{definizione}

Si conviene di denotare con $\pi$ (\emph{pi greco}) la lunghezza di una
semicirconferenza di raggio $1$. Ciò significa che, espresso in radianti, un angolo
piatto misura $\pi$, un angolo retto misura $\pi/2$, mentre un angolo giro misura
$2\pi$ (la lunghezza di una circonferenza di raggio $1$ è $2\pi$). Come tutti sanno, un
valore numerico approssimato di $\pi$ è $3{,}14\ldots$; la conversione tra gradi e
radianti è una proporzione: $\alpha^{\circ}:360^{\circ}=x:2\pi$.

\begin{idea}[Perché i radianti]
La scelta dei radianti non è un vezzo dei matematici. Il grado è un'unità arbitraria
(360 è un'eredità dell'astronomia babilonese), mentre il radiante è una misura
\emph{intrinseca}: un rapporto tra due lunghezze (arco e raggio), dunque un numero puro.
La ricompensa arriverà nel calcolo differenziale: la derivata di $\sen x$ è $\cos x$
\emph{solo se} $x$ è in radianti; in gradi comparirebbe un fastidioso fattore
$\pi/180$. Per un ingegnere il radiante è anche l'unità naturale per le velocità
angolari ($\omega$ in rad/s) e per la relazione $v=\omega r$ tra velocità periferica e
raggio.
\end{idea}

Analogamente a quanto si fa per l'ascissa su una retta, si definisce un'origine e un
verso di rotazione positivo (si suole scegliere il verso antiorario a partire dall'asse
delle $x$) e si considerano in modo naturale anche angoli maggiori di $2\pi$ radianti, o
angoli minori di zero. Così l'angolo $x$ geometricamente corrisponde all'angolo
$x+2\pi$ e anche all'angolo $x+4\pi$, oppure all'angolo $x-2\pi$, o in generale
all'angolo $x+2k\pi$, per ogni $k\in\Z$.

\begin{definizione}[Seno e coseno]\index{seno}\index{coseno}
Le funzioni $\sen x$ e $\cos x$ si definiscono rispettivamente come \emph{l'ordinata} e
\emph{l'ascissa} del punto che si trova sulla circonferenza di centro l'origine e raggio
$1$ e che sottende un angolo orientato di lunghezza $x$, a partire dall'asse delle
ascisse (\cref{fig:radianti}).
\end{definizione}

\begin{figure}[htb]
\centering
\begin{tikzpicture}[scale=1.55]
  \draw[-Latex, grigio] (-1.3,0) -- (1.5,0) node[below, font=\scriptsize] {$x$};
  \draw[-Latex, grigio] (0,-1.3) -- (0,1.5) node[left, font=\scriptsize] {$y$};
  \draw[grigio!70] (0,0) circle (1);
  \node[below right, font=\scriptsize, grigio] at (1,0) {$1$};
  \pgfmathsetmacro{\ang}{40}
  \draw[rame, very thick] (1,0) arc (0:\ang:1) node[midway, right=2pt, font=\scriptsize] {$x$};
  \draw[blu, thick] (0,0) -- (\ang:1);
  \fill[blu] (\ang:1) circle (1.2pt) node[above right, font=\scriptsize] {$P=(\cos x,\sen x)$};
  \draw[tratteggio] (\ang:1) -- ({cos(\ang)},0);
  \draw[tratteggio] (\ang:1) -- (0,{sin(\ang)});
  \draw[verde, very thick] (0,0) -- ({cos(\ang)},0) node[midway, below, font=\scriptsize] {$\cos x$};
  \draw[verde, very thick] ({cos(\ang)},0) -- ({cos(\ang)},{sin(\ang)}) node[pos=0.35, left, font=\scriptsize] {$\sen x$};
\end{tikzpicture}\hspace{8mm}
\begin{tikzpicture}[scale=1.55]
  \draw[-Latex, grigio] (-1.3,0) -- (1.5,0) node[below, font=\scriptsize] {$x$};
  \draw[-Latex, grigio] (0,-1.3) -- (0,1.5) node[left, font=\scriptsize] {$y$};
  \draw[grigio!70] (0,0) circle (1);
  \pgfmathsetmacro{\ang}{140}
  \draw[rame, very thick] (1,0) arc (0:\ang:1) node[pos=0.5, above right=1pt, font=\scriptsize] {$x$};
  \draw[blu, thick] (0,0) -- (\ang:1);
  \fill[blu] (\ang:1) circle (1.2pt) node[above left, font=\scriptsize] {$P$};
  \draw[tratteggio] (\ang:1) -- ({cos(\ang)},0);
  \draw[verde, very thick] (0,0) -- ({cos(\ang)},0) node[midway, below, font=\scriptsize] {$\cos x<0$};
  \draw[verde, very thick] ({cos(\ang)},0) -- ({cos(\ang)},{sin(\ang)}) node[midway, left, font=\scriptsize] {$\sen x$};
\end{tikzpicture}
\caption{Definizione di $\sen x$ e $\cos x$ sulla circonferenza di raggio $1$: $x$ è la
lunghezza dell'arco (in rame) percorso in verso antiorario a partire dal punto $(1,0)$.
A destra, un angolo del secondo quadrante: l'ascissa, e quindi il coseno, è negativa.}
\label{fig:radianti}
\end{figure}

Le funzioni $\sen x$, $\cos x$ sono definite per ogni $x\in\R$, mentre l'immagine delle
due funzioni è compresa tra $-1$ e $1$, cioè:
\begin{equation}\label{eq:sen-cos-limitate}
  -1\le\sen x\le1, \qquad -1\le\cos x\le1, \qquad\perogni x\in\R.
\end{equation}
Inoltre, poiché agli angoli $x$ e $x+2\pi$ corrisponde lo stesso punto $P$, le due
funzioni sono \emph{periodiche} di periodo $2\pi$:
$\sen(x+2\pi)=\sen x$ e $\cos(x+2\pi)=\cos x$. Il grafico delle due funzioni è riportato
nella \cref{fig:sen-cos}. Nel capitolo sullo studio di funzione sarà indicato come
ottenere il grafico; per ora il lettore controlli sul grafico il segno delle due funzioni
in base alla definizione. Le funzioni $\sen x$ e $\cos x$ non sono monotòne su $\R$.

\begin{figure}[htb]
\centering
\begin{tikzpicture}
\begin{axis}[width=\textwidth, height=4.8cm, xmin=-6.9, xmax=6.9, ymin=-1.4, ymax=1.5,
  xtick={-6.2832,-3.1416,-1.5708,1.5708,3.1416,4.7124,6.2832},
  xticklabels={$-2\pi$,$-\pi$,$-\frac{\pi}{2}$,$\frac{\pi}{2}$,$\pi$,$\frac{3}{2}\pi$,$2\pi$},
  ytick={-1,1}, samples=300, axis on top, xticklabel style={fill=white, inner sep=1pt}]
  \addplot[curva, domain=-6.8:6.8] {sin(deg(x))};
  \addplot[curva2, dashed, domain=-6.8:6.8] {cos(deg(x))};
  \node[blu, font=\scriptsize] at (axis cs:2.3,1.2) {$\sen x$};
  \node[rame, font=\scriptsize] at (axis cs:-1.0,1.2) {$\cos x$};
\end{axis}
\end{tikzpicture}
\caption{I grafici di $\sen x$ (linea continua) e $\cos x$ (tratteggiata). Il coseno è
il seno «in anticipo» di $\pi/2$: $\cos x=\sen(x+\pi/2)$.}
\label{fig:sen-cos}
\end{figure}

\subsection{La relazione fondamentale e il teorema di Pitagora}

Delle numerose relazioni tra queste due funzioni, la più importante è senza dubbio
quella che segue dal teorema di Pitagora sui triangoli rettangoli. È utile ricordare
come si ottiene una dimostrazione del teorema di Pitagora, perché è un piccolo
capolavoro di economia.

\begin{teorema}[Pitagora]\label{teo:pitagora}\index{teorema di Pitagora}
In un triangolo rettangolo con cateti lunghi $a$, $b$ e ipotenusa lunga $c$ risulta
\begin{equation}\label{eq:pitagora}
  a^{2}+b^{2}=c^{2}.
\end{equation}
\end{teorema}

\begin{schema}
\Ipotesi{Triangolo rettangolo con cateti lunghi $a$, $b$ e ipotenusa lunga $c$.}
\Tesi{$a^{2}+b^{2}=c^{2}$.}
\end{schema}

\begin{strategia}
L'idea è di calcolare \emph{una stessa area in due modi}. Disponiamo quattro copie del
triangolo dentro un quadrato di lato $a+b$ come nella \cref{fig:pitagora}: la regione
che resta scoperta è un quadrilatero di lato $c$, e un piccolo ragionamento sugli
angoli mostra che è un quadrato. Allora l'area del quadrato grande, $(a+b)^{2}$, è
uguale a quattro triangoli più il quadrato di lato $c$: sviluppando e semplificando
resta proprio la \eqref{eq:pitagora}.
\end{strategia}

\begin{dimostrazione}
Si veda la \cref{fig:pitagora}. Il triangolo rettangolo, ripetuto quattro volte, ha
cateti lunghi $a$, $b$ e angoli $\alpha$, $\beta$, oltre all'angolo retto. Evidentemente
$\alpha+\beta=\pi/2$ (perché la somma degli angoli interni del triangolo vale $\pi$);
ne segue che anche l'angolo $\gamma$ in figura è uguale a $\pi/2$, essendo
$\alpha+\gamma+\beta=\pi$ (angolo piatto). Perciò il rombo di lato $c$ rappresentato nella
parte destra della figura è in realtà un quadrato. L'area del quadrato più grande, di
lato $a+b$, vale $(a+b)^{2}$. Tale area si può ottenere sommando le aree delle figure
componenti (quattro triangoli rettangoli di cateti $a$, $b$, e un quadrato di lato
$c$), cioè $4\cdot(ab/2)+c^{2}$. Perciò
\[
  (a+b)^{2}=a^{2}+2ab+b^{2}=2ab+c^{2},
\]
da cui, semplificando $2ab$ (\cref{prop:semplificazione-somma}), il teorema di Pitagora
\eqref{eq:pitagora}.
\end{dimostrazione}

\begin{figure}[htb]
\centering
\begin{tikzpicture}[scale=0.85, every node/.style={font=\scriptsize}]
  \pgfmathsetmacro{\a}{1.4}\pgfmathsetmacro{\b}{2.4}
  \pgfmathsetmacro{\s}{\a+\b}
  % quadrato di lato a+b con il quadrato inclinato di lato c
  \draw[grigio, thick] (0,0) rectangle (\s,\s);
  \fill[bluchiaro] (\a,0) -- (\s,\a) -- (\b,\s) -- (0,\b) -- cycle;
  \draw[blu, thick] (\a,0) -- (\s,\a) -- (\b,\s) -- (0,\b) -- cycle;
  \draw[rame, thick] (0,0) -- (\a,0) -- (0,\b) -- cycle;
  \node[below] at (\a/2,0) {$a$}; \node[below] at (\a+\b/2,0) {$b$};
  \node[left] at (0,\b/2) {$b$}; \node[left] at (0,\b+\a/2) {$a$};
  \node[right] at (\s,\a/2) {$a$}; \node[right] at (\s,\a+\b/2) {$b$};
  \node[above] at (\b/2,\s) {$b$}; \node[above] at (\b+\a/2,\s) {$a$};
  \node[blu] at (\s/2,\s/2) {$c^{2}$};
  % angoli
  \draw[rame] (\a,0) ++(180:0.35) arc (180:120:0.35); \node[rame] at (\a-0.5,0.25) {$\alpha$};
  \draw[rame] (0,\b) ++(270:0.35) arc (270:300:0.35); \node[rame] at (0.15,\b-0.6) {$\beta$};
  \draw[blu] (\a,0) ++(0:0.35) arc (0:30.26:0.35); \node[blu] at (\a+0.55,0.2) {$\beta$};
  \node[blu] at (\a,0.45) {$\gamma$};
\end{tikzpicture}
\caption{La dimostrazione «per aree» del teorema di Pitagora. Il quadrato grande ha lato
$a+b$; i quattro triangoli rettangoli (uno è evidenziato in rame) lasciano scoperto il
quadrato azzurro di lato $c$. Poiché $\alpha+\beta=\pi/2$, l'angolo $\gamma$ è retto.}
\label{fig:pitagora}
\end{figure}

Applicando il teorema di Pitagora al triangolo rettangolo di cateti $\abs{\sen x}$ e
$\abs{\cos x}$ e ipotenusa uguale a $1$ (\cref{fig:radianti}), si trova la
\emph{relazione fondamentale} (se il triangolo degenera, cioè se $\sen x=0$ oppure
$\cos x=0$, il punto $P$ sta su un asse, l'altra coordinata vale $\pm1$ e la relazione
è ovvia)\index{relazione fondamentale della trigonometria}
\begin{equation}\label{eq:relazione-fondamentale}
  \sen^{2}x+\cos^{2}x=1, \qquad\perogni x\in\R.
\end{equation}
Sono anche importanti le \emph{formule di addizione}\index{formule di addizione}, che
riportiamo senza dimostrazione:
\begin{align}
  \sen(x_{1}\pm x_{2}) &= \sen x_{1}\cos x_{2}\pm\sen x_{2}\cos x_{1}, \label{eq:addizione-sen}\\
  \cos(x_{1}\pm x_{2}) &= \cos x_{1}\cos x_{2}\mp\sen x_{1}\sen x_{2}. \label{eq:addizione-cos}
\end{align}
Se nelle identità sopra scritte scegliamo il segno $+$ e poniamo $x_{1}=x_{2}=x$,
otteniamo le \emph{formule di duplicazione}:
\begin{equation}\label{eq:duplicazione}
  \sen2x=2\sen x\cos x, \qquad \cos2x=\cos^{2}x-\sen^{2}x .
\end{equation}
Dalle \eqref{eq:addizione-sen} e \eqref{eq:addizione-cos} con $x_{1}=0$ e il segno $-$
si ricava anche che il coseno è una funzione \emph{pari} e il seno è \emph{dispari}: $\cos(-x)=\cos x$,
$\sen(-x)=-\sen x$, come si vede del resto dalla simmetria della circonferenza rispetto
all'asse $x$.

\subsection{La funzione tangente}

\begin{definizione}[Tangente]\index{tangente}
A partire dalle funzioni $\sen x$ e $\cos x$ si definisce la \emph{funzione tangente}
\begin{equation}\label{eq:tangente}
  \tg x=\frac{\sen x}{\cos x},
\end{equation}
definita se $\cos x\neq0$, cioè se $x\neq\pi/2+k\pi$, con $k\in\Z$.
\end{definizione}

\begin{figure}[htb]
\centering
\begin{tikzpicture}
\begin{axis}[width=7.2cm, height=5.6cm, xmin=-5, xmax=5, ymin=-3.5, ymax=3.5,
  xtick={-4.7124,-3.1416,-1.5708,1.5708,3.1416,4.7124},
  xticklabels={$-\frac32\pi$,$-\pi$,$-\frac{\pi}{2}$,$\frac{\pi}{2}$,$\pi$,$\frac32\pi$},
  ytick=\empty, samples=200, axis on top, xticklabel style={fill=white, inner sep=1pt}]
  \foreach \k in {-1,0,1} {
    \addplot[curva, domain={-1.29+\k*3.1416}:{1.29+\k*3.1416}] {tan(deg(x))};}
  \foreach \v in {-4.7124,-1.5708,1.5708,4.7124} {\addplot[tratteggio] coordinates {(\v,-3.5) (\v,3.5)};}
  \node[font=\scriptsize, blu] at (axis cs:0.6,2.6) {$\tg x$};
\end{axis}
\end{tikzpicture}\hfill
\begin{tikzpicture}[scale=1.5]
  \draw[-Latex, grigio] (-0.3,0) -- (1.9,0);
  \draw[-Latex, grigio] (0,-0.3) -- (0,1.6);
  \draw[grigio!70] (1,0) arc (0:90:1);
  \pgfmathsetmacro{\ang}{50}
  \draw[blu, thick] (0,0) -- ({1/cos(\ang)*cos(\ang)},{1/cos(\ang)*sin(\ang)});
  \draw[rame, very thick] (1,0) arc (0:\ang:1) node[midway, right=2pt, font=\scriptsize] {$x$};
  \fill[blu] (\ang:1) circle (1.1pt) node[above left, font=\scriptsize] {$P$};
  \draw[tratteggio] (\ang:1) -- ({cos(\ang)},0) node[below, font=\scriptsize] {$A$};
  \draw[verde, very thick] (1,0) -- (1,{tan(\ang)}) node[midway, right, font=\scriptsize] {$\tg x$};
  \fill[verde] (1,{tan(\ang)}) circle (1.1pt) node[above right, font=\scriptsize] {$T$};
  \node[below, font=\scriptsize] at (1,0) {$B$};
  \node[below left, font=\scriptsize] at (0,0) {$O$};
  \node[left, font=\scriptsize] at (0,1) {$1$};
  \draw[grigio] (0,1) -- (-0.04,1);
  \node[font=\scriptsize, grigio, anchor=east] at ({cos(\ang)-0.02},{sin(\ang)*0.2}) {$\sen x$};
\end{tikzpicture}
\caption{A sinistra: il grafico di $\tg x$, con gli asintoti verticali
$x=\pi/2+k\pi$. A destra: la tangente come lunghezza del segmento $BT$ sulla retta
tangente alla circonferenza nel punto $B=(1,0)$.}
\label{fig:tangente}
\end{figure}

Usando le proprietà dei triangoli simili ($OAP$ e $OBT$ nella \cref{fig:tangente}),
$BT=BT/OB=AP/OA=\sen x/\cos x$, si vede che la tangente di un angolo $x$ si può
rappresentare come lunghezza del segmento $BT$. Dalla figura risulta anche chiaro che,
se $0<x<\pi/2$, allora valgono le disuguaglianze
\begin{equation}\label{eq:sen-x-tg}
  0<\sen x<x<\tg x,
\end{equation}
perché la corda è più corta dell'arco (il segmento $AP$ è un cateto del triangolo
rettangolo $APB$, quindi è più corto dell'ipotenusa $PB$, cioè della corda, a sua volta
più corta dell'arco $x$) e l'area del settore circolare $OBP$, che vale $x/2$, è
minore dell'area del triangolo $OBT$, che vale $\tg x/2$. La \eqref{eq:sen-x-tg} è
l'unico fatto «geometrico» di cui avremo bisogno nel seguito: da essa discende il
limite fondamentale $\lim_{x\to0}(\sen x)/x=1$, e da questo tutte le derivate delle
funzioni trigonometriche. La tangente è periodica di periodo $\pi$ e, in ciascun
intervallo $(-\pi/2+k\pi,\pi/2+k\pi)$, è strettamente crescente.

Riportiamo infine una tabella con i valori delle funzioni trigonometriche per alcuni
angoli di uso frequente. Il modo con cui ottenere tali valori è indicato
nell'\cref{es:valori-notevoli}.

\begin{table}[htb]
\centering
\renewcommand{\arraystretch}{1.35}
\setlength{\tabcolsep}{5pt}
\begin{tabular}{l*{8}{c}}
\toprule
$x$ (radianti) & $0$ & $\pi/6$ & $\pi/4$ & $\pi/3$ & $\pi/2$ & $\pi$ & $\tfrac32\pi$ & $2\pi$ \\
$x$ (gradi) & $0$ & $30$ & $45$ & $60$ & $90$ & $180$ & $270$ & $360$ \\
\midrule
$\sen x$ & $0$ & $\tfrac12$ & $\tfrac{\sqrt2}{2}$ & $\tfrac{\sqrt3}{2}$ & $1$ & $0$ & $-1$ & $0$ \\
$\cos x$ & $1$ & $\tfrac{\sqrt3}{2}$ & $\tfrac{\sqrt2}{2}$ & $\tfrac12$ & $0$ & $-1$ & $0$ & $1$ \\
$\tg x$ & $0$ & $\tfrac{\sqrt3}{3}$ & $1$ & $\sqrt3$ & non def. & $0$ & non def. & $0$ \\
\bottomrule
\end{tabular}
\caption{Valori notevoli delle funzioni trigonometriche.}
\label{tab:valori-notevoli}
\end{table}

\begin{esercizio}\label{es:valori-notevoli}
Ricavare i valori della \cref{tab:valori-notevoli} per $x=\pi/4$ e $x=\pi/6$.
\begin{soluzione}
Per $x=\pi/4$ il triangolo $OAP$ della \cref{fig:radianti} è rettangolo isoscele,
quindi $\sen x=\cos x$; dalla relazione fondamentale $2\sen^{2}x=1$, cioè
$\sen x=\cos x=\sqrt2/2$ (positivi perché $P$ è nel primo quadrante), e $\tg x=1$. Per
$x=\pi/6$: il triangolo di vertici $O$, $P$ e il simmetrico $P'$ di $P$ rispetto
all'asse $x$ ha due lati uguali a $1$ e l'angolo in $O$ uguale a $\pi/3$: è
equilatero, dunque $PP'=1$ e $\sen(\pi/6)=PP'/2=1/2$. Dalla relazione fondamentale
$\cos(\pi/6)=\sqrt{1-1/4}=\sqrt3/2$ e $\tg(\pi/6)=(1/2)/(\sqrt3/2)=1/\sqrt3=\sqrt3/3$.
I valori per $\pi/3$ seguono scambiando seno e coseno, perché $\pi/3=\pi/2-\pi/6$ e
$\sen(\pi/2-x)=\cos x$ (formula di addizione). \emph{Scelta logica:} i valori notevoli
si ricavano da figure con molte simmetrie (il triangolo rettangolo isoscele, il triangolo
equilatero), nelle quali una relazione geometrica elementare fornisce un'equazione per
$\sen x$; la relazione fondamentale dà poi il coseno, con il segno fissato dal quadrante.
\end{soluzione}
\end{esercizio}

\begin{python}[Verifica delle identità trigonometriche e delle disuguaglianze \eqref{eq:sen-x-tg}]
import math
# Relazione fondamentale e duplicazione, per angoli "a caso"
for x in (0.3, 1.0, 2.5, -4.0):
    print(math.isclose(math.sin(x)**2 + math.cos(x)**2, 1.0),
          math.isclose(math.sin(2*x), 2*math.sin(x)*math.cos(x)))

# Le disuguaglianze 0 < sen x < x < tg x per 0 < x < pi/2
for x in (0.1, 0.5, 1.0, 1.5):
    s, t = math.sin(x), math.tan(x)
    print(f"x={x:.1f}:  sen x={s:.4f} < x={x:.4f} < tg x={t:.4f}")

# I valori notevoli: math.radians converte i gradi in radianti
for gradi in (0, 30, 45, 60, 90):
    x = math.radians(gradi)
    print(f"{gradi:3d}°  sen={math.sin(x):.4f}"
          f"  cos={math.cos(x):.4f}")
\end{python}
\begin{risultato}
True True
True True
True True
True True
x=0.1:  sen x=0.0998 < x=0.1000 < tg x=0.1003
x=0.5:  sen x=0.4794 < x=0.5000 < tg x=0.5463
x=1.0:  sen x=0.8415 < x=1.0000 < tg x=1.5574
x=1.5:  sen x=0.9975 < x=1.5000 < tg x=14.1014
  0°  sen=0.0000  cos=1.0000
 30°  sen=0.5000  cos=0.8660
 45°  sen=0.7071  cos=0.7071
 60°  sen=0.8660  cos=0.5000
 90°  sen=1.0000  cos=0.0000
\end{risultato}

\begin{attenzione}[$\cos(90^{\circ})$ non è esattamente $0$ al calcolatore]
Nell'ultima riga dell'output il coseno di $90^{\circ}$ appare come $0{,}0000$, ma il
valore calcolato da Python è $6{,}1\cdot10^{-17}$: il numero $\pi/2$ non è
rappresentabile esattamente in virgola mobile, e il calcolatore sta valutando il coseno
di un numero \emph{vicinissimo} a $\pi/2$, ma non uguale. Confrontare numeri in virgola
mobile con \texttt{==} è quasi sempre un errore; per questo nel codice usiamo
\texttt{math.isclose}.
\end{attenzione}

%% ----------------------------------------------------------------------------
%%  §1.10  Appendice: il principio di induzione
%% ----------------------------------------------------------------------------
\section{Appendice: il principio di induzione}
\label{sec:induzione}

Abbiamo utilizzato nel \cref{sec:potenze} la seguente affermazione sulla crescenza
della funzione potenza $x^{n}$:
\begin{equation}\label{eq:potenza-crescente-bis}
  0\le x_{1}<x_{2} \ \implica\ x_{1}^{n}<x_{2}^{n}, \qquad\perogni n\in\N,
\end{equation}
e ne abbiamo dato una dimostrazione per $n=2$ e $n=3$, concludendo con un «e così
via». Vogliamo ora rendere rigoroso quel «così via».

Supponiamo preliminarmente che la \eqref{eq:potenza-crescente-bis} valga per un certo
indice $n$ (ciò che dobbiamo provare è che valga per \emph{tutti} gli $n$; qui stiamo
supponendo che sia vera per \emph{qualche} $n$: ad esempio, per $n=1$ la
\eqref{eq:potenza-crescente-bis} è banalmente vera!). Perciò supponiamo che valgano le
disuguaglianze $0\le x_{1}<x_{2}$, $x_{1}^{n}<x_{2}^{n}$. Otteniamo, moltiplicando la
prima per $x_{1}^{n}\ge0$ e la seconda per $x_{2}>0$:
\begin{equation}\label{eq:passo-potenza}
  x_{1}^{n+1}=x_{1}\,x_{1}^{n}\le x_{2}\,x_{1}^{n}<x_{2}\,x_{2}^{n}=x_{2}^{n+1}.
\end{equation}
Cioè abbiamo provato che, \emph{se} vale la \eqref{eq:potenza-crescente-bis} per un
certo indice $n$, \emph{allora} essa vale anche per l'indice successivo $n+1$. Ma allora
la \eqref{eq:potenza-crescente-bis} vale sempre, perché: sappiamo che la proposizione
vale per $n=1$ (lo verifichiamo banalmente); per quanto sopra detto essa vale anche per
l'indice successivo, cioè $n=2$; ancora, sempre per lo stesso motivo, vale per il
successivo $n=3$, e così via\ldots{} Possiamo raggiungere con questo argomento
qualsiasi naturale $n$. Formuliamo in generale il seguente

\begin{teorema}[Principio di induzione]\label{teo:induzione}\index{principio di induzione}
Supponiamo che una proposizione $\mathcal{P}(n)$ dipendente da un indice $n\in\N$ sia
vera per $n=1$ e che inoltre, supposta vera per $n$, sia vera anche per il successivo
$n+1$. Allora la proposizione è vera per ogni $n\in\N$.
\end{teorema}

\begin{schema}
\Ipotesi{(a) \emph{Base:} $\mathcal{P}(1)$ è vera. \\
(b) \emph{Passo induttivo:} per ogni $n\in\N$, $\mathcal{P}(n)\implica\mathcal{P}(n+1)$.}
\Tesi{$\mathcal{P}(n)$ è vera per ogni $n\in\N$.}
\end{schema}

\begin{strategia}
Tutto sta nel dare un significato preciso alla frase «sommando ripetutamente $1$» con
cui abbiamo definito $\N$ nel \cref{sec:nzq}: $\N$ è il \emph{più piccolo}
sottoinsieme di $\R$ che contiene $1$ e che, se contiene $n$, contiene anche $n+1$.
Basta allora osservare che l'insieme dei naturali per cui $\mathcal{P}(n)$ è vera ha
queste due proprietà: essendo contenuto in $\N$ e non potendo essere più piccolo di
$\N$, coincide con $\N$.
\end{strategia}

\begin{dimostrazione}
Diciamo che un insieme $S\subseteq\R$ è \emph{induttivo} se $1\in S$ e se $n\in S$
implica $n+1\in S$. L'insieme $\R$ è induttivo, e l'intersezione di una famiglia
qualsiasi di insiemi induttivi è ancora un insieme induttivo (se $1$ e $n$ stanno in
tutti, vi stanno anche $1$ e $n+1$). La definizione di $\N$ come insieme dei numeri che
si ottengono da $1$ «sommando ripetutamente $1$» si rende rigorosa ponendo $\N$ uguale
all'intersezione di tutti i sottoinsiemi induttivi di $\R$: così $\N$ è induttivo ed è
contenuto in ogni insieme induttivo.

Sia ora $S=\{n\in\N : \mathcal{P}(n)\text{ è vera}\}$. Per l'ipotesi (a) è $1\in S$.
Se $n\in S$, allora $\mathcal{P}(n)$ è vera e, per l'ipotesi (b), lo è anche
$\mathcal{P}(n+1)$; inoltre $n+1\in\N$, perché $\N$ è induttivo: dunque $n+1\in S$.
Quindi $S$ è induttivo, e perciò $\N\subseteq S$; poiché per costruzione
$S\subseteq\N$, risulta $S=\N$, cioè $\mathcal{P}(n)$ è vera per ogni $n\in\N$.
\end{dimostrazione}

\begin{osservazione}[Il principio di induzione è un assioma o un teorema?]\label{oss:peano}
Nel nostro approccio, che parte dagli assiomi dei numeri reali, il principio di
induzione è un teorema, la cui dimostrazione si riduce a poche righe una volta precisato
che cosa significhi \emph{ripetutamente}. Nell'assiomatizzazione dei numeri naturali
proposta da Giuseppe Peano nel 1889, che non presuppone $\R$, si parte invece da un
insieme $\N$, da un suo elemento $1$ e da una funzione $\sigma$, il «successore», con
gli assiomi:
\begin{enumerate}[label=\textup{(P\arabic*)}, leftmargin=3em]
  \item $1\in\N$;
  \item se $n\in\N$, allora $\sigma(n)\in\N$;
  \item $1$ non è il successore di alcun numero: $\sigma(n)\neq1$ per ogni $n\in\N$;
  \item numeri con lo stesso successore sono uguali: se $\sigma(n)=\sigma(m)$, allora
  $n=m$;
  \item se $S\subseteq\N$ contiene $1$ e contiene $\sigma(n)$ ogni volta che contiene
  $n$, allora $S=\N$.
\end{enumerate}
Per Peano l'induzione è dunque un assioma, il quinto e ultimo. Il nostro $\N\subseteq\R$,
con $\sigma(n)=n+1$, li soddisfa tutti. (P1) e (P2) dicono che $\N$ è induttivo; (P4) è
la legge di semplificazione (\cref{prop:semplificazione-somma}); (P5) è il cuore della
dimostrazione del \cref{teo:induzione}, perché un tale $S$ è induttivo e quindi contiene
$\N$. Per (P3) basta osservare che l'insieme $\{x\in\R : x>0\}$ è induttivo (se $x>0$,
allora $x+1>1>0$), e quindi contiene $\N$: ogni naturale è positivo, mentre $n+1=1$
darebbe $n=0$. I naturali «dentro» $\R$ sono, a tutti gli effetti, i naturali di Peano.
\end{osservazione}

\begin{notastorica}[Peano e l'aritmetica «nova methodo»]
Giuseppe Peano (1858--1932), professore all'Università di Torino, pubblicò nel 1889 un
opuscolo in latino, \emph{Arithmetices principia, nova methodo exposita}, in cui
l'aritmetica è dedotta da pochi assiomi scritti per intero in simboli logici. Gli
assiomi erano nove, ma quattro riguardavano soltanto le proprietà dell'uguaglianza; i
cinque rimanenti sono quelli dell'\cref{oss:peano}, e il primo numero era, come in questo
libro, l'$1$: solo più tardi, nel \emph{Formulario mathematico}, Peano fece cominciare i
naturali dallo zero. Un anno prima Richard Dedekind, in \emph{Was sind und was sollen die
Zahlen?} (1888), aveva caratterizzato i naturali in modo equivalente, con il linguaggio
degli insiemi e delle funzioni, e Peano ne cita il lavoro nella prefazione. Somma e
prodotto si definiscono a partire dal successore per ricorrenza, come aveva già fatto
Hermann Grassmann nel 1861, e le loro proprietà si dimostrano per induzione. Dalla
simbologia di Peano viene anche un segno che usiamo dal \cref{sec:insiemi}: la lettera
$\varepsilon$, iniziale della parola greca \emph{esti} («è»), con cui Peano scriveva
$a\,\varepsilon\,\N$ per dire «$a$ è un numero»; da essa è nato il simbolo $\in$.
\end{notastorica}

\begin{idea}[La fila di tessere del domino]
L'immagine classica: una fila infinita di tessere del domino in piedi. La base (a) dice
che facciamo cadere la prima tessera; il passo induttivo (b) dice che ogni tessera, se
cade, urta la successiva e la fa cadere. Conclusione: cadono tutte. Si noti che
\emph{entrambe} le condizioni sono necessarie: senza (a) nessuna tessera cade mai;
senza (b) la prima cade e la seconda può restare in piedi. Si noti anche che nel passo
induttivo \emph{non} si afferma che $\mathcal{P}(n)$ è vera: si afferma solo che
\emph{se} lo fosse, lo sarebbe anche $\mathcal{P}(n+1)$. È un'implicazione, non
un'asserzione.
\end{idea}

\begin{proposizione}\label{prop:induzione-potenza}
La \eqref{eq:potenza-crescente-bis} vale per ogni $n\in\N$.
\end{proposizione}

\begin{schema}
\Ipotesi{$n\in\N$; \ $0\le x_{1}<x_{2}$.}
\Tesi{$x_{1}^{n}<x_{2}^{n}$.}
\end{schema}

\begin{strategia}
Si applica il \cref{teo:induzione} alla proposizione $\mathcal{P}(n)$ «per ogni
$0\le x_{1}<x_{2}$ risulta $x_{1}^{n}<x_{2}^{n}$». La base è l'ipotesi stessa; il passo
induttivo è il calcolo \eqref{eq:passo-potenza} già svolto: si tratta solo di
riconoscerlo come tale.
\end{strategia}

\begin{dimostrazione}
Sia $\mathcal{P}(n)$ la proposizione «per ogni $0\le x_{1}<x_{2}$ risulta
$x_{1}^{n}<x_{2}^{n}$». $\mathcal{P}(1)$ è vera (afferma che $x_{1}<x_{2}$). Il passo
induttivo è esattamente la \eqref{eq:passo-potenza}. Per il principio di induzione,
$\mathcal{P}(n)$ è vera per ogni $n$.
\end{dimostrazione}

Per chiarire meglio il principio consideriamo altri esempi: alcune formule classiche
che useremo più avanti, in particolare nella teoria delle serie e nella definizione
dell'integrale.

\begin{proposizione}[Somma dei primi $n$ numeri naturali]\label{prop:gauss}
Per ogni $n\in\N$,
\begin{equation}\label{eq:gauss}
  1+2+3+\dots+(n-1)+n=\frac{n(n+1)}{2}.
\end{equation}
\end{proposizione}

\begin{schema}
\Ipotesi{$n\in\N$.}
\Tesi{$1+2+\dots+n=\dfrac{n(n+1)}{2}$.}
\end{schema}

\begin{strategia}
In una dimostrazione per induzione di una \emph{formula} il passo induttivo consiste
sempre nel partire dal membro sinistro con $n+1$, isolare il membro sinistro con $n$,
sostituirlo con il membro destro con $n$ (ipotesi induttiva) e fare i conti fino a
ottenere il membro destro con $n+1$. Qui: $1+\dots+n+(n+1)=\frac{n(n+1)}{2}+(n+1)$, e
basta raccogliere $(n+1)$.
\end{strategia}

\begin{dimostrazione}
La formula è vera per $n=1$; infatti si ha l'identità $1=(1\cdot2)/2$. Supponiamo vera
la \eqref{eq:gauss} per un certo $n$ e dimostriamo la formula analoga con l'indice
$n+1$ al posto di $n$. Per ottenere ciò, è naturale sommare a entrambi i membri il
numero $n+1$:
\begin{equation}\label{eq:gauss-passo}
\begin{split}
  1+2+3+\dots+n&+(n+1)=\frac{n(n+1)}{2}+(n+1)\\
  &=\frac{n(n+1)+2(n+1)}{2}=\frac{(n+1)(n+2)}{2}.
\end{split}
\end{equation}
Abbiamo ottenuto ciò che volevamo: l'ultimo membro è esattamente il secondo membro della
\eqref{eq:gauss} con $n+1$ al posto di $n$. Quindi la \eqref{eq:gauss} risulta vera per
ogni $n\in\N$.
\end{dimostrazione}

\begin{notastorica}[Il piccolo Gauss]
Alla formula \eqref{eq:gauss} è legato un celebre aneddoto su Carl Friedrich Gauss
(1777--1855), raccontato dal suo primo biografo, Wolfgang Sartorius von Waltershausen
(1856): il maestro Büttner avrebbe assegnato alla classe la somma di una lunga
progressione aritmetica, e Gauss, che aveva allora circa nove anni, avrebbe posato quasi
subito sul tavolo la sua lavagnetta con il risultato esatto. Nella versione divenuta
popolare, fissata solo molto più tardi, i numeri da sommare sono quelli da $1$ a $100$ e
il risultato è $5050$. L'idea: si
accoppiano il primo e l'ultimo numero ($1+100=101$), il secondo e il penultimo
($2+99=101$), e così via: $50$ coppie da $101$. La dimostrazione per induzione che
abbiamo dato è meno brillante ma ha un pregio: funziona sempre, anche quando non si ha
un'idea geniale --- purché si conosca già la formula da provare. Il principio di
induzione \emph{verifica}, non \emph{scopre}.
\end{notastorica}

\begin{marginfigure}
\centering
\begin{tikzpicture}[scale=0.3]
  \foreach \i in {1,...,6} \foreach \j in {1,...,\i} \fill[blu!70] (\j-1,\i-1) rectangle ++(0.9,0.9);
  \foreach \i in {1,...,6} \foreach \j in {\i,...,6} \fill[rame!50] (\j,\i-1) rectangle ++(0.9,0.9);
  \draw[grigio, thin] (0,0) rectangle (7,6);
  \node[font=\scriptsize, below] at (3.5,-0.1) {$n+1$};
  \node[font=\scriptsize, left] at (-0.1,3) {$n$};
\end{tikzpicture}
\caption{Una «dimostrazione senza parole» della \eqref{eq:gauss}: due copie di
$1+2+\dots+n$ riempiono un rettangolo $n\times(n+1)$.}
\end{marginfigure}

Un'altra applicazione del principio di induzione è la seguente disuguaglianza, che
sarà fondamentale nello studio delle successioni.

\begin{teorema}[Disuguaglianza di Bernoulli]\label{teo:bernoulli}\index{disuguaglianza di Bernoulli}
Per ogni numero reale $x\ge-1$ e per ogni naturale $n$ risulta
\begin{equation}\label{eq:bernoulli}
  (1+x)^{n}\ge1+nx .
\end{equation}
\end{teorema}

\begin{schema}
\Ipotesi{$x\in\R$, $x\ge-1$; $n\in\N$.}
\Tesi{$(1+x)^{n}\ge1+nx$.}
\end{schema}

\begin{strategia}
Induzione su $n$ con $x$ fissato. Il passo induttivo richiede di moltiplicare
l'ipotesi $(1+x)^{n}\ge1+nx$ per $1+x$: ed è qui che serve l'ipotesi $x\ge-1$, perché
$1+x\ge0$ e la disuguaglianza non cambia verso (\cref{prop:ordine-prodotto}). Poi si
sviluppa il prodotto e si butta via il termine $nx^{2}\ge0$: è il tipico passaggio in
cui si \emph{indebolisce} una stima per ottenere la forma voluta.
\end{strategia}

\begin{dimostrazione}
Per $n=1$ la proposizione è vera (con il segno $=$). Supponiamo vera la
\eqref{eq:bernoulli} per un numero naturale $n$ e moltiplichiamo entrambi i membri per
$1+x$, che è una quantità maggiore o uguale a zero:
\begin{equation}\label{eq:bernoulli-passo}
  (1+x)^{n+1}\ge(1+nx)(1+x)=1+x+nx+nx^{2}\ge1+(n+1)x,
\end{equation}
dove nell'ultimo passaggio abbiamo trascurato il termine $nx^{2}\ge0$. Abbiamo ottenuto
la proposizione con $n+1$ al posto di $n$. Perciò, in base al principio di induzione, la
\eqref{eq:bernoulli} è provata.
\end{dimostrazione}

\begin{marginfigure}
\centering
\begin{tikzpicture}
\begin{axis}[width=5.2cm, height=4.4cm, xmin=-1.3, xmax=1.3, ymin=-2.5, ymax=5,
  xtick={-1,1}, ytick={1}, samples=100]
  \addplot[curva, domain=-1.25:1.25] {(1+x)^4};
  \addplot[curva2, domain=-1.25:1.25] {1+4*x};
  \node[blu, font=\scriptsize, anchor=east] at (axis cs:-0.12,3.9) {$(1+x)^{4}$};
  \node[rame, font=\scriptsize, anchor=west] at (axis cs:0.3,0.8) {$1+4x$};
  \addplot[only marks, mark=*, mark size=1.2pt, grigio] coordinates {(0,1)};
\end{axis}
\end{tikzpicture}
\caption{La disuguaglianza di Bernoulli per $n=4$: la curva $(1+x)^{4}$ sta sopra la
retta $1+4x$ (sua tangente in $x=0$) per ogni $x\ge-1$.}
\end{marginfigure}

\begin{osservazione}
Che cosa succede se $x<-1$? Per $n=2$ la \eqref{eq:bernoulli} vale ancora per ogni
$x$ (equivale a $x^{2}\ge0$), ma per $n=3$ e $x=-4$ si ha
$(1-4)^{3}=-27<-11=1+3\cdot(-4)$: la disuguaglianza è falsa. L'ipotesi $x\ge-1$ non è
dunque un eccesso di prudenza. Il lettore controlli, ripercorrendo la dimostrazione,
qual è esattamente il passaggio che fallisce quando $1+x<0$.
\end{osservazione}

Utilizzando il principio di induzione dimostriamo la formula che esprime la
\emph{somma di una progressione geometrica} di ragione $x\neq1$.

\begin{proposizione}[Somma della progressione geometrica]\label{prop:geometrica}\index{progressione geometrica}
Per ogni $x\in\R$, $x\neq1$, e per ogni $n\in\N$,
\begin{equation}\label{eq:geometrica}
  1+x+x^{2}+\dots+x^{n}=\frac{1-x^{n+1}}{1-x}.
\end{equation}
\end{proposizione}

\begin{schema}
\Ipotesi{$x\in\R$, \ $x\neq1$; \ $n\in\N$.}
\Tesi{$1+x+\dots+x^{n}=\dfrac{1-x^{n+1}}{1-x}$.}
\end{schema}

\begin{strategia}
Induzione su $n$, con $x$ fissato, secondo lo schema della \cref{prop:gauss}: si
aggiunge $x^{n+1}$ ai due membri dell'ipotesi induttiva e si riduce allo stesso
denominatore $1-x$, che non è nullo perché $x\neq1$. I termini $\pm x^{n+1}$ si
cancellano e resta il secondo membro con $n+1$ al posto di $n$.
\end{strategia}

\begin{dimostrazione}
Per $n=1$, a secondo membro abbiamo
\[
  \frac{1-x^{2}}{1-x}=\frac{(1-x)(1+x)}{1-x}=1+x;
\]
quindi la \eqref{eq:geometrica} è vera per $n=1$. Supponendo verificata la
\eqref{eq:geometrica} per un certo $n$, sommiamo a entrambi i membri il termine
$x^{n+1}$:
\begin{equation}\label{eq:geometrica-passo}
\begin{split}
  1+x+x^{2}+\dots+x^{n}+x^{n+1}&=\frac{1-x^{n+1}}{1-x}+x^{n+1}\\
  &=\frac{1-x^{n+1}+x^{n+1}-x^{n+2}}{1-x}=\frac{1-x^{n+2}}{1-x}.
\end{split}
\end{equation}
Abbiamo ottenuto la proposizione con $n+1$ al posto di $n$. Quindi, in base al
principio di induzione, la formula \eqref{eq:geometrica} è dimostrata.
\end{dimostrazione}

\begin{osservazione}[Una dimostrazione diretta]
La \eqref{eq:geometrica} si può anche ottenere senza induzione: posto
$S=1+x+\dots+x^{n}$, si ha $xS=x+x^{2}+\dots+x^{n+1}$, e sottraendo
$S-xS=1-x^{n+1}$ (tutti i termini intermedi si cancellano), da cui $S(1-x)=1-x^{n+1}$.
Ma anche questa dimostrazione nasconde un'induzione: il «si cancellano» per una somma
di $n$ termini è esso stesso un'affermazione che, a rigore, va provata per induzione.
\end{osservazione}

Proviamo mediante il principio di induzione una formula, analoga alla \eqref{eq:gauss},
che verrà utilizzata introducendo gli integrali definiti.

\needspace{16\baselineskip}
\begin{proposizione}[Somma dei primi $n$ quadrati]\label{prop:quadrati}
Per ogni $n\in\N$,
\begin{equation}\label{eq:quadrati}
  1^{2}+2^{2}+3^{2}+\dots+n^{2}=\frac{n(n+1)(2n+1)}{6}.
\end{equation}
\end{proposizione}

\begin{schema}
\Ipotesi{$n\in\N$.}
\Tesi{$1^{2}+2^{2}+\dots+n^{2}=\dfrac{n(n+1)(2n+1)}{6}$.}
\end{schema}

\begin{strategia}
Lo schema è quello della \cref{prop:gauss}, ma i conti sono più lunghi e conviene
organizzarli: dopo aver aggiunto $(n+1)^{2}$ si raccoglie il fattore comune $(n+1)$, si
riduce a un unico denominatore, e si confronta il polinomio di secondo grado ottenuto
con quello che compare nel membro destro con $n+1$ al posto di $n$. Due polinomi sono
uguali se hanno gli stessi coefficienti: è un controllo meccanico.
\end{strategia}

\begin{dimostrazione}
La formula \eqref{eq:quadrati} è vera per $n=1$; risulta infatti $1^{2}=(1\cdot2\cdot3)/6$.
Supponiamo che la \eqref{eq:quadrati} sia verificata per un indice $n$ e deduciamo da
essa la formula analoga con $n+1$ al posto di $n$. Abbiamo
\begin{equation}\label{eq:quadrati-passo}
\begin{split}
  1^{2}+2^{2}&+\dots+n^{2}+(n+1)^{2}
  =\bigl[1^{2}+2^{2}+\dots+n^{2}\bigr]+(n+1)^{2}\\
  &=\frac{n(n+1)(2n+1)}{6}+(n+1)^{2}\\
  &=(n+1)\Bigl[\frac{n(2n+1)}{6}+(n+1)\Bigr]\\
  &=(n+1)\,\frac{2n^{2}+n+6n+6}{6}
   =(n+1)\,\frac{2n^{2}+7n+6}{6}.
\end{split}
\end{equation}
Osservando che la \eqref{eq:quadrati} con $n+1$ al posto di $n$ ha a secondo membro
l'espressione
\begin{equation}\label{eq:quadrati-target}
  \frac{(n+1)(n+2)\bigl[2(n+1)+1\bigr]}{6}=\frac{(n+1)(n+2)(2n+3)}{6},
\end{equation}
rimane soltanto da osservare che
\[
  (n+1)\,\frac{2n^{2}+7n+6}{6}=\frac{(n+1)(n+2)(2n+3)}{6},
\]
che è una relazione soddisfatta perché, da verifica diretta, risulta
$2n^{2}+7n+6=(n+2)(2n+3)$.
\end{dimostrazione}

\begin{proposizione}[Una somma telescopica]\label{prop:telescopica}
Per ogni $n\in\N$,
\begin{equation}\label{eq:telescopica}
  \frac{1}{1\cdot2}+\frac{1}{2\cdot3}+\dots+\frac{1}{n(n+1)}=\frac{n}{n+1}.
\end{equation}
\end{proposizione}

\begin{schema}
\Ipotesi{$n\in\N$.}
\Tesi{$\dfrac{1}{1\cdot2}+\dots+\dfrac{1}{n(n+1)}=\dfrac{n}{n+1}$.}
\end{schema}

\begin{strategia}
Ancora lo schema consueto: si aggiunge il termine $\frac{1}{(n+1)(n+2)}$ ai due membri
dell'ipotesi induttiva, si riduce allo stesso denominatore e si riconosce al numeratore
il quadrato $(n+1)^{2}$, che si semplifica con il fattore $n+1$ del denominatore. Una
dimostrazione diretta, senza induzione, è indicata nell'osservazione che segue.
\end{strategia}

\begin{dimostrazione}
Da verifica diretta la formula \eqref{eq:telescopica} è soddisfatta per $n=1$.
Supponendo poi che la \eqref{eq:telescopica} sia verificata per un indice $n$ generico,
otteniamo
\begin{equation}\label{eq:telescopica-passo}
\begin{split}
  \frac{1}{1\cdot2}&+\dots+\frac{1}{n(n+1)}+\frac{1}{(n+1)(n+2)}
  =\frac{n}{n+1}+\frac{1}{(n+1)(n+2)}\\
  &=\frac{n(n+2)+1}{(n+1)(n+2)}=\frac{n^{2}+2n+1}{(n+1)(n+2)}\\
  &=\frac{(n+1)^{2}}{(n+1)(n+2)}=\frac{n+1}{n+2},
\end{split}
\end{equation}
che corrisponde appunto alla \eqref{eq:telescopica} con $n+1$ al posto di $n$.
\end{dimostrazione}

\begin{osservazione}
Il nome \emph{telescopica} viene dall'identità $\frac{1}{k(k+1)}=\frac1k-\frac{1}{k+1}$:
sommando per $k=1,\dots,n$ i termini intermedi si elidono a coppie, come i tubi di un
cannocchiale che rientrano l'uno nell'altro, e resta $1-\frac{1}{n+1}=\frac{n}{n+1}$.
Ritroveremo questa idea nello studio delle serie.
\end{osservazione}

\begin{python}[Verificare le formule con \texttt{sympy} e per forza bruta]
import sympy as sp
n, k, x = sp.symbols("n k x", positive=True, integer=True)

# sympy calcola le somme in forma chiusa: confrontiamo con le nostre
print(sp.factor(sp.summation(k, (k, 1, n))))          # Gauss
print(sp.factor(sp.summation(k**2, (k, 1, n))))       # quadrati
s = sp.summation(1/(k*(k+1)), (k, 1, n))
print(sp.simplify(s))                                 # telescopica

# Il passo induttivo della somma dei quadrati: due polinomi uguali?
print(sp.expand(2*n**2 + 7*n + 6 - (n + 2)*(2*n + 3)) == 0)

# Forza bruta: controlliamo la formula di Gauss per n = 1, ..., 1000
ok = (sum(range(1, m + 1)) == m*(m + 1)//2 for m in range(1, 1001))
print(all(ok))
\end{python}
\begin{risultato}
n*(n + 1)/2
n*(n + 1)*(2*n + 1)/6
n/(n + 1)
True
True
\end{risultato}

\begin{attenzione}[Mille casi non sono una dimostrazione]
L'ultima riga del programma controlla la \eqref{eq:gauss} per $n$ fino a $1000$:
rassicurante, ma \emph{non} è una dimostrazione. L'affermazione «$n^{2}+n+41$ è primo
per ogni $n$» è vera per $n=0,1,\dots,39$ e falsa per $n=40$ (Eulero); l'affermazione
«$991n^{2}+1$ non è mai un quadrato perfetto» è vera per tutti gli $n$ fino a circa
$10^{28}$, e poi falsa. Solo il principio di induzione (o un'altra dimostrazione)
garantisce \emph{tutti} gli $n$.
\end{attenzione}

\begin{esercizio}
Dimostrare per induzione che un insieme con $n$ elementi ha esattamente $2^{n}$
sottoinsiemi.
\begin{soluzione}
Per $n=1$: $\{a\}$ ha i due sottoinsiemi $\emptyset$ e $\{a\}$. Supponiamo che ogni
insieme con $n$ elementi abbia $2^{n}$ sottoinsiemi e sia $S$ un insieme con $n+1$
elementi; fissiamo $s\in S$. I sottoinsiemi di $S$ si dividono in quelli che non
contengono $s$ (sono i sottoinsiemi di $S\setminus\{s\}$, che ha $n$ elementi: $2^{n}$)
e quelli che contengono $s$ (si ottengono dai precedenti aggiungendo $s$: ancora
$2^{n}$). In tutto $2^{n}+2^{n}=2^{n+1}$. \emph{Scelta logica:} per il passo induttivo
bisogna ricondurre un insieme di $n+1$ elementi a uno di $n$: lo si fa isolando un
elemento $s$ e dividendo i sottoinsiemi in due classi disgiunte, a seconda che
contengano $s$ oppure no; ciascuna classe è in corrispondenza biunivoca con i
sottoinsiemi di $S\setminus\{s\}$.
\end{soluzione}
\end{esercizio}

\begin{esercizio}
Dimostrare per induzione che $2^{n}>n$ per ogni $n\in\N$, e che $n!\ge2^{n-1}$, dove
$n!=1\cdot2\cdots n$ è il \emph{fattoriale} di $n$.
\begin{soluzione}
(a) Per $n=1$: $2^{1}=2>1$. Se $2^{n}>n$, allora $2^{n+1}=2\cdot2^{n}>2n=n+n\ge n+1$,
perché $n\ge1$. (b) Per $n=1$: $1!=1=2^{0}$. Se $n!\ge2^{n-1}$, allora
$(n+1)!=(n+1)\,n!\ge(n+1)\,2^{n-1}\ge2\cdot2^{n-1}=2^{n}$, perché $n+1\ge2$ e
$2^{n-1}>0$. \emph{Scelta logica:} in entrambi i casi il passo induttivo confronta il
fattore per cui si moltiplica il primo membro passando da $n$ a $n+1$ ($2$ per $2^{n}$,
$n+1$ per $n!$) con quello del secondo membro; come nella disuguaglianza di Bernoulli,
si conclude \emph{indebolendo} la stima ($n\ge1$, $n+1\ge2$).
\end{soluzione}
\end{esercizio}

%% ----------------------------------------------------------------------------
%%  §1.11  Esercizi di riepilogo
%% ----------------------------------------------------------------------------
\section{Esercizi di riepilogo}
\label{sec:esercizi-cap1}

Gli esercizi di questo paragrafo non servono a «fare pratica» nel senso scolastico: ogni
soluzione è accompagnata da un commento sulle \emph{scelte logiche} che la guidano,
perché è lì che si impara a dimostrare. Il lettore provi sempre prima da solo.

\begin{esercizio}[Unicità di $\sqrt2$]
Dimostrare che esiste al più un numero reale positivo $c$ tale che $c^{2}=2$.
\begin{soluzione}
Siano $c_{1},c_{2}>0$ con $c_{1}^{2}=c_{2}^{2}=2$. Allora
$0=c_{1}^{2}-c_{2}^{2}=(c_{1}-c_{2})(c_{1}+c_{2})$; per la legge di annullamento del
prodotto (\cref{prop:annullamento}) uno dei fattori è nullo; ma $c_{1}+c_{2}>0$, quindi
$c_{1}=c_{2}$. \emph{Scelta logica:} per provare l'unicità si suppongono due oggetti e si
mostra che coincidono; la fattorizzazione $a^{2}-b^{2}=(a-b)(a+b)$ trasforma un
problema di uguaglianza in un problema di annullamento, dove abbiamo un teorema
pronto. Si noti che l'\emph{esistenza} di $c$ richiede invece l'assioma di completezza
(\cref{sec:nzq}): esistenza e unicità sono due questioni logicamente distinte.
\end{soluzione}
\end{esercizio}

\begin{esercizio}[Dove si usa la completezza?]
Si consideri l'insieme $A=\{x\in\Q : x<1\}$ e $B=\{x\in\Q : x\ge1\}$. Esiste un
razionale $c$ che separa $A$ e $B$? È in contraddizione con il
\cref{teo:Q-non-completo}?
\begin{soluzione}
Sì: $c=1$ separa $A$ e $B$ (anzi è l'unico separatore). Nessuna contraddizione: il
\cref{teo:Q-non-completo} afferma che l'assioma di completezza fallisce in $\Q$ per
\emph{almeno una} coppia di insiemi, non per tutte. \emph{Scelta logica:} la negazione di
«per ogni coppia $A,B$ esiste $c$» è «esiste una coppia $A,B$ per cui non esiste $c$»;
per confutare un'affermazione universale basta un controesempio, e per quello abbiamo
dovuto scegliere con cura gli insiemi della \eqref{eq:AB-radice2}.
\end{soluzione}
\end{esercizio}

\begin{esercizio}[Disuguaglianza tra media aritmetica e geometrica]
Dimostrare che per ogni $a,b\ge0$ si ha $\sqrt{ab}\le\dfrac{a+b}{2}$, con uguaglianza se
e solo se $a=b$.
\begin{soluzione}
Poiché entrambi i membri sono $\ge0$ e la funzione $t\mapsto t^{2}$ è strettamente
crescente su $[0,+\infty)$ (\cref{prop:potenza-crescente}), la disuguaglianza equivale
a $ab\le\frac{(a+b)^{2}}{4}$, cioè a $4ab\le a^{2}+2ab+b^{2}$, cioè a
$0\le a^{2}-2ab+b^{2}=(a-b)^{2}$, che è vera per \ref{ass:ordine-prodotto} (e la regola
dei segni). L'uguaglianza vale se e solo se $(a-b)^{2}=0$, cioè $a=b$. \emph{Scelta
logica:} la catena di «equivale a» è lecita solo perché a ogni passo abbiamo usato
un'equivalenza: la stretta crescenza del quadrato sui numeri non negativi, che vale nei
due versi, la moltiplicazione per $4>0$ (\cref{prop:ordine-prodotto}) e l'assioma
\ref{ass:ordine-somma}, che permette di sommare o sottrarre uno stesso termine ai due
membri; se avessimo usato una semplice implicazione, avremmo dovuto percorrere la
catena a ritroso.
\end{soluzione}
\end{esercizio}

\begin{marginfigure}
\centering
\begin{tikzpicture}[scale=0.9]
  \draw[grigio!70] (0,0) -- (4,0);
  \draw[grigio!70] (4,0) arc (0:180:2);
  \draw[blu, thick] (1.2,0) -- (1.2,{sqrt(1.2*2.8)});
  \draw[rame, thick] (2,0) -- (2,2);
  \draw[tratteggio] (0,0) -- (1.2,{sqrt(1.2*2.8)}) -- (4,0);
  \node[font=\scriptsize, below] at (0.6,0) {$a$};
  \node[font=\scriptsize, below] at (2.6,0) {$b$};
  \node[font=\scriptsize, blu, left] at (1.2,0.45) {$\sqrt{ab}$};
  \node[font=\scriptsize, rame, right] at (2,0.55) {$\tfrac{a+b}{2}$};
\end{tikzpicture}
\caption{Media aritmetica e geometrica: in una semicirconferenza di diametro $a+b$,
l'altezza $\sqrt{ab}$ (secondo teorema di Euclide) non supera mai il raggio.}
\end{marginfigure}

\begin{esercizio}[Funzioni e monotonia]
Sia $f\colon\R\to\R$ strettamente crescente. Dimostrare che per ogni $y\in\R$
l'equazione $f(x)=y$ ha \emph{al più} una soluzione. Mostrare con un esempio che può
non averne nessuna.
\begin{soluzione}
Se $f(x_{1})=f(x_{2})=y$ con $x_{1}\neq x_{2}$, diciamo $x_{1}<x_{2}$, la stretta
crescenza darebbe $f(x_{1})<f(x_{2})$: assurdo. Per l'esempio: $f(x)=e^{x}$ è
strettamente crescente, ma $e^{x}=-1$ non ha soluzioni. \emph{Scelta logica:} «al più
una» è un'affermazione di unicità e si tratta come sempre; l'esempio mostra che
l'unicità non implica l'esistenza, cioè che iniettività e suriettività sono proprietà
indipendenti (\cref{fig:iniettiva-suriettiva}).
\end{soluzione}
\end{esercizio}

\begin{esercizio}[Valore assoluto]
Dimostrare che per ogni $x,y\in\R$ vale $\abs{x\,y}=\abs{x}\,\abs{y}$ \emph{senza}
distinguere i casi, usando l'identità $\abs{t}=\sqrt{t^{2}}$.
\begin{soluzione}
Si ha $\abs{t}^{2}=t^{2}$ per ogni $t$ (in entrambi i casi della definizione), e
$\abs{t}\ge0$; poiché la radice quadrata è l'inversa del quadrato su $[0,+\infty)$,
$\abs{t}=\sqrt{t^{2}}$. Allora
$\abs{xy}=\sqrt{(xy)^{2}}=\sqrt{x^{2}y^{2}}=\sqrt{x^{2}}\sqrt{y^{2}}=\abs{x}\abs{y}$,
dove abbiamo usato $\sqrt{uv}=\sqrt u\sqrt v$ per $u,v\ge0$ (che segue dall'unicità
della radice: $(\sqrt u\sqrt v)^{2}=uv$). \emph{Scelta logica:} rappresentare un
oggetto definito per casi con una formula unica elimina la distinzione dei casi; il
prezzo è dover usare proprietà della radice, cioè in ultima analisi la completezza.
Entrambe le dimostrazioni sono corrette: quella per casi usa meno strumenti, questa è
più corta.
\end{soluzione}
\end{esercizio}

\begin{esercizio}[Induzione: dove sta l'errore?]
Si consideri la seguente «dimostrazione» che tutti i cavalli hanno lo stesso colore.
\emph{Per $n=1$ è ovvio. Supponiamo che in ogni insieme di $n$ cavalli tutti abbiano lo
stesso colore, e consideriamo $n+1$ cavalli $c_{1},\dots,c_{n+1}$. I primi $n$ hanno lo
stesso colore per ipotesi induttiva; anche gli ultimi $n$, $c_{2},\dots,c_{n+1}$, hanno lo
stesso colore. Poiché $c_{2}$ sta in entrambi i gruppi, tutti gli $n+1$ cavalli hanno lo
stesso colore.} Dove sta l'errore?
\begin{soluzione}
Il passo induttivo fallisce per $n=1$: i due gruppi $\{c_{1}\}$ e $\{c_{2}\}$ non hanno
elementi in comune, e il «ponte» $c_{2}$ non esiste. Il passo $\mathcal{P}(n)\implica
\mathcal{P}(n+1)$ è corretto per ogni $n\ge2$, ma la base $\mathcal{P}(1)$ non si
collega a $\mathcal{P}(2)$: la prima tessera del domino cade, ma non urta la seconda.
\emph{Scelta logica:} nel verificare il passo induttivo bisogna assicurarsi che il
ragionamento valga \emph{per ogni} $n$ a partire dalla base, inclusi i primi valori,
che sono spesso quelli critici.
\end{soluzione}
\end{esercizio}

\needspace{14\baselineskip}
\begin{esercizio}[Con Python]
Scrivere una funzione Python \texttt{radice(y, n)} che, dato $y\ge0$, restituisca
un'approssimazione di $\sqrt{y}$ con errore minore di $10^{-n}$ usando il metodo di
bisezione del \cref{sec:nzq}: si parte dall'intervallo $[0,\max\{1,y\}]$ e lo si dimezza
finché la sua lunghezza è minore di $10^{-n}$. Quante iterazioni servono per $n=12$?
Perché il metodo funziona (quale teorema di questo capitolo garantisce che il punto
medio $m$ soddisfa $m^{2}<y$ oppure $m^{2}\ge y$, e che la scelta è coerente con
l'ordine)?
\begin{soluzione}
\begin{python}
def radice(y, n):
    a, b = 0.0, max(1.0, y)          # a^2 <= y <= b^2
    iterazioni = 0
    while b - a >= 10**(-n):
        m = (a + b) / 2
        if m * m < y:
            a = m                    # m è "troppo piccolo"
        else:
            b = m                    # m è "troppo grande"
        iterazioni += 1
    return (a + b) / 2, iterazioni

print(radice(2, 12))                 # (1.414213562372879, 41)
\end{python}
La dicotomia \ref{ass:dicotomia} garantisce che vale $m^{2}<y$ oppure $m^{2}\ge y$;
la \cref{prop:potenza-crescente} (stretta crescenza di $t\mapsto t^{2}$ sui non
negativi) garantisce che se $m^{2}<y$ allora la radice cercata è maggiore di $m$, e
viceversa. Servono $41$ iterazioni: l'intervallo iniziale $[0,2]$ ha lunghezza $2$, dopo
$k$ dimezzamenti ha lunghezza $2^{1-k}$, e $2^{-40}\approx9{,}1\cdot10^{-13}$ è la prima
potenza sotto $10^{-12}$. Ogni passo guadagna una cifra binaria, cioè circa $0{,}3$
cifre decimali. \emph{Scelta logica:} la bisezione usa soltanto l'ordinamento, e
l'invariante $a^{2}\le y\le b^{2}$, conservato a ogni passo, garantisce che la radice
cercata stia sempre in $[a,b]$: l'errore è quindi controllato a priori dalla lunghezza
dell'intervallo, che si dimezza, e il numero di iterazioni si calcola in anticipo.
\end{soluzione}
\end{esercizio}

%% ----------------------------------------------------------------------------
%%  §1.12  Punti fondamentali del capitolo
%% ----------------------------------------------------------------------------
\section{Punti fondamentali del capitolo}
\label{sec:sintesi-cap01}

Il capitolo ha fissato le regole del gioco: gli undici assiomi dei numeri reali, tra i
quali l'assioma di completezza è l'unico che distingue $\R$ da $\Q$. Su questa base
abbiamo costruito il linguaggio che useremo in tutto il libro (insiemi, funzioni,
induzione) e il catalogo delle funzioni elementari.

\begin{puntichiave}
\item \textbf{Gli assiomi.} $\R$ soddisfa le proprietà di campo (A1)--(A6), le proprietà
dell'ordinamento (O1)--(O4) e l'assioma di completezza (\cref{ass:completezza}): se $A$ e
$B$ sono non vuoti e $a\le b$ per ogni $a\in A$ e $b\in B$, esiste $c\in\R$ con
$a\le c\le b$ per ogni $a\in A$ e $b\in B$. Ogni altra proprietà si dimostra.
\item \textbf{Le regole dell'algebra sono teoremi.} La legge di annullamento del prodotto
(\cref{prop:annullamento}), la regola dei segni (\cref{prop:segni}) e il cambio di verso
di una disuguaglianza moltiplicata per $c\le0$ (\cref{prop:ordine-prodotto}) discendono
dagli assiomi; la divisione per zero è impossibile perché porterebbe a $0=1$.
\item \textbf{Il linguaggio degli insiemi.} Due insiemi si dimostrano uguali per doppia
inclusione; le classi di una relazione di equivalenza formano una partizione
(\cref{prop:classi}); la relazione $\le$ è un ordine totale su $\R$.
\item \textbf{$\Q$ non è completo.} Nessun razionale ha quadrato $2$
(\cref{teo:radice2}), e gli insiemi \eqref{eq:AB-radice2} non sono separati da alcun
razionale (\cref{teo:Q-non-completo}): $\Q$ soddisfa (A1)--(A6) e (O1)--(O4), ma non la
completezza, che in $\R$ garantisce l'esistenza di $\sqrt2$.
\item \textbf{Funzioni.} Una funzione associa a ogni elemento del dominio \emph{uno e un
solo} valore; è invertibile se e solo se è biunivoca, e allora
$f^{-1}\circ f=\mathrm{id}_{A}$ e $f\circ f^{-1}=\mathrm{id}_{B}$.
\item \textbf{Criterio di invertibilità.} Una funzione $f\colon A\to B$ suriettiva e
strettamente monotòna è invertibile, e $f^{-1}$ è strettamente monotòna dello stesso tipo
(\cref{teo:monotona-invertibile}); senza ipotesi di continuità il viceversa è falso.
\item \textbf{Valore assoluto.} Per $r\ge0$ si ha $\abs{x}\le r\sse-r\le x\le r$
(\cref{prop:modulo-r}); valgono la disuguaglianza triangolare
$\abs{x_{1}+x_{2}}\le\abs{x_{1}}+\abs{x_{2}}$ (\cref{teo:triangolare}) e la
$\bigl|\abs{x_{1}}-\abs{x_{2}}\bigr|\le\abs{x_{1}-x_{2}}$ (\cref{cor:triangolare-inversa}).
\item \textbf{Potenze, esponenziali, logaritmi.} La potenza $x^{n}$ è strettamente
crescente per $x\ge0$ (\cref{prop:potenza-crescente}); per $a>0$, $a\neq1$,
l'esponenziale $a^{x}$ è strettamente monotòna e il logaritmo ne è l'inversa,
$y=\log_{a}x\sse a^{y}=x$ per $x>0$; il logaritmo trasforma prodotti in somme
(\cref{prop:log-proprieta}), ed è questo il principio del regolo calcolatore.
\item \textbf{Trigonometria.} Gli angoli si misurano in radianti; $\sen x$ e $\cos x$
hanno periodo $2\pi$, $\tg x$ ha periodo $\pi$; $\sen^{2}x+\cos^{2}x=1$ e, per
$0<x<\pi/2$, $0<\sen x<x<\tg x$ \eqref{eq:sen-x-tg}, l'unico fatto geometrico che
useremo per derivare le funzioni trigonometriche.
\item \textbf{Il principio di induzione.} Se $\mathcal{P}(1)$ è vera e
$\mathcal{P}(n)\implica\mathcal{P}(n+1)$ per ogni $n\in\N$, allora $\mathcal{P}(n)$ è vera
per ogni $n\in\N$ (\cref{teo:induzione}); con esso si provano le formule delle somme
(\cref{prop:gauss}, \cref{prop:geometrica}, \cref{prop:quadrati},
\cref{prop:telescopica}).
\item \textbf{Disuguaglianza di Bernoulli.} Per $x\ge-1$ e $n\in\N$ si ha
$(1+x)^{n}\ge1+nx$ (\cref{teo:bernoulli}); l'ipotesi $x\ge-1$ è essenziale: per $n=3$ e
$x=-4$ la disuguaglianza è falsa.
\item \textbf{Errori tipici.} $f^{-1}$ non è il reciproco $1/f$; la controimmagine
$f^{-1}(Y)$ ha senso anche se $f$ non è invertibile; la relazione $y^{2}=x$ non definisce
una funzione di $x$; un'affermazione verificata in mille casi non è dimostrata.
\item \textbf{Il calcolatore e i numeri reali.} I numeri in virgola mobile sono un insieme
finito che non soddisfa gli assiomi (\texttt{0.1 + 0.2 == 0.3} è falso): si confrontano
con una tolleranza (\texttt{math.isclose}), e i grafici «per punti» vanno interpretati
alla luce della definizione.
\item \textbf{Per l'ingegnere.} Il radiante è l'unità naturale delle velocità angolari
($v=\omega r$); la bisezione approssima $\sqrt2$ con un errore controllato a priori, che
si dimezza a ogni passo, ed è il metodo che ritroveremo per gli zeri delle funzioni
continue.
\end{puntichiave}

\paragraph{Il filo del discorso.} Il capitolo lascia aperte tre questioni, tutte legate
alla completezza: l'esistenza della radice $n$-esima di ogni numero positivo, finora
soltanto annunciata; la proprietà di Archimede, aggirata nella dimostrazione del
\cref{teo:Q-non-completo}; il significato di $a^{x}$ per $x$ irrazionale. Il prossimo
capitolo trasforma l'assioma di completezza in uno strumento di lavoro, l'estremo
superiore, primo anello della catena dei teoremi di esistenza
(\cref{fig:schema-teoremi}), e con esso risolve tutte e tre le questioni.

